\chapter{往年真题整理}

本章按年份升序整理 \texttt{elegantbook2/试卷/} 中的往年工科数学分析试卷与答案资料。每份材料均已在 \texttt{word/past-exams/manifest.md} 中登记；PDF 先经 \texttt{pdftotext -layout} 抽取，扫描型材料则配合人工校对。

\begin{note}
本章正文已经直接并入主文件；\texttt{word/past-exams/reviewed/} 下的校对底稿与 \texttt{raw/} 下的原始抽取文本仅保留为中间审计材料。
\end{note}

\paragraph*{源文件审计总表}
下列源文件均来自 \texttt{elegantbook2/试卷/}，并已按年份与卷别并入本章相应小节；答案文件或学生答案文件不另开节，而是并入同卷的“答案与解析”。
\begin{enumerate}
  \item 2006 A 试题：2006数学分析下A卷.pdf。
  \item 2006 B 试题：2006数学分析下B卷.pdf。
  \item 2007 A 试题：2007数学分析下A卷.pdf。
  \item 2007 A 答案：2007数学分析下A卷及答案.pdf。
  \item 2009 A 试题：2009级计算机工科数学分析下 期末A.doc。
  \item 2009 B 试题：2009级计算机工科数学分析下 期末B.doc。
  \item 2012 A 试题：2012级计算机工科数学分析下 期末考试题(A卷).doc。
  \item 2012 B 试题：2012级计算机工科数学分析下 期末考试题(B卷) .doc。
  \item 2013 A 试题：2013级软件 工科数学分析下A.doc。
  \item 2013 A 答案：2013级软件 工科数学分析下A附解答.doc。
  \item 2013 B 试题：2013级软件 工科数学分析下B.doc。
  \item 2014 A 试题：2014级软件 工科数学分析下A.doc。
  \item 2014 B 试题：2014级软件 工科数学分析下B.doc。
  \item 2015 A 试题：2015级软件 工科数学分析下A.doc。
  \item 2015 B 试题：2015级软件 工科数学分析下B.doc。
  \item 2016 A 试题：2016级工科数学分析试卷A.doc。
  \item 2016 B 试题：2016级工科数学分析试卷B.doc。
  \item 2016 B 答案：2016级工科数学分析期末试卷B解答.pdf。
  \item 2017 A 试题：2017级工科数学分析期末试卷A.doc。
  \item 2017 A 答案：2017级工科数学分析期末试卷A答案.pdf。
  \item 2017 B 试题：2017级工科数学分析期末试卷B.doc。
  \item 2017 B 答案：2017级工科数学分析期末试卷B答案.pdf。
  \item 2018 A 试题：2018级工科数学分析（二）A.doc。
  \item 2018 A 答案：2018级工科数学分析（二）A附解答.doc。
  \item 2018 B 试题：2018级工科数学分析（二）B.doc。
  \item 2018 B 答案：2018级工科数学分析（二）B附解答.doc。
  \item 2019 A 试题：2019级 工科数学分析（二）A无答案.doc。
  \item 2019 A 答案一：2019级 工科数学分析（二）A答案1 .doc。
  \item 2019 A 答案二：2019级 工科数学分析（二）A答案2 .doc。
  \item 2019 B 补考试题：2019级 工科数学分析（二）B 补考用.doc。
  \item 2020 A 试题：2020级工数下A无答案.pdf。
  \item 2020 B 试题：2020级B卷无答案.pdf。
  \item 2021 A 试题：2021级工科数学分析下 期末考试题（A）卷.pdf。
  \item 2021 A 答案：2021级工科数学分析A答案.pdf。
  \item 2021 B 试题：2021级工科数学分析下 期末考试题（B）卷 .pdf。
  \item 2022 A 试题：2022级A卷无答案.pdf。
  \item 2022 A 学生答案：2022级A卷学生答案.pdf。
  \item 2022 B 试题：2022级工科数学分析下期末B无答案.pdf。
  \item 2022 B 答案：2022级工科数学分析下期末B有答案.pdf。
  \item 2023 A 试题：2023级工科数学分析下 期末考试题（A）卷 .pdf。
  \item 2023 B 试题：2023级工科数学分析下 期末考试题（B）卷 .pdf。
\end{enumerate}

% 2006--2007 reviewed draft.
% Source PDFs:
% - elegantbook2/试卷/2006数学分析下A卷.pdf
% - elegantbook2/试卷/2006数学分析下B卷.pdf
% - elegantbook2/试卷/2007数学分析下A卷.pdf
% - elegantbook2/试卷/2007数学分析下A卷及答案.pdf
% Raw OCR/text sources under word/past-exams/raw were used as auxiliary references;
% low-confidence OCR points have been resolved against page images where needed.

\section{2006--2007 第二学期 A 卷}

\subsection{试题}

\paragraph*{一、单项选择题（每小题 3 分，共 15 分）}

\begin{enumerate}
  \item 设有直线
  \[
    L:\begin{cases}
      x+3y+2z+1=0,\\
      2x-y-10z+3=0
    \end{cases}
  \]
  及平面 $\pi:4x-2y+z-2=0$，则直线 $L$（\quad）。

  A. 平行于平面 $\pi$；\quad
  B. 在平面 $\pi$ 上；\quad
  C. 垂直于平面 $\pi$；\quad
  D. 与平面 $\pi$ 斜交。

  \item 二元函数
  \[
    f(x,y)=
    \begin{cases}
      \dfrac{x^2y}{x^2+y^2}, & (x,y)\ne(0,0),\\
      0, & (x,y)=(0,0)
    \end{cases}
  \]
  在点 $(0,0)$ 处（\quad）。

  A. 偏导数存在，但不可微；\quad
  B. 连续、偏导数不存在；\quad
  C. 不连续、偏导数不存在；\quad
  D. 偏导数不存在。

  \item 设平面区域 \(\displaystyle D=\{(x,y)\mid -a\le x\le a,\ x\le y\le a\},\qquad
    D_1=\{(x,y)\mid 0\le x\le a,\ x\le y\le a\}\)，
  则
  \(\displaystyle \iint_D (xy+\cos x\sin y)\,dxdy\)
  等于（\quad）。

  A. $2\iint_{D_1}\cos x\sin y\,dxdy$；\quad
  B. $2\iint_{D_1}xy\,dxdy$；\quad
  C. $4\iint_{D_1}(xy+\cos x\sin y)\,dxdy$；\quad
  D. $0$。

  \item 若级数 $\sum_{n=1}^{\infty}a_n^2$，$\sum_{n=1}^{\infty}b_n^2$ 都收敛，则下列结论不成立的是（\quad）。

  A. $\sum_{n=1}^{\infty}(a_n+b_n)^2$ 收敛；\quad
  B. $\sum_{n=1}^{\infty}\dfrac{|a_n|}{n}$ 收敛；\quad
  C. $\sum_{n=1}^{\infty}a_nb_n$ 收敛；\quad
  D. $\sum_{n=1}^{\infty}a_nb_n$ 发散。

  \item 微分方程 \(\displaystyle y'+\frac{1}{x}y=\frac{\sin x}{x}\)
  的通解为（\quad）。

  A. $y=\dfrac{1}{x}(-\cos x+C)$；\quad
  B. $y=x(-\cos x+C)$；\quad
  C. $y=\dfrac{1}{x}(-\sin x+C)$；\quad
  D. $y=\dfrac{1}{x}(\cos x+C)$。
\end{enumerate}

\paragraph*{二、填空题（每小题 3 分，共 15 分）}

\begin{enumerate}
  \item 过点 $(0,2,4)$ 且同时平行于平面 $x+2z=1$ 和 $y-3z=2$ 的直线方程是 \underline{\hspace{4cm}}。

  \item 将积分 \(\displaystyle \int_1^e dx\int_0^{\ln x} f(x,y)\,dy\)
  改变积分顺序为 \underline{\hspace{4cm}}。
  \textit{（校勘：原卷内层微分疑似误印为 $dx$；按二重积分换序题意应为 $dy$。）}

  \item 级数 \(\displaystyle \sum_{n=1}^{\infty}\frac{3^n+(-2)^n}{n}(x+1)^n\)
  的收敛域为 \underline{\hspace{4cm}}。

  \item 设函数 $u=xyz$，则 $u$ 沿方向 $\ell=\{2,-1,3\}$ 的方向导数为 \underline{\hspace{4cm}}。

  \item 设 $L$ 为取逆时针方向的圆周 $x^2+y^2=9$，求曲线积分 \(\displaystyle \oint_L (2xy-2y)\,dx+(x^2-4x)\,dy\)
  的值为 \underline{\hspace{4cm}}。
\end{enumerate}

\paragraph*{三、计算题（每小题 7 分，共 35 分）}

\begin{enumerate}
  \item 将函数 \(\displaystyle f(x)=\ln\frac{1}{2+2x+x^2}\)
  按 $(x+1)$ 的正整数次幂展开成幂级数，即在 $x_0=-1$ 处展开成泰勒级数。

  \item 利用高斯公式计算曲面积分 \(\displaystyle \iint_{\Sigma}x^2\,dydz+2y^2\,dzdx+(3z^2-4x^2y^2)\,dxdy\)，
  其中 $\Sigma$ 为 $z=\sqrt{x^2+y^2}$ 与 $z=2$ 围成的立体表面，取外侧。

  \item 设二阶可微函数 $\varphi(x)$ 满足方程 \(\displaystyle \varphi(x)=e-\int_0^x (x-u)\varphi(u)\,du\)，
  求 $\varphi(x)$。

  \item 设 \(\displaystyle z=xf\left(2x,\frac{y^2}{x}\right)\)，
  其中 $f$ 具有二阶连续偏导数，求 $\dfrac{\partial^2 z}{\partial x\partial y}$。

  \item 设三重积分 \(\displaystyle I=\iiint_V (x^2+y^2+z^2)\,dxdydz\)，
  其中积分区域 $V$ 是由球面 $x^2+y^2+z^2=R^2$（$z\ge 0$）及锥面 $x^2+y^2=z^2$ 所围成的区域。
\end{enumerate}

\paragraph*{四、证明题（第 1 题 7 分，第 2 题 8 分，共 15 分）}

\begin{enumerate}
  \item 证明：一般项级数 \(\displaystyle \sum_{n=1}^{\infty}\sin\left(\pi\sqrt{n^2+1}\right)\)
  条件收敛。

  \item 证明：曲面 \(\displaystyle F\left(\frac{x-a}{z-c},\frac{y-b}{z-c}\right)=0\)
  上任意一点的切平面都通过一定点。
\end{enumerate}

\paragraph*{五、应用题（每小题 10 分，共 20 分）}

\begin{enumerate}
  \item 设函数 $y(x)$（$x\ge 0$）二阶可导，且 $y'(x)>0$，$y(0)=1$。
  过曲线 $y=y(x)$ 上任意一点 $P(x,y)$ 作该曲线的切线及 $x$ 轴的垂线，上述两直线与 $x$ 轴所围成的三角形面积记为 $S_1$。
  区间 $(0,x]$ 上以 $y=y(x)$ 为曲边的曲边梯形面积记为 $S_2$。
  若 $2S_1-S_2$ 恒为 $1$，求此曲线 $y=y(x)$ 的方程。

  \item 将长为 $a$ 的线段截成两段，一段围成一个正方形，另一段围成一个圆。问这两段的长度分别为多少时，它们所围成的正方形和圆的面积之和最大？
\end{enumerate}

\subsection{答案与解析}

\paragraph*{一、单项选择题}
\begin{enumerate}
  \item 由两平面法向量 $(1,3,2)$ 与 $(2,-1,-10)$ 叉乘，得直线方向向量可取 $(4,-2,1)$，它与平面 $\pi$ 的法向量 $(4,-2,1)$ 平行，故直线垂直于平面。选 C。
  \item 函数在原点连续，且两个偏导数均存在；但沿 $y=x$ 时 \(\displaystyle \frac{f(x,x)-f(0,0)}{\sqrt{x^2+x^2}}=\frac{x/2}{\sqrt2|x|}\)
  不趋于 $0$，故不可微。选 A。
  \item 直接积分可得
  \(\displaystyle \iint_Dxy\,dxdy=0,\qquad
  \iint_D\cos x\sin y\,dxdy=2\iint_{D_1}\cos x\sin y\,dxdy\)。
  选 A。
  \item A、B、C 均可由 Cauchy 不等式或比较判别推出收敛，D 与 C 矛盾，故不成立。选 D。
  \item 方程的积分因子为 $x$，于是 \(\displaystyle (xy)'=\sin x,\qquad xy=-\cos x+C\)。
  选 A。
\end{enumerate}

\paragraph*{二、填空题}
\begin{enumerate}
  \item 两平面的法向量为 $(1,0,2)$ 与 $(0,1,-3)$，直线方向向量可取 \(\displaystyle (1,0,2)\times(0,1,-3)=(-2,3,1)\)。因此直线方程为 \(\displaystyle \frac{x}{-2}=\frac{y-2}{3}=z-4\)。
  \item 积分区域为 $1\le x\le e,\ 0\le y\le\ln x$，即 \(\displaystyle 0\le y\le1,\ e^y\le x\le e\)。故 \(\displaystyle \int_1^e dx\int_0^{\ln x}f(x,y)\,dy=\int_0^1dy\int_{e^y}^e f(x,y)\,dx\)。
  \item 半径由 $3^n$ 项决定，为 $1/3$，中心为 $-1$。当 $x+1=-1/3$ 时为交错调和级数加收敛级数；当 $x+1=1/3$ 时含调和级数发散。收敛域为
  \(\displaystyle \left[-\frac43,-\frac23\right)\)。
  \item \(\displaystyle \nabla u=(yz,xz,xy),\quad \frac{\ell}{|\ell|}=\frac{(2,-1,3)}{\sqrt{14}}\)，
  故方向导数为 \(\displaystyle D_\ell u=\frac{2yz-xz+3xy}{\sqrt{14}}\)。
  \item 由 Green 公式，\(\displaystyle Q_x-P_y=(2x-4)-(2x-2)=-2\)。
  圆盘面积为 $9\pi$，故积分为 \(\displaystyle -18\pi\)。
\end{enumerate}

\paragraph*{三、计算题}
\begin{enumerate}
  \item 令 $t=x+1$，则 \(\displaystyle f(x)=\ln\frac1{1+t^2}=-\ln(1+t^2)
  =\sum_{n=1}^{\infty}\frac{(-1)^n}{n}t^{2n}\)。
  因此
  \(\displaystyle f(x)=\sum_{n=1}^{\infty}\frac{(-1)^n}{n}(x+1)^{2n},\qquad |x+1|<1\)。
  \item 设向量场
  \(\displaystyle \mathbf F=(x^2,2y^2,3z^2-4x^2y^2)\)。
  其散度为 \(\displaystyle \nabla\cdot\mathbf F=2x+4y+6z\)。圆锥体关于 $x,y$ 对称，$2x,4y$ 的积分为 $0$。在柱坐标下区域为 $0\le z\le2,\ 0\le r\le z$，故
  \[
  \iint_{\Sigma}\mathbf F\cdot d\mathbf S
  =\int_0^{2\pi}\int_0^2\int_0^z6z\,r\,dr\,dz\,d\theta
  =24\pi.
  \]
  \item 对方程 \(\displaystyle \varphi(x)=e-\int_0^x(x-u)\varphi(u)\,du\)
  求导，得 $\varphi'(x)=-\int_0^x\varphi(u)\,du$，再求导得
  \(\displaystyle \varphi''(x)=-\varphi(x),\qquad \varphi(0)=e,\quad \varphi'(0)=0\)。因此
  \(\displaystyle \varphi(x)=e\cos x\)。
  \item 设 $u=2x,\ v=y^2/x$。先对 $y$ 求导：
  \(\displaystyle z_y=x f_2(u,v)\frac{2y}{x}=2y f_2(u,v)\)。
  再对 $x$ 求导：\(\displaystyle z_{xy}=2y\left(2f_{21}(u,v)-\frac{y^2}{x^2}f_{22}(u,v)\right)\)。即
  \[
  \frac{\partial^2z}{\partial x\partial y}
  =4y f_{12}\left(2x,\frac{y^2}{x}\right)
  -\frac{2y^3}{x^2}f_{22}\left(2x,\frac{y^2}{x}\right).
  \]
  \item 用球坐标。锥面 $x^2+y^2=z^2$ 对应 $\varphi=\pi/4$，区域为
  \(\displaystyle 0\le\rho\le R,\qquad 0\le\varphi\le\frac{\pi}{4},\qquad 0\le\theta\le2\pi\)。
  因 $x^2+y^2+z^2=\rho^2$，故
  \[
  I=\int_0^{2\pi}\int_0^{\pi/4}\int_0^R\rho^4\sin\varphi\,d\rho d\varphi d\theta
  =\frac{\pi R^5}{5}(2-\sqrt2).
  \]
\end{enumerate}

\paragraph*{四、证明题}
\begin{enumerate}
  \item 记 \(\displaystyle \delta_n=\pi\left(\sqrt{n^2+1}-n\right)\sim\frac{\pi}{2n}\)。
  则 \(\displaystyle \sin\left(\pi\sqrt{n^2+1}\right)=(-1)^n\sin\delta_n\)。由交错级数判别法，原级数收敛；又
  \(\displaystyle \left|\sin\left(\pi\sqrt{n^2+1}\right)\right|\sim\frac{\pi}{2n}\)，绝对值级数发散，故原级数条件收敛。
  \item 设 \(\displaystyle u=\frac{x-a}{z-c},\qquad v=\frac{y-b}{z-c}\)。
  曲面由过定点 $(a,b,c)$ 的直线族
  \(\displaystyle (x,y,z)=(a,b,c)+t(u,v,1)\) 生成。经过曲面上一点并沿该母线的方向仍在曲面上，所以该母线方向属于该点切平面，切平面必通过定点 $(a,b,c)$。
\end{enumerate}

\paragraph*{五、应用题}
\begin{enumerate}
  \item 切线与 $x$ 轴交点到 $x$ 的水平距离为 $y/y'$，故
  \(\displaystyle S_1=\frac12y\cdot\frac{y}{y'}=\frac{y^2}{2y'},\qquad
  S_2=\int_0^x y(t)\,dt\)。
  条件 $2S_1-S_2=1$ 化为
  \(\displaystyle \frac{y^2}{y'}-\int_0^xy(t)\,dt=1\)。求导得 \(\displaystyle yy''=(y')^2\)。因此 $(y'/y)'=0$，即 $y=Ce^{kx}$。由 $y(0)=1$ 得 $C=1$；再代入 $x=0$ 的原条件得 $1/y'(0)=1$，故 $k=1$。所以
  \(\displaystyle y=e^x\)。
  \item 设围正方形的线段长为 $s$，围圆的线段长为 $a-s$。面积和为
  \(\displaystyle A(s)=\frac{s^2}{16}+\frac{(a-s)^2}{4\pi},\qquad 0\le s\le a\)。
  这是开口向上的二次函数，最大值在端点取得。比较
  \(\displaystyle A(0)=\frac{a^2}{4\pi},\qquad A(a)=\frac{a^2}{16}\)，因 $1/(4\pi)>1/16$，故最大时正方形段为 $0$，圆段为 $a$。若要求两段都为正，则没有最大值，只有上确界 $a^2/(4\pi)$。
\end{enumerate}

\section{2006--2007 第二学期 B 卷}

\subsection{试题}

\paragraph*{一、单项选择题（每小题 3 分，共 15 分）}

\begin{enumerate}
  \item 下列说法不正确的是（\quad）。

  A. 如果函数 $f(x,y)$ 在区域 $D$ 中连续，则 $|f(x,y)|$ 也在 $D$ 中连续；\quad
  B. 函数 $f(x,y)$ 沿射线 $\ell$ 的方向导数，等于梯度在 $\ell$ 上的投影；\quad
  C. 如果函数 $f(x,y)$ 在点 $P(x_0,y_0)$ 取极值，则必有 $f_x(x_0,y_0)=0,\ f_y(x_0,y_0)=0$；\quad
  D. 如果函数 $f(x,y)$ 在点 $P(x_0,y_0)$ 存在两个连续的偏导数，则在该点一定可微。

  \item 函数 $z=\ln(1+x^2+y^2)$ 当 $x=1,\ y=2$ 时的全微分为（\quad）。

  A. $\dfrac13dx+\dfrac13dy$；\quad
  B. $\dfrac23dx+\dfrac13dy$；\quad
  C. $\dfrac23dx+\dfrac23dy$；\quad
  D. $\dfrac13dx+\dfrac23dy$。

  \item 设 $\Omega$ 是曲面 $z=x^2+y^2$ 与平面 $z=1$ 所围成的闭区域，则 $\iiint_{\Omega}f(x,y,z)\,dV$ 化为三次积分为（\quad）。

  A. $\displaystyle\int_{-1}^{1}dy\int_{-\sqrt{1-y^2}}^{\sqrt{1-y^2}}dx\int_1^{x^2+y^2}f(x,y,z)\,dz$；

  B. $\displaystyle4\int_0^1dx\int_0^{\sqrt{1-x^2}}dy\int_{x^2+y^2}^{1}f(x,y,z)\,dz$；

  C. $\displaystyle4\int_0^1dx\int_0^{\sqrt{1-x^2}}dy\int_0^{x^2+y^2}f(x,y,z)\,dz$；

  D. $\displaystyle\int_{-1}^{1}dx\int_{-\sqrt{1-x^2}}^{\sqrt{1-x^2}}dy\int_{x^2+y^2}^{1}f(x,y,z)\,dz$。

  \item 设曲线 $L$ 为沿顺时针方向的圆周 $x^2+y^2=9$，则第二型闭路曲线积分 \(\displaystyle \oint_L (2xy-2y)\,dx+(x^2-4x)\,dy\)
  等于（\quad）。

  A. $0$；\quad B. $2\pi$；\quad C. $6\pi$；\quad D. $18\pi$。

  \item 已知幂级数 $\sum_{n=1}^{\infty}a_nx^n$ 的收敛域为 $[-8,8]$，则 \(\displaystyle \sum_{n=2}^{\infty}\frac{a_nx^n}{n(n-1)}\)
  的收敛半径为（\quad）。
  \textit{（校勘：原卷下标疑似为 $n=1$，但 $n(n-1)$ 使首项无定义；此处按 $n=2$ 处理。）}

  A. $R>8$；\quad B. $R<8$；\quad C. $R=8$；\quad D. 不能确定。
\end{enumerate}

\paragraph*{二、填空题（每小题 3 分，共 15 分）}

\begin{enumerate}
  \item 点 $(2,1,0)$ 到平面 $3x+4y+5z=0$ 的距离为 \underline{\hspace{4cm}}。

  \item 微分方程 \(\displaystyle \frac{1}{y}\frac{dy}{dx}=2x+\frac{x-x^3}{y}\)
  （$y\ne0$）的通解为 \underline{\hspace{4cm}}。

  \item 设 \(\displaystyle f(x,y)=\ln\left(x+\frac{y}{2x}\right)\)，
  则 $\left.\dfrac{\partial f}{\partial y}\right|_{x=1,\ y=0}=$ \underline{\hspace{4cm}}。

  \item 设
  \[
    f(x)=
    \begin{cases}
      x^2, & -2\le x<0,\\
      4-x^2, & 0\le x<2
    \end{cases}
  \]
  则 $f(x)$ 以 $4$ 为周期的傅里叶级数的和函数 $s(x)$ 在 $[-2,2)$ 上的表达式为 \underline{\hspace{4cm}}。

  \item 设曲线积分 \(\displaystyle \int_L [f(x)-e^x]\sin y\,dx-f(x)\cos y\,dy\)
  与积分路径无关，其中 $f(x)$ 具有一阶连续导数，且 $f(0)=0$，则 $f(x)=$ \underline{\hspace{4cm}}。
\end{enumerate}

\paragraph*{三、计算题（每小题 7 分，共 35 分）}

\begin{enumerate}
  \item 判断交错级数 \(\displaystyle \sum_{n=2}^{\infty}\frac{(-1)^n}{\sqrt n+(-1)^n}\)
  的敛散性。
  \textit{（校勘：原卷下标疑似为 $n=1$，但首项分母为 $0$；此处按 $n=2$ 处理。）}

  \item 利用高斯公式计算曲面积分 \(\displaystyle \iint_{\Sigma}x^2\,dydz+2y^2\,dzdx+(3z^2-4x^2y^2)\,dxdy\)，
  其中 $\Sigma$ 为 $z=\sqrt{x^2+y^2}$ 与 $z=2$ 围成的立体表面，取外侧。

  \item 已知 $y_1(x)=e^x$ 是齐次线性微分方程 \(\displaystyle (2x-1)y''-(2x+1)y'+2y=0\)
  的一个解，求此方程的通解。

  \item 设 \(\displaystyle z=xf\left(\frac{y}{x}\right)+yg\left(x,\frac{x}{y}\right)\)，
  其中 $f,g$ 均为二阶可微函数，求 $\dfrac{\partial^2 z}{\partial x\partial y}$。

  \item 计算二重积分 \(\displaystyle \iint_D (x^2+y^2)\,d\sigma\)，
  其中 $D$ 是由 $y=\sqrt{1-x^2}$ 和 $x$ 轴所围成的平面区域。
\end{enumerate}

\paragraph*{四、证明题（第 1 题 7 分，第 2 题 8 分，共 15 分）}

\begin{enumerate}
  \item 证明：级数 $\sum_{n=1}^{\infty}(a_{n+1}-a_n)$ 收敛的充分必要条件是数列 $\{a_n\}$ 收敛。

  \item 证明：曲面 \(\displaystyle F\left(\frac{x-a}{z-c},\frac{y-b}{z-c}\right)=0\)
  上任意一点的切平面都通过一定点。
\end{enumerate}

\paragraph*{五、应用题（每小题 10 分，共 20 分）}

\begin{enumerate}
  \item 设有力场 \(\displaystyle \mathbf F=\left(-\frac1y,\frac{x}{y^2}\right)\)，
  证明力 $\mathbf F$ 在上半平面内对质点所做的功与路径无关，并计算质点从点 $A(1,2)$ 移动到点 $B(2,1)$ 时力 $\mathbf F$ 所做的功。

  \item 设周长为 $2p$ 的矩形绕它的一边旋转构成圆柱形，求矩形的边长各为多少时，圆柱的体积最大。
\end{enumerate}

\subsection{答案与解析}

\paragraph*{一、单项选择题}
\begin{enumerate}
  \item C 不正确。极值点若偏导不存在，不能推出 $f_x=f_y=0$。选 C。
  \item \(\displaystyle dz=\frac{2x}{1+x^2+y^2}dx+\frac{2y}{1+x^2+y^2}dy\)，
  在 $(1,2)$ 处为 $\dfrac13dx+\dfrac23dy$。选 D。
  \item 区域满足 $x^2+y^2\le1,\ x^2+y^2\le z\le1$，故选 D。
  \item 对顺时针方向，Green 公式取负号。因
  \(\displaystyle Q_x-P_y=(2x-4)-(2x-2)=-2\)，
  逆时针积分为 $-18\pi$，顺时针积分为 $18\pi$。选 D。
  \item 除以 $n(n-1)$ 不改变幂级数收敛半径，故 $R=8$。选 C。
\end{enumerate}

\paragraph*{二、填空题}
\begin{enumerate}
  \item \(\displaystyle d=\frac{|3\cdot2+4\cdot1+5\cdot0|}{\sqrt{3^2+4^2+5^2}}=\sqrt2\)。
  \item 原方程化为
  \(\displaystyle y'=2xy+x-x^3\)。
  即 \(\displaystyle y'-2xy=x-x^3\)。取积分因子 $e^{-x^2}$，得
  \(\displaystyle (ye^{-x^2})'=(x-x^3)e^{-x^2}\)。积分后通解为
  \(\displaystyle y=\frac{x^2}{2}+Ce^{x^2}\)。
  \item \(\displaystyle f_y(x,y)=\frac{1}{x+y/(2x)}\cdot\frac1{2x}\)，
  故
  \(\displaystyle \left.f_y\right|_{(1,0)}=\frac12\)。
  \item 傅里叶级数在连续点收敛到函数值，在跳跃点收敛到左右极限平均值。因此
  \[
  s(x)=
  \begin{cases}
  x^2,&-2<x<0,\\
  4-x^2,&0<x<2,\\
  2,&x=-2,0.
  \end{cases}
  \]
  \item 令 \(\displaystyle P=(f(x)-e^x)\sin y,\qquad Q=-f(x)\cos y\)。
  路径无关要求 $P_y=Q_x$，即
  \(\displaystyle (f-e^x)\cos y=-f'(x)\cos y\)。故 $f'+f=e^x$，结合 $f(0)=0$ 得
  \(\displaystyle f(x)=\frac{e^x-e^{-x}}2\)。
\end{enumerate}

\paragraph*{三、计算题}
\begin{enumerate}
  \item 当 $n\ge2$ 时，
  \[
  \frac{(-1)^n}{\sqrt n+(-1)^n}
  =
  \frac{(-1)^n(\sqrt n-(-1)^n)}{n-1}
  =
  \frac{(-1)^n\sqrt n}{n-1}-\frac1{n-1}.
  \]
  第一项为交错型且趋于 $0$，第二项为负调和型，故级数发散到 $-\infty$。
  \item 与 A 卷同题。由 Gauss 公式，
  \(\displaystyle \iint_{\Sigma}x^2\,dydz+2y^2\,dzdx+(3z^2-4x^2y^2)\,dxdy=24\pi\)。
  \item 已知一解 $y_1=e^x$。令 $y=ve^x$，代入方程可化为
  \(\displaystyle (2x-1)v''+(2x-3)v'=0\)。
  设 $w=v'$，则
  \(\displaystyle (2x-1)w'+(2x-3)w=0\)，得 $w=C(2x-1)e^{-x}$，积分得另一解可取 $2x+1$。故通解为
  \(\displaystyle y=C_1e^x+C_2(2x+1)\)。
  \item 直接链式求导可得
  \[
  \frac{\partial^2z}{\partial x\partial y}
  =-\frac{y}{x^2}f''\left(\frac yx\right)
  +g_1\left(x,\frac xy\right)
  -\frac{x}{y}g_{12}\left(x,\frac xy\right)
  -\frac{x}{y^2}g_{22}\left(x,\frac xy\right).
  \]
  \item 区域为上半单位圆盘，取极坐标：
  \[
  \iint_D(x^2+y^2)\,d\sigma
  =\int_0^\pi\int_0^1 r^2\cdot r\,dr\,d\theta
  =\frac{\pi}{4}.
  \]
\end{enumerate}

\paragraph*{四、证明题}
\begin{enumerate}
  \item 部分和为 \(\displaystyle \sum_{n=1}^N(a_{n+1}-a_n)=a_{N+1}-a_1\)。
  因此级数收敛当且仅当 $\{a_{N+1}\}$ 收敛，也即 $\{a_n\}$ 收敛。
  \item 同 A 卷：曲面由过定点 $(a,b,c)$ 的母线族生成，母线方向属于切平面，所以任意切平面都通过 $(a,b,c)$。
\end{enumerate}

\paragraph*{五、应用题}
\begin{enumerate}
  \item 设 \(\displaystyle P=-\frac1y,\qquad Q=\frac{x}{y^2}\)。
  有 $P_y=Q_x=1/y^2$，上半平面单连通，故路径无关。势函数可取
  \(\displaystyle \phi(x,y)=-\frac{x}{y}\)。因此从 $A(1,2)$ 到 $B(2,1)$ 所做的功为
  \(\displaystyle \phi(B)-\phi(A)=-2-\left(-\frac12\right)=-\frac32\)。
  \item 设矩形一边为 $x$，另一边为 $p-x$。若绕边 $x$ 旋转，则体积
  \(\displaystyle V=\pi(p-x)^2x\)。
  求导得最大值在 $x=p/3$ 处取得，此时另一边为 $2p/3$。故矩形边长为
  \(\displaystyle \frac p3,\qquad \frac{2p}{3}\)，且绕较短边旋转时体积最大，最大体积为
  \(\displaystyle \frac{4\pi p^3}{27}\)。
\end{enumerate}

\section{2007--2008 第二学期 A 卷}

\subsection{试题}

\paragraph*{一、填空题（共 5 小题，每小题 3 分，共 15 分）}

\begin{enumerate}
  \item 函数 $z=\arctan\dfrac{y}{x}$ 在点 $(1,0)$ 处的梯度为 \underline{\hspace{3cm}}。

  \item 积分 \(\displaystyle \int_0^2 dx\int_x^2 e^{-y^2}\,dy=\underline{\hspace{3cm}}\)。

  \item 幂级数 \(\displaystyle \sum_{n=1}^{\infty}\frac{(x-2)^n}{n4^n}\)
  的收敛域为 \underline{\hspace{3cm}}。

  \item 设数量场 $u=\ln\sqrt{x^2+y^2}$，则 \(\displaystyle \operatorname{div}(\operatorname{grad}u)=\underline{\hspace{3cm}}\)。

  \item 函数 $z=xy$ 在条件 $x+y=1$ 下的极大值为 \underline{\hspace{3cm}}。
\end{enumerate}

\paragraph*{二、选择题（共 5 小题，每小题 3 分，共 15 分）}

\begin{enumerate}
  \item 抛物线 $y^2=4x$ 上与直线 $x-y+4=0$ 相距最近的点是（\quad）。

  A. $(0,0)$；\quad B. $(1,1)$；\quad C. $(1,2)$；\quad D. $(2,2\sqrt2)$。

  \item 设区域 $D:0\le y\le x,\ 0\le x\le \pi$，则 \(\displaystyle \iint_D \sqrt{1-\cos^2 x}\,dxdy\)
  等于（\quad）。

  A. $0$；\quad B. $\dfrac{\pi}{2}$；\quad C. $2$；\quad D. $\pi$。

  \item 设 $\alpha$ 是常数，则级数 \(\displaystyle \sum_{n=1}^{\infty}\left[\frac{\sin(n\alpha)}{n^2}-\frac1{\sqrt n}\right]\)
  是（\quad）。

  A. 绝对收敛；\quad B. 发散；\quad C. 条件收敛；\quad D. 敛散性与 $\alpha$ 取值有关。

  \item 若 $0\le a_n<\dfrac1n\ (n=1,2,\cdots)$，则下列级数中肯定收敛的是（\quad）。

  A. $\sum_{n=1}^{\infty}a_n$；\quad
  B. $\sum_{n=1}^{\infty}(-1)^na_n$；\quad
  C. $\sum_{n=1}^{\infty}\sqrt{a_n}$；\quad
  D. $\sum_{n=1}^{\infty}(-1)^na_n^2$。

  \item 设 $y_1$ 与 $y_2$ 是非齐次线性微分方程 \(\displaystyle y''+ay'+by=f(x)\)
  的两个特解，则下列说法正确的是（\quad）。

  A. $y_1+y_2$ 是非齐次方程的解；\quad
  B. $y_1-y_2$ 是非齐次方程的解；\quad
  C. $y_1+y_2$ 是对应齐次方程的解；\quad
  D. $y_1-y_2$ 是对应齐次方程的解。
\end{enumerate}

\paragraph*{三、计算题（共 2 小题，每小题 7 分，共 14 分）}

\begin{enumerate}
  \item 设 $u(x,y,z)=x+yz$，其中 $z=z(x,y)$ 由 $x+y+z=0$ 确定，求 \(\displaystyle \frac{\partial^2u}{\partial x\partial y}\)。

  \item 求曲面 $x^2+2y^2+3z^2=21$ 上平行于平面 $x+4y+6z=0$ 的各个切平面。
\end{enumerate}

\paragraph*{四、计算题（共 2 小题，每小题 7 分，共 14 分）}

\begin{enumerate}
  \item 计算曲线积分 \(\displaystyle \int_L xy\,dx+(y-x)\,dy\)，
  其中 $L$ 为 $y=x^2$ 上由点 $O(0,0)$ 到点 $A(1,1)$ 的一段。

  \item 设 $\Sigma$ 是锥面 $z=\sqrt{x^2+y^2}$ 被平面 $z=1$ 所截下的有限部分下侧，计算 \(\displaystyle \iint_{\Sigma}x\,dydz+y\,dzdx+(4-2z)\,dxdy\)。
\end{enumerate}

\paragraph*{五、证明题（8 分）}

将函数 $f(x)=x$ 在 $[0,\pi]$ 上展开成余弦级数，并证明
\(\displaystyle \sum_{n=1}^{\infty}\frac1{(2n-1)^2}=\frac{\pi^2}{8}\)。

\paragraph*{六、计算题（8 分）}

求幂级数 \(\displaystyle \sum_{n=1}^{\infty}n x^{2n}\) 的收敛域及和函数。

\paragraph*{七、证明题（8 分）}

设 $f(x)\in C[0,2]$，证明
\(\displaystyle \int_0^2 dx\int_x^2 f(x)f(y)\,dy
  =\frac12\left[\int_0^2 f(x)\,dx\right]^2\)。

\paragraph*{八、证明题（8 分）}

证明函数 \(\displaystyle F(y)=\int_1^{+\infty}\frac{dx}{\sqrt[3]{x^4+y^2}}\) 在 $-\infty<y<+\infty$ 上连续。

\paragraph*{九、应用题（10 分）}

设 \(\displaystyle f(x)=\sin x-\int_0^x (x-t)f(t)\,dt\)，其中 $f(x)$ 为连续函数，求 $f(x)$。

\subsection{答案与解析}

\paragraph*{一、填空题}

\begin{enumerate}
  \item $\operatorname{grad}z(1,0)=(0,1)$。
  \item 交换积分次序：
  \[
    \int_0^2dx\int_x^2e^{-y^2}\,dy
    =
    \int_0^2dy\int_0^y e^{-y^2}\,dx
    =
    \int_0^2ye^{-y^2}\,dy
    =
    \frac{1-e^{-4}}{2}
  \]
  。
  \textit{（校勘：答案版给出的数值对应 $e^{y^2}$，但试题 PDF 图像为 $e^{-y^2}$；此处按题面计算。）}
  \item 令 $t=\dfrac{x-2}{4}$，则级数为 $\displaystyle\sum_{n=1}^\infty \frac{t^n}{n}$。当 $t=-1$ 时收敛，当 $t=1$ 时发散，故收敛域为 $[-2,6)$。
  \item $\operatorname{div}(\operatorname{grad}u)=0$（在 $(x,y)\ne(0,0)$ 上）。
  \item 极大值为 $\dfrac14$。
\end{enumerate}

\paragraph*{二、选择题}

\begin{enumerate}
  \item 设抛物线上点为 $\left(\dfrac{y^2}{4},y\right)$，到直线的距离与
  \[
    \frac{y^2}{4}-y+4=\frac{(y-2)^2}{4}+3
  \]
  同时取最小。故 $y=2,\ x=1$，选 C。
  \item 在 $0\le x\le\pi$ 上有 $\sqrt{1-\cos^2x}=\sin x$，故
  \[
    \iint_D\sqrt{1-\cos^2x}\,dxdy=\int_0^\pi x\sin x\,dx=\pi.
  \]
  选 D。
  \item 因 $\displaystyle\sum_{n=1}^\infty \left|\frac{\sin(n\alpha)}{n^2}\right|$ 收敛，而 $\displaystyle\sum_{n=1}^\infty\frac1{\sqrt n}$ 发散，所以原级数发散，选 B。
  \item 由 $0\le a_n<\dfrac1n$ 得 $0\le a_n^2<\dfrac1{n^2}$，所以 $\displaystyle\sum(-1)^na_n^2$ 绝对收敛，选 D。
  \item 设线性算子 $L[y]=y''+ay'+by$。因 $L[y_1]=L[y_2]=f(x)$，故 $L[y_1-y_2]=0$，所以 $y_1-y_2$ 是对应齐次方程的解，选 D。
\end{enumerate}

\paragraph*{三、计算题}

\begin{enumerate}
  \item 由 $x+y+z=0$ 得 $z_x=z_y=-1$，$z_{xy}=0$。又 $u=x+yz$，故
  \[
    \frac{\partial^2u}{\partial x\partial y}
    =
    \frac{\partial}{\partial y}\left(1+yz_x\right)
    =
    \frac{\partial}{\partial y}(1-y)=-1.
  \]

  \item 设切点为 $(x_0,y_0,z_0)$。椭球面法向量为
  \[
    (2x_0,4y_0,6z_0),
  \]
  与平面 $x+4y+6z=0$ 的法向量 $(1,4,6)$ 平行，因此
  \(\displaystyle (2x_0,4y_0,6z_0)=\lambda(1,4,6)\)。代入 $x_0^2+2y_0^2+3z_0^2=21$ 得 $\lambda=\pm2$，故切平面为
  \(\displaystyle x+4y+6z=\pm 21\)。
\end{enumerate}

\paragraph*{四、计算题}

\begin{enumerate}
  \item 令 $y=x^2$，$0\le x\le 1$，则 $dy=2x\,dx$，所以
  \[
    \int_L xy\,dx+(y-x)\,dy
    =
    \int_0^1\left(x^3+(x^2-x)2x\right)\,dx
    =
    \int_0^1(3x^3-2x^2)\,dx
    =
    \frac1{12}.
  \]

  \item 添上顶面 $\Sigma_0:z=1,\ x^2+y^2\le1$ 组成闭曲面。向量场 \(\displaystyle \mathbf F=(x,y,4-2z)\)
  的散度为 $1+1-2=0$，故闭曲面积分为 $0$。底面外法向量为 $(0,0,1)$，底面积分为
  \(\displaystyle \iint_{\Sigma_0}(4-2)\,dxdy=2\pi\)。因此锥面下侧部分积分为 $-2\pi$。
\end{enumerate}

\paragraph*{五、证明题}

将 $f(x)=x$ 偶延拓到 $[-\pi,\pi]$，其余弦级数为
\[
  x\sim \frac{\pi}{2}
  +\sum_{n=1}^{\infty}\frac{2\big((-1)^n-1\big)}{\pi n^2}\cos nx
  =
  \frac{\pi}{2}
  -\frac{4}{\pi}\sum_{n=1}^{\infty}\frac{\cos(2n-1)x}{(2n-1)^2}.
\]
取 $x=0$ 得
\[
  0=\frac{\pi}{2}
  -\frac{4}{\pi}\sum_{n=1}^{\infty}\frac1{(2n-1)^2},
  \qquad
  \sum_{n=1}^{\infty}\frac1{(2n-1)^2}=\frac{\pi^2}{8}.
\]

\paragraph*{六、计算题}

当 $|x|<1$ 时，令 $u=x^2$，有
\[
  \sum_{n=1}^{\infty}nx^{2n}
  =
  \sum_{n=1}^{\infty}nu^n
  =
  \frac{u}{(1-u)^2}
  =
  \frac{x^2}{(1-x^2)^2}.
\]
在 $x=\pm1$ 时级数发散，故收敛域为 $(-1,1)$。

\paragraph*{七、证明题}

记 \(\displaystyle I=\int_0^2 dx\int_x^2 f(x)f(y)\,dy\)。由积分区域关于直线 $x=y$ 对称，交换 $x,y$ 得
\[
I=\int_0^2 dy\int_0^y f(x)f(y)\,dx.
\]
两式相加，覆盖正方形 $[0,2]\times[0,2]$，因此
\[
2I=\int_0^2\int_0^2 f(x)f(y)\,dxdy=\left[\int_0^2 f(x)\,dx\right]^2,
\]
命题得证。

\paragraph*{八、证明题}

任取 $N>0$。当 $y\in[-N,N]$ 且 $x\ge1$ 时，\(\displaystyle 0\le \frac1{\sqrt[3]{x^4+y^2}}\le \frac1{\sqrt[3]{x^4}}=\frac1{x^{4/3}}\)，而 $\int_1^{+\infty}x^{-4/3}\,dx$ 收敛。因此由含参广义积分的一致收敛判别，$F(y)$ 在 $[-N,N]$ 上连续。由于 $N$ 任意，$F(y)$ 在 $(-\infty,+\infty)$ 上连续。

\paragraph*{九、应用题}

由 \(\displaystyle f(x)=\sin x-\int_0^x(x-t)f(t)\,dt\) 对 $x$ 求导，得
\(\displaystyle f'(x)=\cos x-\int_0^x f(t)\,dt\)。再求导得 \(\displaystyle f''(x)=-\sin x-f(x)\)。即
\(\displaystyle f''+f=-\sin x,\qquad f(0)=0,\quad f'(0)=1\)。解得
\(\displaystyle f(x)=\frac12\sin x+\frac{x}{2}\cos x\)。

\section{2009级工科数学分析下 A卷}

\subsection{资料来源}
\begin{itemize}
  \item 试题：2009级计算机工科数学分析下 期末A.doc，DOC 正文已抽取；公式对象已按 Word 公式图片逐项恢复。
  \item 公式图片审计：\texttt{tmp/past-exam-equations/2009-A/manifest.csv}。
\end{itemize}

\subsection{试题}
诚信应考，考试作弊将带来严重后果！

华南理工大学 2009--2010 第二学期期末考试

《工科数学分析》（下）试卷（A）

注意事项：
\begin{enumerate}
  \item 考前请将密封线内填写清楚；
  \item 所有答案请直接答在试卷上；
  \item 考试形式：闭卷；
  \item 本试卷共五大题，满分 100 分，考试时间 120 分钟。
\end{enumerate}

\subsubsection*{一、填空题（共 5 小题，每小题 4 分，共 20 分）}
\begin{enumerate}
  \item $(y''')^3+(y')^4+x^2+x+1=0$ 是 \underline{\hspace{2cm}} 阶微分方程。
  \item 已知 $(a\times b)\cdot c=2$，则 $[(a+b)\times(b+c)]\cdot(c+a)=$ \underline{\hspace{3cm}}。
  \item 函数 $z=\arctan\dfrac{y}{x}+\ln\sqrt{x^2+y^2}$ 在点 $(1,0)$ 处的梯度 $\operatorname{grad}z\big|_{(1,0)}=$ \underline{\hspace{3cm}}。
  \item 交换积分顺序，则 $\displaystyle\int_1^e dx\int_0^{\ln x}f(x,y)\,dy=$ \underline{\hspace{3cm}}。
  \item 设 $f(x)$ 是以 $2\pi$ 为周期的函数，它在 $[-\pi,\pi]$ 上的表达式是
  \[
  f(x)=
  \begin{cases}
  x+1,&-\pi\le x<0,\\
  x^2,&0\le x<\pi.
  \end{cases}
  \]
  若它的傅里叶级数的和函数为 $S(x)$，则 $S(0)=$ \underline{\hspace{3cm}}。
\end{enumerate}

\subsubsection*{二、选择题（共 5 小题，每小题 4 分，共 20 分）}
\begin{enumerate}
  \item 抛物线 $y^2=4x$ 上与直线 $x-y+4=0$ 相距最近的点是（\quad）。

  A. $(0,0)$；\quad B. $(1,1)$；\quad C. $(1,2)$；\quad D. $(2,2\sqrt2)$。

  \item 直线 $L:\dfrac{x-2}{2}=\dfrac{y+1}{-2}=\dfrac{z-3}{1}$ 与平面 $\pi:x+2y-2z=6$ 的关系是（\quad）。

  A. 平行；\quad B. 垂直；\quad C. 相交但不垂直；\quad D. 重合。

  \item 设 $f(x,y)$ 在点 $(x_0,y_0)$ 处的两个偏导数都存在，则（\quad）。

  A. $f(x,y)$ 在点 $(x_0,y_0)$ 必连续；\quad
  B. $f(x,y)$ 在点 $(x_0,y_0)$ 必可微；\quad
  C. $\displaystyle\lim_{x\to x_0}f(x,y_0)$ 与 $\displaystyle\lim_{y\to y_0}f(x_0,y)$ 都存在；\quad
  D. $\displaystyle\lim_{\substack{x\to x_0\\y\to y_0}}f(x,y)$ 存在。

  \item 设区域 $D:0\le y\le x,\ 0\le x\le\pi$，则二重积分 $\displaystyle\iint_D\sqrt{1-\sin^2x}\,dxdy$（\quad）。

  A. $0$；\quad B. $\dfrac{\pi}{2}$；\quad C. $2$；\quad D. $\pi$。

  \item 设 $\alpha$ 为常数，则级数 $\displaystyle\sum_{n=1}^{\infty}\left[\dfrac{\sin(n\alpha)}{n^2}-\dfrac1n\right]$（\quad）。

  A. 绝对收敛；\quad B. 发散；\quad C. 条件收敛；\quad D. 收敛性与 $\alpha$ 的取值有关。
\end{enumerate}

\subsubsection*{三、计算题（本大题共 4 小题，每小题 8 分，共 32 分）}
\begin{enumerate}
  \item 求微分方程 $xy''=y'\ln\dfrac{y'}{x}$ 的通解。
  \item 计算三重积分 $\displaystyle I=\iiint_{\Omega}(x^2+y^2+z^2)\,dV$，其中 $\Omega:x^2+y^2+z^2\le 2z$。
  \item 计算曲线积分 \(\displaystyle I=\int_L (x+y)^2\,dx-(x^2+y^2\sin y)\,dy\)，
  其中 $L$ 是抛物线 $y=x^2$ 上从点 $(-1,1)$ 到点 $(1,1)$ 的那一段。
  \item 求幂级数 $\displaystyle\sum_{n=0}^{\infty}(2n+1)x^{2n}$ 的收敛域及和函数。
\end{enumerate}

\subsubsection*{四、应用题（本大题共 2 小题，每小题 7 分，共 14 分）}
\begin{enumerate}
  \item 造一个容积为 $V$ 的长方形无盖铝盒，怎样设计尺寸才能使它的表面积最小？
  \item 求密度 $\rho\equiv1$ 的均匀球壳 $x^2+y^2+z^2=a^2\ (z\ge0)$ 对于 $Oz$ 轴的转动惯量。
\end{enumerate}

\subsubsection*{五、证明题（本大题共 2 小题，每小题 7 分，共 14 分）}
\begin{enumerate}
  \item 设函数 $f(x)$ 在区间 $[a,b]$ 上连续，证明：
  \(\displaystyle \int_a^b dx\int_a^x f(y)\,dy=\int_a^b f(y)(b-y)\,dy\)。
  \item 设正数序列 $\{x_n\}$ 单调上升且有界，证明级数 \(\displaystyle \sum_{n=1}^{\infty}\left(1-\frac{x_n}{x_{n+1}}\right)\)
  收敛。
\end{enumerate}

\subsection{答案与解析}

\subsubsection*{一、填空题}
\begin{enumerate}
  \item 最高阶导数为 $y'''$，故为 $3$ 阶微分方程。
  \item 展开得 \(\displaystyle [(a+b)\times(b+c)]\cdot(c+a)
  =(a\times b)\cdot c+(b\times c)\cdot a=2+2=4\)。
  \item \(\displaystyle z_x=\frac{x-y}{x^2+y^2},\qquad z_y=\frac{x+y}{x^2+y^2}\)。
  在 $(1,0)$ 处 $\operatorname{grad}z=(1,1)$。
  \item 积分区域为 $1\le x\le e,\ 0\le y\le\ln x$，即 $0\le y\le1,\ e^y\le x\le e$。因此
  \(\displaystyle \int_1^e dx\int_0^{\ln x}f(x,y)\,dy
  =\int_0^1dy\int_{e^y}^e f(x,y)\,dx\)。
  \item $x=0$ 是跳跃点，傅里叶级数在该点收敛到左右极限平均值：
  \(\displaystyle S(0)=\frac{f(0-)+f(0+)}2=\frac{1+0}{2}=\frac12\)。
\end{enumerate}

\subsubsection*{二、选择题}
\begin{enumerate}
  \item 令 $y=t,\ x=t^2/4$，到直线的距离只需最小化 \(\displaystyle \left(\frac{t^2}{4}-t+4\right)^2\)。
  在选项中最小点为 $(1,2)$，故选 C。
  \item 直线方向向量 $v=(2,-2,1)$，平面法向量 $n=(1,2,-2)$。因 $v\cdot n=-4\ne0$，直线与平面相交；且 $v$ 不平行于 $n$，故不垂直于平面。选 C。
  \item 偏导数存在只保证沿坐标方向的一元极限存在，不保证连续、可微或二重极限存在。选 C。
  \item \(\displaystyle \iint_D\sqrt{1-\sin^2x}\,dxdy
  =\int_0^\pi x|\cos x|\,dx=\pi\)。
  故选 D。
  \item $\sum\sin(n\alpha)/n^2$ 绝对收敛，而 $-\sum1/n$ 发散，所以原级数发散。选 B。
\end{enumerate}

\subsubsection*{三、计算题}
\begin{enumerate}
  \item 设 $p=y'$，则 \(\displaystyle xp'=p\ln\frac{p}{x}\)。
  令 $q=\ln\dfrac{p}{x}$，有 $p=xe^q$，代入得
  \(\displaystyle 1+xq'=q,\qquad q=1+C_1x\)。因而 $y'=ex e^{C_1x}$。故
  \(\displaystyle y=C_2+e\int x e^{C_1x}\,dx\)。即当 $C_1\ne0$ 时
  \(\displaystyle y=C_2+\frac{e\,e^{C_1x}(C_1x-1)}{C_1^2}\)；当 $C_1=0$ 时
  \(\displaystyle y=C_2+\frac e2x^2\)。
  \item 区域为球心 $(0,0,1)$、半径 $1$ 的球。令 $z=w+1$，则
  \(\displaystyle x^2+y^2+z^2=x^2+y^2+w^2+1+2w\)。
  对称性使 $2w$ 的积分为 $0$，所以
  \[
  I=\iiint_{r\le1}r^2\,dV+\iiint_{r\le1}1\,dV
  =\frac{4\pi}{5}+\frac{4\pi}{3}
  =\frac{32\pi}{15}.
  \]
  \item 取参数 $x=t,\ y=t^2,\ -1\le t\le1$，则 $dx=dt,\ dy=2t\,dt$。于是
  \[
  I=\int_{-1}^{1}\left[(t+t^2)^2-2t(t^2+t^4\sin t^2)\right]dt
  =\int_{-1}^{1}(t^2+t^4)\,dt=\frac{16}{15}.
  \]
  \item 令 $r=x^2$，则
  \[
  \sum_{n=0}^{\infty}(2n+1)x^{2n}
  =2\sum_{n=0}^{\infty}nr^n+\sum_{n=0}^{\infty}r^n
  =\frac{1+r}{(1-r)^2}.
  \]
  因此收敛域为 $(-1,1)$，和函数为
  \(\displaystyle S(x)=\frac{1+x^2}{(1-x^2)^2}\)。
\end{enumerate}

\subsubsection*{四、应用题}
\begin{enumerate}
  \item 设底面长宽为 $x,y$，高为 $h$，则 $xyh=V$，表面积
  \[
  S=xy+2xh+2yh.
  \]
  极小值由对称性取 $x=y=a$。于是 $h=V/a^2$，
  \(\displaystyle S=a^2+\frac{4V}{a},\qquad S'=2a-\frac{4V}{a^2}\)。得 $a^3=2V$。故底面为正方形，边长
  \[
  a=(2V)^{1/3},\qquad h=\frac a2=\left(\frac V4\right)^{1/3}.
  \]
  \item 在上半球面上取球坐标，$x^2+y^2=a^2\sin^2\varphi$，$dS=a^2\sin\varphi\,d\varphi d\theta$。于是
  \[
  I_{Oz}=\iint_\Sigma(x^2+y^2)\,dS
  =2\pi a^4\int_0^{\pi/2}\sin^3\varphi\,d\varphi
  =\frac{4\pi a^4}{3}.
  \]
\end{enumerate}

\subsubsection*{五、证明题}
\begin{enumerate}
  \item 左端积分区域为 $a\le y\le x\le b$。交换积分次序：
  \[
  \int_a^b dx\int_a^x f(y)\,dy
  =\int_a^b f(y)\,dy\int_y^b dx
  =\int_a^b f(y)(b-y)\,dy.
  \]
  \item 因 $\{x_n\}$ 为正数且单调上升，有 $x_{n+1}\ge x_1>0$。于是
  \(\displaystyle 0\le1-\frac{x_n}{x_{n+1}}
  =\frac{x_{n+1}-x_n}{x_{n+1}}
  \le\frac{x_{n+1}-x_n}{x_1}\)。
  部分和满足
  \(\displaystyle \sum_{n=1}^N\left(1-\frac{x_n}{x_{n+1}}\right)\le\frac{x_{N+1}-x_1}{x_1}\)，右端有界，故正项级数收敛。
\end{enumerate}

\section{2009级工科数学分析下 B卷}

\subsection{资料来源}
\begin{itemize}
  \item 试题：2009级计算机工科数学分析下 期末B.doc，DOC 正文已抽取；公式对象已按 Word 公式图片逐项恢复。
  \item 公式图片审计：\texttt{tmp/past-exam-equations/2009-B/manifest.csv}。
\end{itemize}

\subsection{试题}
诚信应考，考试作弊将带来严重后果！

华南理工大学 2009--2010 第二学期期末考试

《工科数学分析》（下）试卷（B）

注意事项：
\begin{enumerate}
  \item 考前请将密封线内填写清楚；
  \item 所有答案请直接答在试卷上；
  \item 考试形式：闭卷；
  \item 本试卷共五大题，满分 100 分，考试时间 120 分钟。
\end{enumerate}

\subsubsection*{一、填空题（共 5 小题，每小题 4 分，共 20 分）}
\begin{enumerate}
  \item 方程 $y''+5y'+4=0$ 的通解为 \underline{\hspace{4cm}}。
  \item 原点到平面 $x+\dfrac{y}{-4}+\dfrac{z}{3}=1$ 的距离为 \underline{\hspace{3cm}}。
  \item 函数 $z=\arctan\dfrac{y}{x}+\ln\sqrt{x^2+y^2}$ 在点 $(1,0)$ 处的梯度 $\operatorname{grad}z\big|_{(1,0)}=$ \underline{\hspace{3cm}}。
  \item 设 $L$ 是 $xOy$ 平面上任意一条分段光滑的闭曲线，则第二类曲线积分
  \(\displaystyle \oint_L(2xy+\cos x)\,dx+(x^2+\sin y)\,dy=\underline{\hspace{3cm}}\)。
  \item 设 $f(x)$ 是以 $2\pi$ 为周期的函数，它在 $[-\pi,\pi]$ 上的表达式是
  \[
  f(x)=
  \begin{cases}
  x+1,&-\pi\le x<0,\\
  x^2,&0\le x<\pi.
  \end{cases}
  \]
  若它的傅里叶级数的和函数为 $S(x)$，则 $S(0)=$ \underline{\hspace{3cm}}。
\end{enumerate}

\subsubsection*{二、选择题（共 5 小题，每小题 4 分，共 20 分）}
\begin{enumerate}
  \item 在 $R^3$ 中，方程 $x^2=4y$ 的图形是（\quad）。

  A. 抛物线；\quad B. 抛物柱面；\quad C. 椭圆抛物面；\quad D. 旋转抛物面。

  \item 曲面 $3x^2+y^2+z^2=12$ 上点 $M(-1,0,3)$ 处的切平面与平面 $z=0$ 的夹角是（\quad）。

  A. $\dfrac{\pi}{6}$；\quad B. $\dfrac{\pi}{4}$；\quad C. $\dfrac{\pi}{3}$；\quad D. $\dfrac{\pi}{2}$。

  \item 设 $f(x,y)$ 在点 $(x_0,y_0)$ 处的两个偏导数都存在，则（\quad）。

  A. $f(x,y)$ 在点 $(x_0,y_0)$ 必连续；\quad
  B. $f(x,y)$ 在点 $(x_0,y_0)$ 必可微；\quad
  C. $\displaystyle\lim_{x\to x_0}f(x,y_0)$ 与 $\displaystyle\lim_{y\to y_0}f(x_0,y)$ 都存在；\quad
  D. $\displaystyle\lim_{\substack{x\to x_0\\y\to y_0}}f(x,y)$ 存在。

  \item 设平面曲线 $L$ 为下半圆周 $y=-\sqrt{1-x^2}$，则曲线积分 $\displaystyle\int_L(x^2+y^2)\,ds$（\quad）。

  A. $0$；\quad B. $-1$；\quad C. $2$；\quad D. $\pi$。

  \item 幂级数 $\displaystyle\sum_{n=1}^{\infty}\frac{\ln(n+1)}{n}x^n$ 的收敛域为（\quad）。

  A. $(-1,1)$；\quad B. $[-1,1]$；\quad C. $[-1,1)$；\quad D. $(-1,1]$。
\end{enumerate}

\subsubsection*{三、计算题（本大题共 4 小题，每小题 8 分，共 32 分）}
\begin{enumerate}
  \item 求微分方程 $y''+y'+2y=0$ 的通解。
  \item 计算二重积分 $\displaystyle I=\iint_D\frac{x}{x^2+y^2}\,dxdy$，其中 $D$ 由曲线 $y=\dfrac{x^2}{2}$ 与直线 $y=x$ 围成。
  \item 计算曲线积分 \(\displaystyle I=\int_L (x+y)^2\,dx-(x^2+y^2\sin y)\,dy\)，
  其中 $L$ 是抛物线 $y=x^2$ 上从点 $(-1,1)$ 到点 $(1,1)$ 的那一段。
  \item 求幂级数 $\displaystyle\sum_{n=0}^{\infty}\frac{n}{n+1}x^n$ 的收敛域及和函数。
\end{enumerate}

\subsubsection*{四、应用题（本大题共 2 小题，每小题 7 分，共 14 分）}
\begin{enumerate}
  \item 欲造一个无盖的长方体容器，已知底部造价为 $3$ 元/平方米，侧面造价为 $1$ 元/平方米。现想用 36 元造一个容积最大的容器，问应该怎样设计尺寸？
  \item 设 $\mathbf F(x,y)=xy^2\mathbf i+yf(x)\mathbf j$ 为保守场，其中 $f(x)$ 具有连续导数，且 $f(0)=0$。求曲线积分
  \(\displaystyle I=\int_{(0,0)}^{(1,1)}\mathbf F\cdot d\mathbf s\)。
\end{enumerate}

\subsubsection*{五、证明题（本大题共 2 小题，每小题 7 分，共 14 分）}
\begin{enumerate}
  \item 证明：圆柱螺线 $x=a\cos t,\ y=a\sin t,\ z=bt$ 的切线与 $z$ 轴的夹角为常数。
  \item 若正项级数 $\displaystyle\sum_{n=1}^{\infty}a_n$ 收敛，证明正项级数 \(\displaystyle \sum_{n=1}^{\infty}\frac{\sqrt{a_n}}{n}\)
  也收敛。
\end{enumerate}

\subsection{答案与解析}

\subsubsection*{一、填空题}
\begin{enumerate}
  \item 令 $p=y'$，则 $p'+5p=-4$。解得
  \(\displaystyle p=C_1e^{-5x}-\frac45\)，
  故通解为
  \(\displaystyle y=C_2+C_1e^{-5x}-\frac45x\)。
  \item 平面化为 $x-\dfrac14y+\dfrac13z-1=0$，原点到平面的距离
  \(\displaystyle d=\frac1{\sqrt{1+\frac1{16}+\frac1{9}}}=\frac{12}{13}\)。
  \item 与 A 卷同题，
  \(\displaystyle \operatorname{grad}z\big|_{(1,0)}=(1,1)\)。
  \item 由格林公式，
  \[
  \frac{\partial}{\partial x}(x^2+\sin y)-\frac{\partial}{\partial y}(2xy+\cos x)=2x-2x=0.
  \]
  对任意闭曲线积分为 $0$。
  \item $x=0$ 是跳跃点，故
  \(\displaystyle S(0)=\frac{f(0-)+f(0+)}2=\frac{1+0}{2}=\frac12\)。
\end{enumerate}

\subsubsection*{二、选择题}
\begin{enumerate}
  \item 方程与 $z$ 无关，是沿 $z$ 方向延伸的抛物柱面，选 B。
  \item 曲面法向量为 $\nabla(3x^2+y^2+z^2)=(-6,0,6)$，平面 $z=0$ 的法向量为 $(0,0,1)$。两平面夹角满足
  \(\displaystyle \cos\theta=\frac{|(-6,0,6)\cdot(0,0,1)|}{6\sqrt2}=\frac1{\sqrt2}\)，
  故 $\theta=\dfrac{\pi}{4}$，选 B。
  \item 偏导数存在只保证沿坐标方向的一元极限存在。选 C。
  \item 下半圆周是单位圆的一半，且 $x^2+y^2=1$，故 \(\displaystyle \int_L(x^2+y^2)\,ds=\int_L1\,ds=\pi\)。
  选 D。
  \item 半径为 $1$。$x=1$ 时为 $\sum\ln(n+1)/n$ 发散；$x=-1$ 时为交错级数 $\sum(-1)^n\ln(n+1)/n$，因 $\ln(n+1)/n\downarrow0$ 而收敛。故收敛域为 $[-1,1)$，选 C。
\end{enumerate}

\subsubsection*{三、计算题}
\begin{enumerate}
  \item 特征方程 \(\displaystyle r^2+r+2=0,\qquad r=\frac{-1\pm i\sqrt7}{2}\)。
  故通解为
  \(\displaystyle y=e^{-x/2}\left(C_1\cos\frac{\sqrt7}{2}x+C_2\sin\frac{\sqrt7}{2}x\right)\)。
  \item 区域为 $0\le x\le2,\ x^2/2\le y\le x$。于是
  \[
  I=\int_0^2\int_{x^2/2}^{x}\frac{x}{x^2+y^2}\,dy\,dx
  =\int_0^2\left[\arctan\frac{y}{x}\right]_{x^2/2}^{x}dx.
  \]
  因此 \(\displaystyle I=\int_0^2\left(\frac{\pi}{4}-\arctan\frac{x}{2}\right)dx
  =\ln2\)。
  \item 与 A 卷同题，取 $x=t,\ y=t^2,\ -1\le t\le1$，得
  \(\displaystyle I=\int_{-1}^{1}(t^2+t^4)\,dt=\frac{16}{15}\)。
  \item 当 $|x|<1$，
  \[
  \sum_{n=0}^{\infty}\frac{n}{n+1}x^n
  =\sum_{n=0}^{\infty}x^n-\sum_{n=0}^{\infty}\frac{x^n}{n+1}
  =\frac1{1-x}+\frac{\ln(1-x)}{x}.
  \]
  端点 $x=1$ 时发散，$x=-1$ 时收敛，故收敛域为 $[-1,1)$。和函数为
  \[
  S(x)=\frac1{1-x}+\frac{\ln(1-x)}{x}\quad(-1<x<1),
  \]
  并有 $S(0)=0$，$S(-1)=\dfrac12-\ln2$。
\end{enumerate}

\subsubsection*{四、应用题}
\begin{enumerate}
  \item 设底面长宽为 $x,y$，高为 $h$。费用约束为 \(\displaystyle 3xy+2xh+2yh=36\)。
  最大化 $V=xyh$。由对称性取 $x=y=a$，则
  \(\displaystyle 3a^2+4ah=36,\qquad V=a^2h=\frac{a(36-3a^2)}4\)。令 $V'=9-\dfrac94a^2=0$，得 $a=2$，进而 $h=3$。故底面为 $2\times2$，高为 $3$。
  \item 保守场条件为 $P_y=Q_x$，即 \(\displaystyle 2xy=yf'(x)\)。
  故 $f'(x)=2x$，又 $f(0)=0$，得 $f(x)=x^2$。势函数满足
  \(\displaystyle \varphi_x=xy^2,\qquad \varphi_y=x^2y\)。取 $\varphi=\dfrac12x^2y^2$。于是
  \(\displaystyle I=\varphi(1,1)-\varphi(0,0)=\frac12\)。
\end{enumerate}

\subsubsection*{五、证明题}
\begin{enumerate}
  \item 螺线切向量为 \(\displaystyle r'(t)=(-a\sin t,\ a\cos t,\ b)\)。
  与 $z$ 轴方向 $k=(0,0,1)$ 的夹角 $\theta$ 满足
  \(\displaystyle \cos\theta=\frac{|r'(t)\cdot k|}{|r'(t)|}=\frac{|b|}{\sqrt{a^2+b^2}}\)，为常数。
  \item 由不等式 $2\sqrt{uv}\le u+v$，取 $u=a_n$，$v=1/n^2$，得
  \(\displaystyle \frac{\sqrt{a_n}}{n}\le\frac12a_n+\frac1{2n^2}\)。
  因 $\sum a_n$ 与 $\sum1/n^2$ 均收敛，由比较判别法知
  \(\displaystyle \sum_{n=1}^{\infty}\frac{\sqrt{a_n}}{n}\) 收敛。
\end{enumerate}

\section{2012级工科数学分析下 A卷}

\subsection{资料来源}
\begin{itemize}
  \item 试题：2012级计算机工科数学分析下期末考试题 A 卷，DOC 正文已抽取；公式对象已按 Word 公式图片逐项恢复。
  \item 公式图片审计：\texttt{tmp/past-exam-equations/2012-A/manifest.csv}。
\end{itemize}

\subsection{试题}
诚信应考，考试作弊将带来严重后果！

华南理工大学期末考试

《工科数学分析》2012--2013 第二学期期末考试试卷（A）

注意事项：
\begin{enumerate}
  \item 考前请将密封线内填写清楚；
  \item 所有答案请直接答在试卷上（或答题纸上）；
  \item 考试形式：闭卷；
  \item 本试卷共 5 个大题，满分 100 分，考试时间 120 分钟。
\end{enumerate}

\subsubsection*{一、单项选择题（每小题 3 分，共 15 分）}
\begin{enumerate}
  \item 下述函数中有可能成为二阶微分方程
  \[
  f(x,y,y'')=0
  \]
  的通解的是（\quad）。

  A. $C_1x+C_2y=0$；\quad B. $y-C_1^2=e^x+C_2$；\quad
  C. $C_1y=C_2e^{C_3x}$；\quad D. $y=C_1\ln x+C_2x+C_3$。

  \item 设 $f(x,y)$ 是可微函数，且 \(\displaystyle f(0,0)=0,\ f_x(0,0)=a,\ f_y(0,0)=b\)，令 $\varphi(t)=f(t,f(t,t))$，则 $\varphi'(0)=$（\quad）。

  A. $a$；\quad B. $a+b(a+b)$；\quad C. $a+1$；\quad D. $\dfrac{a}{1-b}$。

  \item
  \[
  \int_0^1dy\int_{-y}^{\sqrt{2y-y^2}} f(x,y)\,dx
  \]
  交换积分顺序后，正确的结果是（\quad）。

  A. $\displaystyle \int_{-1}^0dx\int_{-x}^1f(x,y)\,dy+\int_0^1dx\int_{1-\sqrt{1-x^2}}^1f(x,y)\,dy$；

  B. $\displaystyle \int_0^1dx\int_{1-\sqrt{1-x^2}}^1f(x,y)\,dy$；

  C. $\displaystyle \int_{-1}^0dx\int_{-x}^1f(x,y)\,dy$；

  D. $\displaystyle \int_{-1}^1dx\int_{1-\sqrt{1-x^2}}^1f(x,y)\,dy$。

  \item 设 $C$ 是右半圆周 \(\displaystyle x^2+y^2=a^2,\ x\ge0,\ a>0\)，则曲线积分 \(\displaystyle \int_C(x+y)\,ds\)（\quad）。

  A. $0$；\quad B. $2a^2$；\quad C. $a^2$；\quad D. $4a$。

  \item 设 $f(x)=x^2\ (0\le x<1)$，而
  \[
  S(x)=\sum_{n=1}^{\infty}b_n\sin n\pi x\quad(-\infty<x<+\infty),
  \]
  其中 \(\displaystyle b_n=2\int_0^1f(x)\sin n\pi x\,dx\quad(n=1,2,\cdots)\)，则 $S\!\left(-\dfrac12\right)$（\quad）。

  A. $-\dfrac12$；\quad B. $-\dfrac14$；\quad C. $\dfrac14$；\quad D. $\dfrac12$。
\end{enumerate}

\subsubsection*{二、填空题（每小题 3 分，共 15 分）}
\begin{enumerate}
  \item 微分方程 \(\displaystyle yy''-y'^2=0\) 的通解为 \underline{\hspace{4cm}}。
  \item 设 \(\displaystyle u=x+(y-1)\arcsin\frac{x}{y}\)，则 $\left.\dfrac{\partial u}{\partial x}\right|_{(1,2)}=$ \underline{\hspace{3cm}}。
  \item 设区域 \(\displaystyle D:0\le y\le x,\ 0\le x\le\pi\)，则二重积分
  \[
  \iint_D \sqrt{1-\sin^2x}\,dx\,dy=
  \]
  \underline{\hspace{3cm}}。
  \item 已知曲线 $C$ 为 \(\displaystyle x^2+y^2=ax\)，则曲线积分
  \[
  \int_C\sqrt{x^2+y^2}\,ds=
  \]
  \underline{\hspace{3cm}}。
  \item 幂级数 \(\displaystyle \sum_{n=1}^{\infty}\frac{(x-2)^n}{n4^n}\) 的收敛域为 \underline{\hspace{4cm}}。
\end{enumerate}

\subsubsection*{三、计算题（每小题 10 分，共 50 分）}
\begin{enumerate}
  \item 计算三重积分
  \[
  I=\iiint_\Omega (x^2+y^2+z^2)\,dV,
  \]
  其中 $\Omega:x^2+y^2+z^2\le2z$。
  \item 设 $\Sigma$ 是锥面 $z=\sqrt{x^2+y^2}$ 被平面 $z=0$ 及 $z=1$ 所截部分的外侧，计算第二类曲面积分
  \[
  I=\iint_\Sigma xy\,dy\,dz+y\,dz\,dx+(z^2-2z)\,dx\,dy.
  \]
  \item 求幂级数 \(\displaystyle \sum_{n=0}^{\infty}(2n+1)x^n\) 的和函数，并据此求数项级数 \(\displaystyle \sum_{n=0}^{\infty}\frac{2n+1}{2^n}\) 的值。
  \item 求 $a,b$ 的值，使得包含圆周
  \[
  (x-1)^2+y^2=1
  \]
  在其内部的椭圆 \(\displaystyle \frac{x^2}{a^2}+\frac{y^2}{b^2}=1\quad(a>0,\ b>0,\ a\ne b)\) 有最小的面积。
  \item 求微分方程
  \[
  y''-y=2xe^x
  \]
  的通解。
\end{enumerate}

\subsubsection*{四、证明题（本题 10 分）}
证明数项级数 \(\displaystyle \sum_{n=1}^{\infty}\sin(\pi\sqrt{n^2+1})\) 条件收敛。

\subsubsection*{五、应用题（本题 10 分）}
设某山峰可由曲面 \(\displaystyle z=5-x^2-2y^2\) 表示。位于点
\(\displaystyle \left(\frac12,-\frac12,\frac{17}{4}\right)\) 处的登山者发现其供氧面具漏气，必须返回。他应该沿哪个方向才能最快到达山底？

\subsection{答案与解析}

\subsubsection*{一、单项选择题}
\begin{enumerate}
  \item 二阶微分方程通解通常含两个独立任意常数，只有 B 项含两个独立常数且能表示 $y$，故选 B。
  \item 链式法则给出 \(\displaystyle \varphi'(0)=f_x(0,0)+f_y(0,0)\bigl(f_x(0,0)+f_y(0,0)\bigr)=a+b(a+b)\)。故选 B。
  \item 区域由 $-y\le x$ 与 $x^2+(y-1)^2\le1$ 组成，故选 A。
  \item 右半圆可取 $x=a\cos t,\ y=a\sin t,\ -\dfrac\pi2\le t\le\dfrac\pi2$。于是
  \[
  \int_C(x+y)\,ds=a^2\int_{-\pi/2}^{\pi/2}(\cos t+\sin t)\,dt=2a^2.
  \]
  故选 B。
  \item 正弦级数对应奇延拓且周期为 $2$，故 \(\displaystyle S\!\left(-\frac12\right)=-f\!\left(\frac12\right)=-\frac14\)。故选 B。
\end{enumerate}

\subsubsection*{二、填空题}
\begin{enumerate}
  \item 设 $p=y'$，由 \(\displaystyle yy''-(y')^2=0\) 得 $\left(\dfrac{y'}{y}\right)'=0$，故 \(\displaystyle y=Ce^{kx}\)。
  \item
  \[
  u_x=(y-1)\frac1{\sqrt{1-(x/y)^2}}\cdot\frac1y
  =\frac{y-1}{\sqrt{y^2-x^2}}.
  \]
  在 $(1,2)$ 处为 $\dfrac1{\sqrt3}$。
  \item 因 $x\in[0,\pi]$，有 $\sqrt{1-\sin^2x}=|\cos x|$。所以
  \[
  \iint_D|\cos x|\,dx\,dy=\int_0^\pi x|\cos x|\,dx=\pi.
  \]
  \item 圆 $x^2+y^2=ax$ 上可用极坐标 $r=a\cos\theta$，结果为 \(\displaystyle 2a^2\)。
  \item 半径 $R=4$，端点 $x=-2,6$。在 $x=6$ 时为 $\sum 1/n$ 发散；在 $x=-2$ 时为 $\sum(-1)^n/n$ 收敛，故收敛域为
  \(\displaystyle [-2,6)\)。
\end{enumerate}

\subsubsection*{三、计算题}
\begin{enumerate}
  \item 区域是球心 $(0,0,1)$、半径 $1$ 的球。令 $z=w+1$，则 \(\displaystyle x^2+y^2+z^2=x^2+y^2+w^2+1+2w\)。
  对称性使 $2w$ 积分为 $0$，故
  \[
  I=\iiint_{r\le1}r^2\,dV+\iiint_{r\le1}1\,dV
  =\frac{4\pi}{5}+\frac{4\pi}{3}=\frac{32\pi}{15}.
  \]
  \item 外侧锥面 $z=r$ 对应远离 $z$ 轴且 $z$ 分量为负的法向 \(\displaystyle (z_x,z_y,-1)=\left(\frac{x}{r},\frac{y}{r},-1\right)\)。
  因而
  \[
  dy\,dz=\frac{x}{r}\,dx\,dy,\quad dz\,dx=\frac{y}{r}\,dx\,dy,\quad dx\,dy=-dx\,dy.
  \]
  在 $D:x^2+y^2\le1$ 上积分得
  \[
  I=\iint_D\left(\frac{x^2y}{r}+\frac{y^2}{r}-r^2+2r\right)\,dx\,dy
  =\frac{\pi}{3}-\frac{\pi}{2}+\frac{4\pi}{3}
  =\frac{7\pi}{6}.
  \]
  \item 当 $|x|<1$，
  \[
  \sum_{n=0}^{\infty}(2n+1)x^n
  =2\sum_{n=0}^{\infty}nx^n+\sum_{n=0}^{\infty}x^n
  =\frac{2x}{(1-x)^2}+\frac1{1-x}
  =\frac{1+x}{(1-x)^2}.
  \]
  令 $x=\dfrac12$，得
  \[
  \sum_{n=0}^{\infty}\frac{2n+1}{2^n}=6.
  \]
  \item 与 2014 A 卷同型，最小面积椭圆为 \(\displaystyle a=2,\ b=\sqrt2\)。
  \item 齐次解为 $C_1e^x+C_2e^{-x}$。右端为一次多项式乘 $e^x$，且 $1$ 是一重特征根，设 \(\displaystyle y_p=e^x(Ax^2+Bx)\)。
  代入得 $A=\dfrac12,\ B=-\dfrac12$，故
  \[
  y=C_1e^x+C_2e^{-x}+\frac12e^x(x^2-x).
  \]
\end{enumerate}

\subsubsection*{四、证明题}
  因 \(\displaystyle \sqrt{n^2+1}=n+\frac1{2n}+O\!\left(\frac1{n^3}\right)\)，故
\[
\sin(\pi\sqrt{n^2+1})
=(-1)^n\sin\left(\frac{\pi}{2n}+O\!\left(\frac1{n^3}\right)\right)
\sim (-1)^n\frac{\pi}{2n}.
\]
原级数由交错级数判别法收敛；而绝对值级数与 $\sum 1/n$ 同阶发散，故条件收敛。

\subsubsection*{五、应用题}
  山峰高度函数为 $z=5-x^2-2y^2$。最快下山方向为负梯度方向：
  \(\displaystyle -\nabla z=(2x,4y)\)。在 $\left(\dfrac12,-\dfrac12\right)$ 处，
  \(\displaystyle -\nabla z=(1,-2)\)。故应沿水平面内方向 $(1,-2)$ 前进，单位方向为 \(\displaystyle \frac{(1,-2)}{\sqrt5}\)。

\section{2012级工科数学分析下 B卷}

\subsection{资料来源}
\begin{itemize}
  \item 试题：2012级计算机工科数学分析下期末考试题 B 卷，DOC 正文已抽取；公式对象已按 Word 公式图片逐项恢复。
  \item 公式图片审计：\texttt{tmp/past-exam-equations/2012-B/manifest.csv}。
\end{itemize}

\subsection{试题}
诚信应考，考试作弊将带来严重后果！

华南理工大学期末考试

《工科数学分析》2012--2013 第二学期期末考试试卷（B）

注意事项：
\begin{enumerate}
  \item 考前请将密封线内填写清楚；
  \item 所有答案请直接答在试卷上（或答题纸上）；
  \item 考试形式：闭卷；
  \item 本试卷共 5 个大题，满分 100 分，考试时间 120 分钟。
\end{enumerate}

\subsubsection*{一、单项选择题（每小题 3 分，共 15 分）}
\begin{enumerate}
  \item 初值问题 \(\displaystyle y''+4y'+4y=0,\qquad y(0)=1,\quad y'(0)=3\)
  的解为（\quad）。

  A. $y=(C_1+C_2x)e^{-2x}$；\quad B. $y=(1+5x)e^{-2x}$；\quad
  C. $y=(1+3x)e^{-2x}$；\quad D. $y=e^{-2x}+5e^{2x}$。

  \item 对于二元函数 $z=f(x,y)$，下列说法正确的是（\quad）。

  A. 若函数 $z=f(x,y)$ 在 $(x_0,y_0)$ 连续，则函数在该点一定存在偏导数；

  B. 若函数 $z=f(x,y)$ 在 $(x_0,y_0)$ 存在偏导数，则函数在该点一定连续；

  C. 若函数 $z=f(x,y)$ 在 $(x_0,y_0)$ 的偏导数连续，则函数在该点一定可微；

  D. 若函数 $z=f(x,y)$ 在 $(x_0,y_0)$ 可微，则函数在该点的偏导数一定连续。

  \item 曲面 \(\displaystyle 3x^2+y^2+z^2=12\)
  上点 $M(-1,0,3)$ 处的切平面与平面 $z=0$ 的夹角是（\quad）。

  A. $\dfrac{\pi}{6}$；\quad B. $\dfrac{\pi}{4}$；\quad C. $\dfrac{\pi}{3}$；\quad D. $\dfrac{\pi}{2}$。

  \item 已知 $L$ 是球面 \(\displaystyle x^2+y^2+z^2=1\)
  与平面 $x+y+z=0$ 的交线，则第一类曲线积分 \(\displaystyle \int_Lx^2\,ds\)
  （\quad）。

  A. $0$；\quad B. $\dfrac{2\pi}{3}$；\quad C. $\dfrac{\pi}{3}$；\quad D. $\dfrac{4\pi}{3}$。

  \item 幂级数 \(\displaystyle \sum_{n=1}^{\infty}\frac{x^n}{n}\)
  的和函数是（\quad）。

  A. $s(x)=-\ln(1-x),\ (-1\le x<1)$；\quad
  B. $s(x)=-\ln(x-1),\ (-1<x<1)$；

  C. $s(x)=-\ln(1+x),\ (-1\le x<1)$；\quad
  D. $s(x)=\ln(1-x),\ (-1\le x<1)$。
\end{enumerate}

\subsubsection*{二、填空题（每小题 3 分，共 15 分）}
\begin{enumerate}
  \item 已知 \(\displaystyle f(x,y)=\sin x\ln(y+1)+\cos y\ln(1-x)\)，
  则 $\left.\dfrac{\partial f}{\partial x}\right|_{(0,0)}=$ \underline{\hspace{3cm}}。
  \item 微分方程 \(\displaystyle y^{(4)}-y=0\)
  的通解为 \underline{\hspace{4cm}}。
  \item 设 $D=[0,1;0,2]$，则二重积分
  \(\displaystyle \iint_D\sqrt{x+y}\,dx\,dy=\underline{\hspace{3cm}}\)。
  \item 设 $L$ 是任意一条光滑闭曲线，则
  \(\displaystyle \oint_L(2xy+\cos x)\,dx+(x^2+\sin y)\,dy=\underline{\hspace{3cm}}\)。
  \item 函数项级数 \(\displaystyle \sum_{n=1}^{\infty}\frac1n\left(\frac{x-1}{x}\right)^n\)
  的收敛域为 \underline{\hspace{4cm}}。
\end{enumerate}

\subsubsection*{三、计算题（每小题 10 分，共 50 分）}
\begin{enumerate}
  \item 求曲面 \(\displaystyle x^2+y^2+z^2=x\)
  的切平面，使它垂直于平面 $x-y-z=2$ 和 $x-y-\dfrac12z=2$。
  \item 求微分方程 \(\displaystyle y''+(4x+e^{2y})(y')^3=0\)
  的通解。
  \item 计算三重积分 \(\displaystyle I=\iiint_D(x^2+y^2)\,dx\,dy\,dz\)，
  其中 $D$ 是由圆锥面 $x^2+y^2=z^2$ 与平面 $z=1$ 所围成的区域。
  \item 设 $\Sigma$ 为曲面 \(\displaystyle x^2+y^2+z^2=2ax\quad(a>0)\)，
  计算曲面积分 \(\displaystyle I=\iint_\Sigma(x^2+y^2+z^2)\,dS\)。
  \item 将函数 \(\displaystyle f(x)=\frac1{x^2+4x+3}\)
  在 $x=1$ 处展成幂级数。
\end{enumerate}

\subsubsection*{四、证明题（本题 10 分）}
设
\[
x=x(y,z),\qquad y=y(x,z),\qquad z=z(x,y)
\]
都是由方程
\[
F(x,y,z)=0
\]
所确定的具有连续偏导数的函数，证明
\[
\frac{\partial x}{\partial y}\frac{\partial y}{\partial z}\frac{\partial z}{\partial x}=-1.
\]

\subsubsection*{五、应用题（本题 10 分）}
设有平面力场
\[
\mathbf F(x,y)=(2xy^3-y^2\cos x)\mathbf i+(1-2y\sin x+3x^2y^2)\mathbf j,
\]
求一质点沿曲线
\[
L:2x=\pi y^2
\]
从点 $O(0,0)$ 运动到点 $A\!\left(\dfrac\pi2,1\right)$ 时变力 $\mathbf F$ 所作的功。

\subsection{答案与解析}

\subsubsection*{一、单项选择题}
\begin{enumerate}
  \item 特征方程 $(r+2)^2=0$，故 $y=(C_1+C_2x)e^{-2x}$。由初值 $C_1=1,\ C_2=5$，选 B。
  \item 偏导连续可推出可微，选 C。
  \item 曲面法向量为 $\nabla F=(6x,2y,2z)$，在 $M$ 处为 $(-6,0,6)$。切平面法向量与 $z=0$ 法向量 $(0,0,1)$ 的夹角余弦为 $\dfrac1{\sqrt2}$，两平面夹角为 $\dfrac\pi4$，选 B。
  \item 交线为单位球上的大圆。由对称性，圆上 $x^2$ 平均值为 $1/3$，弧长为 $2\pi$，故积分为 $\dfrac{2\pi}{3}$，选 B。
  \item 由经典展开
  \(\displaystyle -\ln(1-x)=\sum_{n=1}^{\infty}\frac{x^n}{n}\)
  且在 $x=-1$ 收敛、$x=1$ 发散，选 A。
\end{enumerate}

\subsubsection*{二、填空题}
\begin{enumerate}
  \item \(\displaystyle f_x=\cos x\ln(y+1)-\frac{\cos y}{1-x}\)，
  故 $f_x(0,0)=-1$。
  \item 特征根为 $1,-1,i,-i$，故通解为
  \(\displaystyle y=C_1e^x+C_2e^{-x}+C_3\cos x+C_4\sin x\)。
  \item 先对 $x$ 积分：
  \[
  \iint_D\sqrt{x+y}\,dx\,dy
  =\frac23\int_0^2\bigl((y+1)^{3/2}-y^{3/2}\bigr)\,dy
  =\frac4{15}(3^{5/2}-2).
  \]
  \item 由 Green 公式，旋度为 $2x-2x=0$，故积分为 $0$。
  \item 令 $t=\dfrac{x-1}{x}=1-\dfrac1x$。级数 $\sum t^n/n$ 的收敛域为 $[-1,1)$，即
  \(\displaystyle -1\le1-\frac1x<1\)。
  解得
  \(\displaystyle x\in\left[\frac12,+\infty\right)\)。
\end{enumerate}

\subsubsection*{三、计算题}
\begin{enumerate}
  \item 两平面法向量分别为 $\mathbf n_1=(1,-1,-1)$ 与 $\mathbf n_2=(1,-1,-1/2)$。所求切平面同时垂直于它们，其法向量应平行于
  \(\displaystyle \mathbf n_1\times\mathbf n_2=(-1,-1,0)\)。
  曲面法向量为 $(2x-1,2y,2z)$，令其平行于 $(1,1,0)$，得 $z=0,\ 2x-1=2y$。与曲面联立，得两点
  \(\displaystyle \left(0,-\frac12,0\right),\qquad \left(\frac45,\frac3{10},0\right)\)。对应切平面为
  \(\displaystyle x+y+\frac12=0,\qquad x+y-\frac{11}{10}=0\)。
  \item 令 $p=y'$，把 $x$ 看作 $y$ 的函数，利用
  \(\displaystyle y''=-\frac{x''}{(x')^3},\qquad y'=\frac1{x'}\)。
  原方程化为 \(\displaystyle -x''+4x+e^{2y}=0\)，即
  \(\displaystyle x''-4x=e^{2y}\)。以 $y$ 为自变量求解得
  \(\displaystyle x=C_1e^{2y}+C_2e^{-2y}+\frac14 y e^{2y}\)。该式为原方程通解的隐式形式。
  \item 用柱坐标，$0\le z\le1,\ 0\le r\le z$。于是
  \(\displaystyle I=\int_0^{2\pi}\int_0^1\int_0^z r^2\cdot r\,dr\,dz\,d\theta
  =\frac{\pi}{6}\)。
  \item 曲面为以 $(a,0,0)$ 为球心、半径为 $a$ 的球面，且曲面上 $x^2+y^2+z^2=2ax$。由对称性
  \(\displaystyle \iint_\Sigma x\,dS=a\cdot4\pi a^2\)，
  故
  \(\displaystyle I=2a\iint_\Sigma x\,dS=8\pi a^4\)。
  \item
  \[
  \frac1{x^2+4x+3}=\frac1{(x+1)(x+3)}
  =\frac12\left(\frac1{x+1}-\frac1{x+3}\right).
  \]
  令 $t=x-1$，则
  \[
  f(x)=\frac12\left(\frac1{t+2}-\frac1{t+4}\right)
  =\sum_{n=0}^{\infty}(-1)^n\left(\frac1{2^{n+2}}-\frac1{2^{2n+3}}\right)t^n,
  \]
  其中 $|t|<2$，即 $|x-1|<2$。
\end{enumerate}

\subsubsection*{四、证明题}
由 $F(x,y,z)=0$ 隐函数求导：
\[
\frac{\partial x}{\partial y}=-\frac{F_y}{F_x},\qquad
\frac{\partial y}{\partial z}=-\frac{F_z}{F_y},\qquad
\frac{\partial z}{\partial x}=-\frac{F_x}{F_z}.
\]
三式相乘立即得
\[
\frac{\partial x}{\partial y}\frac{\partial y}{\partial z}\frac{\partial z}{\partial x}=-1.
\]

\subsubsection*{五、应用题}
记 \(\displaystyle P=2xy^3-y^2\cos x,\qquad Q=1-2y\sin x+3x^2y^2\)。有
\(\displaystyle P_y=6xy^2-2y\cos x,\qquad Q_x=-2y\cos x+6xy^2\)，故力场保守。求势函数：\(\displaystyle \Phi_x=P\)。对 $x$ 积分得
\(\displaystyle \Phi=x^2y^3-y^2\sin x+h(y)\)。再由 $\Phi_y=Q$ 得 $h'(y)=1$，故
\(\displaystyle \Phi=x^2y^3-y^2\sin x+y\)。所作功为
\(\displaystyle W=\Phi\!\left(\frac\pi2,1\right)-\Phi(0,0)
=\frac{\pi^2}{4}-1+1=\frac{\pi^2}{4}\)。

\section{2013级工科数学分析下 A卷}

\subsection{资料来源}
\begin{itemize}
  \item 试题：2013 级软件工科数学分析下 A 卷，已导出页面图校对。
  \item 答案：同卷附解答文件。
\end{itemize}

\subsection{试题}

\paragraph*{试卷信息}
诚信应考，考试作弊将带来严重后果！

华南理工大学本科生期末考试《工科数学分析》2013--2014学年第二学期期末考试试卷（A）卷。

注意事项：开考前请将密封线内各项信息填写清楚；所有答案请直接答在试卷上（或答题纸上）；考试形式：闭卷；本试卷共5个大题，满分100分，考试时间120分钟。

\paragraph*{一、单项选择题（每小题3分，共15分）}
\begin{enumerate}
  \item 若 $y_1,y_2$ 是非齐次线性方程 $y''+ay'+by=f(x)$ 的两个特解，则下面结论中正确的是（\quad）。
  \begin{enumerate}
    \item[A.] $y_1+y_2$ 是非齐次线性方程的解。
    \item[B.] $y_1-y_2$ 是非齐次线性方程的解。
    \item[C.] $y_1+y_2$ 是 $y''+ay'+by=0$ 的解。
    \item[D.] $y_1-y_2$ 是 $y''+ay'+by=0$ 的解。
  \end{enumerate}

  \item 设 $z=f(x,y)$ 可微，且 $f(x,x^2)=\frac12,\ \left.f_x(x,y)\right|_{y=x^2}=\frac1x$，则 $\left.f_y(x,y)\right|_{y=x^2}$（\quad）。
  \begin{enumerate}
    \item[A.] $-\dfrac{1}{x^2}$
    \item[B.] $\dfrac{1}{x^2}$
    \item[C.] $\dfrac{1}{2x^2}$
    \item[D.] $-\dfrac{1}{2x^2}$
  \end{enumerate}

  \item 设 $C:x^2+y^2=a^2$ 取逆时针方向，则 $\displaystyle\oint_C \frac{(x+y)\,dx-(x-y)\,dy}{x^2+y^2}$（\quad）。
  \begin{enumerate}
    \item[A.] $-2\pi$
    \item[B.] $2\pi$
    \item[C.] $0$
    \item[D.] $-\pi$
  \end{enumerate}

  \item 幂级数 $\displaystyle\sum_{n=1}^{\infty}\frac{(-1)^{n-1}}{n}(x-1)^n$ 的收敛域是（\quad）。
  \begin{enumerate}
    \item[A.] $(0,2)$
    \item[B.] $(0,2]$
    \item[C.] $[0,2)$
    \item[D.] $[0,2]$
  \end{enumerate}

  \item 设 $f(x)=x^2\ (0\le x<1)$，而
  \[
    S(x)=\sum_{n=1}^{\infty}b_n\sin n\pi x\quad(-\infty<x<+\infty),
    \qquad
    b_n=2\int_0^1 f(x)\sin n\pi x\,dx.
  \]
  则 $S\!\left(-\dfrac12\right)$（\quad）。
  \begin{enumerate}
    \item[A.] $-\dfrac12$
    \item[B.] $\dfrac14$
    \item[C.] $-\dfrac14$
    \item[D.] $\dfrac12$
  \end{enumerate}
\end{enumerate}

\paragraph*{二、填空题（每小题3分，共15分）}
\begin{enumerate}
  \item 微分方程 $y^{(4)}-y=0$ 的通解为 \underline{\hspace{5cm}}。
  \item 函数 $z=2x^2+y^2$ 在点 $(1,1)$ 处沿该点梯度方向的方向导数为 \underline{\hspace{5cm}}。
  \item 交换积分的顺序，则 $\displaystyle\int_1^e dx\int_0^{\ln x} f(x,y)\,dy=$ \underline{\hspace{5cm}}。
  \item 设 $\Sigma:x^2+y^2+z^2=a^2$，则 $\displaystyle\iint_{\Sigma}(x^2+y^2+z^2)\,dS=$ \underline{\hspace{5cm}}。
  \item 设幂级数 $\displaystyle\sum_{n=1}^{\infty}a_nx^n$ 的收敛半径为 $3$，则幂级数
  \[
    \sum_{n=1}^{\infty}n a_n(x-1)^{n+1}
  \]
  的收敛区间为 \underline{\hspace{5cm}}。
\end{enumerate}

\paragraph*{三、计算题（每小题10分，共50分）}
\begin{enumerate}
  \item 设函数 $f(u)$ 具有二阶连续导数，而 $z=f(e^x\sin y)$ 满足方程 $\dfrac{\partial^2 z}{\partial x^2}+\dfrac{\partial^2 z}{\partial y^2}=e^{2x}z$，求 $f(u)$。
  \item 计算三重积分 $I=\displaystyle\iiint_{\Omega}(x+z)\,dV$，其中 $\Omega$ 是由曲面 $z=x^2+y^2$ 与曲面 $z=1-x^2-y^2$ 所围成的区域。
  \item 设 $f(x)$ 在 $(-\infty,+\infty)$ 内有连续导数，$L$ 是从点 $A(3,\frac23)$ 到点 $B(1,2)$ 的直线段，计算曲线积分
  \[
    I=\int_L \frac{1+y^2f(xy)}{y}\,dx
    +\frac{x}{y^2}\bigl[y^2f(xy)-1\bigr]\,dy.
  \]
  \item 设 $\Sigma$ 是锥面 $z=\sqrt{x^2+y^2}$ 被平面 $z=0$ 及 $z=1$ 所截部分的下侧，计算第二类曲面积分 $I=\displaystyle\iint_{\Sigma}x\,dy\,dz+y\,dz\,dx+(z^2-2z)\,dx\,dy$。
  \item 求幂级数 $\displaystyle\sum_{n=1}^{\infty}\frac{2n+1}{n!}x^{2n}$ 的和函数。
\end{enumerate}

\paragraph*{四、证明题（本题10分）}
\begin{enumerate}
  \setcounter{enumi}{5}
  \item 证明级数 $\displaystyle\sum_{n=1}^{\infty}\frac{(-1)^n(1-e^{-nx})}{n^2+x^2}$ 在 $[0,+\infty)$ 上一致收敛。
\end{enumerate}

\paragraph*{五、应用题（本题10分）}
\begin{enumerate}
  \setcounter{enumi}{6}
  \item 在第一象限内作椭球面 $\dfrac{x^2}{a^2}+\dfrac{y^2}{b^2}+\dfrac{z^2}{c^2}=1$ 的切平面，使切平面与三个坐标面所围成的四面体的体积最小，求切点的坐标。
\end{enumerate}

\subsection{答案与解析}

\paragraph*{一、单项选择题}
\begin{enumerate}
  \item 两个非齐次方程特解之差满足对应齐次方程，即 $y_1-y_2$ 是 $y''+ay'+by=0$ 的解。选 D。
  \item 对恒等式 $f(x,x^2)=\frac12$ 求导，得 $f_x(x,x^2)+2xf_y(x,x^2)=0$，代入 $f_x(x,x^2)=\frac1x$，得 $f_y(x,x^2)=-\dfrac1{2x^2}$。选 D。
  \item 在圆周上分解
  \(\displaystyle \frac{(x+y)\,dx-(x-y)\,dy}{x^2+y^2}
  =d\ln r-d\theta\)，闭曲线第一项积分为 $0$，逆时针一周第二项为 $-2\pi$。选 A。
  \item 令 $t=x-1$，该幂级数为 \(\displaystyle \ln(1+t)\)，其收敛区间为 $-1<t\le1$，即 $0<x\le2$。选 B。
  \item 该正弦级数对应 $[0,1]$ 上 $f(x)=x^2$ 的奇延拓，故 $S(-1/2)=-f(1/2)=-\dfrac14$。选 C。
\end{enumerate}

\paragraph*{二、填空题}
\begin{enumerate}
  \item 微分方程 $y^{(4)}-y=0$ 的通解为 $y=C_1e^x+C_2e^{-x}+C_3\cos x+C_4\sin x$。
  \item 函数 $z=2x^2+y^2$ 在点 $(1,1)$ 处沿该点梯度方向的方向导数为 $2\sqrt5$。
  \item 交换积分的顺序，则 $\displaystyle\int_1^e dx\int_0^{\ln x} f(x,y)\,dy$ $=\displaystyle\int_0^1dy\int_{e^y}^{e}f(x,y)\,dx$。
  \item 设 $\Sigma:x^2+y^2+z^2=a^2$，则 $\displaystyle\iint_{\Sigma}(x^2+y^2+z^2)\,dS$ $=4\pi a^4$。
  \item 设幂级数 $\displaystyle\sum_{n=1}^{\infty}a_nx^n$ 的收敛半径为 $3$，则幂级数 $\displaystyle\sum_{n=1}^{\infty}n a_n(x-1)^{n+1}$ 的收敛区间为 $(-2,4)$。
\end{enumerate}

\paragraph*{三、计算题}
\begin{enumerate}
  \item 解：由 $u=e^x\sin y$；又由 $z=f(u)$ 满足方程
  \[
    \frac{\partial^2 z}{\partial x^2}
    +\frac{\partial^2 z}{\partial y^2}=e^{2x}z,
  \]
  得到 $f''(u)=f(u)$，解得 $f(u)=C_1e^u+C_2e^{-u}$。
  \item 解：根据对称性，$\displaystyle\iiint_{\Omega}x\,dV=0$（单独2分）。利用柱面坐标可得
  \[
    I=2\pi\int_0^{1/\sqrt2}dr\int_{r^2}^{1-r^2}zr\,dz
    =\frac{\pi}{8};
  \]
  或利用球面坐标可得 $I=\dfrac{\pi}{8}$。
  \item 解：
  \[
    d\left(F(xy)+\frac{x}{y}\right)
    =\frac{1+y^2f(xy)}{y}\,dx
    +\frac{x}{y^2}\bigl[y^2f(xy)-1\bigr]\,dy.
  \]
  法一：由原积分为全微分积分（6分），则
  \(\displaystyle I=\left[F(xy)+\frac{x}{y}\right]_{(3,2/3)}^{(1,2)}=-4\)。法二：记 $F'(t)=f(t)$，则 $I=-4$。
  \item 解：法一：取参数
  \[
    x=z\cos\theta,\qquad y=z\sin\theta,\qquad
    0\le z\le1,\quad 0\le\theta\le2\pi.
  \]
  下侧对应面积向量 $(z\cos\theta,z\sin\theta,-z)\,d\theta dz$，故
  \[
    I=\int_0^{2\pi}\int_0^1(3z^2-z^3)\,dz\,d\theta
    =\frac{3\pi}{2}.
  \]
  法二：利用参数曲面积分同样可得 $I=\dfrac{3\pi}{2}$。法三：补上下底面后利用 Gauss 公式可得同一结果。
  \item 解：
  \[
    \sum_{n=1}^{\infty}\frac{2n+1}{n!}x^{2n}
    =2\sum_{n=1}^{\infty}\frac{n(x^2)^n}{n!}
    +\sum_{n=1}^{\infty}\frac{(x^2)^n}{n!}
    =2x^2e^{x^2}+e^{x^2}-1.
  \]
  因此 $S(x)=(2x^2+1)e^{x^2}-1$。
\end{enumerate}

\paragraph*{四、证明题}
证明：因为
\(\displaystyle \left|\frac{(-1)^n(1-e^{-nx})}{n^2+x^2}\right|\le \frac{1}{n^2}\)，
且 $\displaystyle\sum_{n=1}^{\infty}\frac1{n^2}$ 收敛，所以级数
\(\displaystyle \sum_{n=1}^{\infty}\frac{(-1)^n(1-e^{-nx})}{n^2+x^2}\)
在 $[0,+\infty)$ 上一致收敛（10分）。

\paragraph*{五、应用题}
解：设 $P(x_0,y_0,z_0)$ 为椭球面上在第一象限的一点。过此点的切平面方程为
\[
  \frac{x_0x}{a^2}+\frac{y_0y}{b^2}+\frac{z_0z}{c^2}=1,
\]
即
\[
  \frac{x}{a^2/x_0}+\frac{y}{b^2/y_0}+\frac{z}{c^2/z_0}=1.
\]
此切平面与坐标面围成四面体的体积为
\(\displaystyle V=\frac{a^2b^2c^2}{6x_0y_0z_0}\)。
（3分）
要求 $V$ 在条件
\(\displaystyle \frac{x_0^2}{a^2}+\frac{y_0^2}{b^2}+\frac{z_0^2}{c^2}=1\)
下的最小值，只需求 $x_0y_0z_0$ 在同一条件下的最大值。由拉格朗日乘数法，只需求函数
\[
  F=x_0y_0z_0+\lambda\left(\frac{x_0^2}{a^2}
  +\frac{y_0^2}{b^2}+\frac{z_0^2}{c^2}-1\right)
\]
的驻点（6分），即 $F_{x_0}=F_{y_0}=F_{z_0}=0$（8分）。由驻点方程两两相除得
\(\displaystyle \frac{x_0}{a}=\frac{y_0}{b}=\frac{z_0}{c}\)。
由此得
\(\displaystyle \frac{x_0^2}{a^2}=\frac{y_0^2}{b^2}=\frac{z_0^2}{c^2}=\frac13\)，
所以
\(\displaystyle x_0=\frac{a}{\sqrt3},\qquad
  y_0=\frac{b}{\sqrt3},\qquad
  z_0=\frac{c}{\sqrt3}\)。
（9分）
由驻点唯一，又根据实际意义最小体积存在，所以切点坐标为
\(\displaystyle \left(\frac{a}{\sqrt3},\frac{b}{\sqrt3},\frac{c}{\sqrt3}\right)\)。
（10分）

\section{2013级工科数学分析下 B卷}

\subsection{资料来源}
\begin{itemize}
  \item 试题：2013 级软件工科数学分析下 B 卷，已导出页面图校正公式。
  \item 未发现配套答案或解析文本；下列答案与解析为据题面补写。
\end{itemize}

\subsection{试题}

\paragraph*{试卷信息}
诚信应考，考试作弊将带来严重后果！

华南理工大学本科生期末考试《工科数学分析》2013--2014学年第二学期期末考试试卷（B）卷。

注意事项：开考前请将密封线内各项信息填写清楚；所有答案请直接答在试卷上（或答题纸上）；考试形式：闭卷；本试卷共5个大题，满分100分，考试时间120分钟。

\paragraph*{一、单项选择题（每小题3分，共15分）}
\begin{enumerate}
  \item 方程 \((y-\ln x)\,\mathrm{d}x+x\,\mathrm{d}y=0\) 是（\quad）。
  \begin{enumerate}
    \item[A.] 可分离变量微分方程。
    \item[B.] 一阶线性非齐次方程。
    \item[C.] 一阶线性齐次方程。
    \item[D.] 非线性方程。
  \end{enumerate}

  \item 设 \(f(x,y)\) 是可微函数，且 \(f(0,0)=0\)，\(f_x(0,0)=a\)，\(f_y(0,0)=b\)，令
  \(\varphi(t)=f(t,f(t,t))\)，则 \(\varphi'(0)=\)（\quad）。
  \begin{enumerate}
    \item[A.] \(a\)
    \item[B.] \(a+b(a+b)\)
    \item[C.] \(a+1\)
    \item[D.] \(\dfrac{a}{1-b}\)
  \end{enumerate}

  \item 设 \(C:x^2+y^2=a^2\)，取逆时针方向，则
  \(\displaystyle \oint_C xy^2\,\mathrm{d}y-x^2y\,\mathrm{d}x=\)（\quad）。
  \begin{enumerate}
    \item[A.] \(\dfrac{\pi a^4}{2}\)
    \item[B.] \(-\dfrac{\pi a^4}{2}\)
    \item[C.] \(\pi a^4\)
    \item[D.] \(-\pi a^4\)
  \end{enumerate}

  \item 幂级数 \(\displaystyle \sum_{n=1}^{\infty}\frac{x^n}{n2^n}\) 的收敛域是（\quad）。
  \begin{enumerate}
    \item[A.] \((-2,2)\)
    \item[B.] \([-2,2]\)
    \item[C.] \((-2,2]\)
    \item[D.] \([-2,2)\)
  \end{enumerate}

  \item 设 \(f(x)\) 是周期为 \(2\) 的周期函数，它在区间 \((-1,1]\) 上的表达式为
  \[
    f(x)=
    \begin{cases}
      2, & -1<x\le 0,\\
      x^3, & 0<x\le 1,
    \end{cases}
  \]
  则 \(f(x)\) 的 Fourier 级数在 \(x=1\) 处收敛于（\quad）。
  \begin{enumerate}
    \item[A.] \(1\)
    \item[B.] \(0\)
    \item[C.] \(\dfrac32\)
    \item[D.] \(-\dfrac32\)
  \end{enumerate}
\end{enumerate}

\paragraph*{二、填空题（每小题3分，共15分）}
\begin{enumerate}
  \item 微分方程 \(y''-2y'+y=0\) 的通解为 \underline{\hspace{5cm}}。
  \item 函数 \(\displaystyle z=\arctan\frac{y}{x}+\ln\sqrt{x^2+y^2}\) 在点 \((1,2)\) 处的梯度为 \underline{\hspace{5cm}}。
  \item 交换积分的顺序，则 \(\displaystyle \int_0^2\mathrm{d}x\int_x^2 \mathrm{e}^{-y^2}\,\mathrm{d}y=\) \underline{\hspace{5cm}}。
  \item 设 \(\Sigma:x^2+y^2+z^2=a^2\)，则
  \(\displaystyle \iint_\Sigma (x^2+y^2+z^2)\,\mathrm{d}S=\) \underline{\hspace{5cm}}。
  \item 设幂级数
  \[
    \sum_{n=1}^{\infty}a_n(x-1)^n
  \]
  在 \(x=-1\) 处收敛，则此级数在 \(x=2\) 的敛散性是 \underline{\hspace{5cm}}。
\end{enumerate}

\paragraph*{三、计算题（每小题10分，共50分）}
\begin{enumerate}
  \item 设 \(f(\xi,\eta)\) 具有连续的二阶偏导数，且满足 Laplace 方程
  \(\displaystyle \frac{\partial^2 f}{\partial \xi^2}+\frac{\partial^2 f}{\partial \eta^2}=0\)，令
  \(z=f(x^2-y^2,2xy)\)，求
  \(\displaystyle \frac{\partial^2z}{\partial x^2}+\frac{\partial^2z}{\partial y^2}\)。

  \item 计算三重积分
  \[
    I=\iiint_\Omega z\mathrm{e}^{-(x^2+y^2+z^2)}\,\mathrm{d}V,
    \qquad
    \Omega:\begin{cases}
      1\le x^2+y^2+z^2\le 4,\\
      x\ge 0,\ y\ge 0,\ z\ge 0.
    \end{cases}
  \]

  \item 已知 \(\varphi(\pi)=1\)，试确定 \(\varphi(x)\)，使曲线积分
  \[
    I=\int_A^B\left[\sin x-\varphi(x)\right]\frac{y}{x}\,\mathrm{d}x+\varphi(x)\,\mathrm{d}y
  \]
  与路径无关，并求当 \(A,B\) 两点分别为 \((1,0),(\pi,\pi)\) 时积分 \(I\) 的值。

  \item 计算曲面积分
  \[
    I=\iint_\Sigma (8y+1)x\,\mathrm{d}y\mathrm{d}z+2(1-y^2)\,\mathrm{d}z\mathrm{d}x-4yz\,\mathrm{d}x\mathrm{d}y,
  \]
  其中 \(\Sigma:y=1+x^2+z^2\)，介于 \(1\le y\le 3\) 之间，且其法向量与 \(y\) 轴正向的夹角恒大于 \(\dfrac{\pi}{2}\)。

  \item 求级数 \(\displaystyle \sum_{n=1}^{\infty}\frac{(-1)^n n}{(2n+1)!}\) 的和。
\end{enumerate}

\paragraph*{四、证明题（本题10分）}
\begin{enumerate}
  \setcounter{enumi}{5}
  \item 证明级数 \(\displaystyle \sum_{n=1}^{\infty}\frac{\sin(nx)}{n^2}\) 在 \((-\infty,+\infty)\) 上一致收敛。
\end{enumerate}

\paragraph*{五、应用题（本题10分）}
\begin{enumerate}
  \setcounter{enumi}{6}
  \item 求内接于半径为 \(a\) 的球且有最大体积的长方体。
\end{enumerate}

\subsection{答案与解析}
\paragraph*{一、单项选择题}
\begin{enumerate}
  \item 选 B。原方程可化为 \(\displaystyle y'+\frac1x y=\frac{\ln x}{x}\)，因而是一阶线性非齐次方程。

  \item 选 B。由链式法则，\(\displaystyle \varphi'(t)=f_x(t,f(t,t))+f_y(t,f(t,t))\frac{\mathrm{d}}{\mathrm{d}t}f(t,t)\)，且 \(\left.\dfrac{\mathrm{d}}{\mathrm{d}t}f(t,t)\right|_{t=0}=f_x(0,0)+f_y(0,0)=a+b\)，所以
  \(\varphi'(0)=a+b(a+b)\)。

  \item 选 A。令 \(P=-x^2y\)，\(Q=xy^2\)。由 Green 公式，
  \[
    \oint_C P\,\mathrm{d}x+Q\,\mathrm{d}y
    =\iint_D(Q_x-P_y)\,\mathrm{d}A
    =\iint_D(x^2+y^2)\,\mathrm{d}A
    =\frac{\pi a^4}{2}.
  \]

  \item 选 D。由比值判别法得收敛半径 \(R=2\)。当 \(x=2\) 时为调和级数 \(\sum 1/n\)，发散；当 \(x=-2\) 时为交错调和级数，收敛。因此收敛域为 \([-2,2)\)。

  \item 选 C。周期为 \(2\)，在 \(x=1\) 处按 Fourier 收敛定理取左右极限平均值：\(\displaystyle \frac{f(1-0)+f(1+0)}2=\frac{1+2}{2}=\frac32\)。
\end{enumerate}

\paragraph*{二、填空题}
\begin{enumerate}
  \item 特征方程 \((r-1)^2=0\)，故通解为 \(\displaystyle y=(C_1+C_2x)\mathrm{e}^x\)。

  \item 有
  \[
    z_x=-\frac{y}{x^2+y^2}+\frac{x}{x^2+y^2},\qquad
    z_y=\frac{x}{x^2+y^2}+\frac{y}{x^2+y^2}.
  \]
  所以
  \[
    \nabla z(1,2)=\left(-\frac15,\frac35\right).
  \]

  \item 积分区域为 \(0\le x\le y\le 2\)，故
  \[
    \int_0^2\mathrm{d}x\int_x^2 \mathrm{e}^{-y^2}\,\mathrm{d}y
    =\int_0^2\mathrm{d}y\int_0^y \mathrm{e}^{-y^2}\,\mathrm{d}x
    =\int_0^2y\mathrm{e}^{-y^2}\,\mathrm{d}y
    =\frac{1-\mathrm{e}^{-4}}2.
  \]

  \item 在球面上 \(x^2+y^2+z^2=a^2\)，故 \(\displaystyle \iint_\Sigma (x^2+y^2+z^2)\,\mathrm{d}S=a^2\cdot 4\pi a^2=4\pi a^4\)。

  \item 该幂级数以 \(x=1\) 为中心，在 \(x=-1\) 处收敛，说明收敛半径至少为 \(2\)。点 \(x=2\) 到中心的距离为 \(1<2\)，故级数在 \(x=2\) 处绝对收敛。
\end{enumerate}

\paragraph*{三、计算题}
\begin{enumerate}
  \item 记 \(\xi=x^2-y^2,\eta=2xy\)。由链式法则直接计算可得
  \[
    \frac{\partial^2 z}{\partial x^2}+\frac{\partial^2 z}{\partial y^2}
    =4(x^2+y^2)\left(\frac{\partial^2 f}{\partial \xi^2}+\frac{\partial^2 f}{\partial \eta^2}\right)=0.
  \]

  \item 用球坐标
  \(x=r\sin\varphi\cos\theta,\ y=r\sin\varphi\sin\theta,\ z=r\cos\varphi\)。区域为
  \(1\le r\le 2,\ 0\le\varphi\le\pi/2,\ 0\le\theta\le\pi/2\)，因此
  \[
    I=\int_0^{\pi/2}\!\!\int_0^{\pi/2}\!\!\int_1^2
    r\cos\varphi\,\mathrm{e}^{-r^2}r^2\sin\varphi\,\mathrm{d}r\,\mathrm{d}\varphi\,\mathrm{d}\theta
    =\frac{\pi}{8}\left(\frac2{\mathrm{e}}-\frac5{\mathrm{e}^4}\right).
  \]

  \item 设 \(\displaystyle P=\left[\sin x-\varphi(x)\right]\frac{y}{x},\ Q=\varphi(x)\)。路径无关要求 \(P_y=Q_x\)，即 \(\displaystyle \frac{\sin x-\varphi(x)}{x}=\varphi'(x),\ (x\varphi(x))'=\sin x\)。
  积分得 \(x\varphi(x)=-\cos x+C\)。由 \(\varphi(\pi)=1\) 得 \(C=\pi-1\)，所以
  \[
    \varphi(x)=\frac{\pi-1-\cos x}{x}.
  \]
  势函数可取 \(F(x,y)=y\varphi(x)\)，故
  \(\displaystyle I=F(\pi,\pi)-F(1,0)=\pi\)。

  \item 取参数 \((x,z)\)，令 \(y=1+x^2+z^2\)。题设法向量的 \(y\) 分量为负，故取 \(\displaystyle \mathbf r_x\times\mathbf r_z=(2x,-1,2z)\)。
  投影区域为 \(D:x^2+z^2\le 2\)。代入第二类曲面积分，
  \[
    I=\iint_D\left[(8y+1)x\cdot2x-2(1-y^2)-4yz\cdot2z\right]\mathrm{d}x\mathrm{d}z,
    \qquad y=1+x^2+z^2.
  \]
  化简后利用圆盘对称性：
  \[
    I=10\iint_Dx^2\,\mathrm{d}A+4\iint_Dr^2\,\mathrm{d}A+6\iint_Dr^4\,\mathrm{d}A
    =10\pi+8\pi+16\pi=34\pi.
  \]

  \item 由
  \[
    \sin t=\sum_{n=0}^{\infty}\frac{(-1)^n t^{2n+1}}{(2n+1)!}
  \]
  得
  \[
    \sum_{n=1}^{\infty}\frac{(-1)^n n}{(2n+1)!}
    =\frac12\sum_{n=1}^{\infty}\frac{(-1)^n(2n+1)}{(2n+1)!}
     -\frac12\sum_{n=1}^{\infty}\frac{(-1)^n}{(2n+1)!}
    =\frac{\cos1-\sin1}{2}.
  \]
\end{enumerate}

\paragraph*{四、证明题}
\begin{enumerate}
  \setcounter{enumi}{5}
  \item 对任意 \(x\in(-\infty,+\infty)\)，都有 \(\displaystyle \left|\frac{\sin(nx)}{n^2}\right|\le \frac1{n^2}\)。
  又数项级数 \(\sum_{n=1}^{\infty}1/n^2\) 收敛，由 Weierstrass 判别法，级数
  \(\sum_{n=1}^{\infty}\sin(nx)/n^2\) 在 \((-\infty,+\infty)\) 上一致收敛。
\end{enumerate}

\paragraph*{五、应用题}
\begin{enumerate}
  \setcounter{enumi}{6}
  \item 设长方体以球心为中心，三个半边长为 \(x,y,z>0\)，则约束为 \(\displaystyle x^2+y^2+z^2=a^2\)，体积为 \(V=8xyz\)。由对称性或 Lagrange 乘子法，最大值在 \(x=y=z=a/\sqrt3\) 处取得。因此长方体三条棱长均为 \(\displaystyle \frac{2a}{\sqrt3}\)，最大体积为 \(\displaystyle V_{\max}=8\left(\frac{a}{\sqrt3}\right)^3=\frac{8a^3}{3\sqrt3}\)。
\end{enumerate}

\section{2014级工科数学分析下 A卷}

\subsection{资料来源}
\begin{itemize}
  \item 试题：2014级软件工科数学分析下 A 卷，DOC 正文已抽取；公式对象已按 Word 公式图片逐项恢复。
  \item 公式图片审计：\texttt{tmp/past-exam-equations/2014-A/manifest.csv}。
\end{itemize}

\subsection{试题}
诚信应考，考试作弊将带来严重后果！

华南理工大学本科生期末考试

《工科数学分析》2014--2015 学年第二学期期末考试试卷（A）卷

注意事项：
\begin{enumerate}
  \item 开考前请将密封线内各项信息填写清楚；
  \item 所有答案请直接答在试卷上（或答题纸上）；
  \item 考试形式：闭卷；
  \item 本试卷共 5 个大题，满分 100 分，考试时间 120 分钟。
\end{enumerate}

\subsubsection*{一、填空题（每小题 4 分，共 20 分）}
\begin{enumerate}
  \item 微分方程 \(\displaystyle y''-4y'=e^{4x}\) 的特解形式为 \underline{\hspace{4cm}}。
  \item 设 \(\displaystyle xyz+\sqrt{x^2+y^2+z^2}=\sqrt2\)，则 $\left.dz\right|_{(1,0,-1)}=$ \underline{\hspace{3cm}}。
  \item 交换积分的顺序，则
  \[
  \int_0^1dy\int_y^{\sqrt{2y-y^2}}f(x,y)\,dx=
  \]
  \underline{\hspace{3cm}}。
  \item 设 $f(x)=x^2\ (0\le x<1)$，而
  \[
  S(x)=\sum_{n=1}^{\infty}b_n\sin n\pi x\quad(-\infty<x<+\infty),
  \]
  其中
  \[
  b_n=2\int_0^1 f(x)\sin n\pi x\,dx\quad(n=1,2,\cdots),
  \]
  则 $S\!\left(-\dfrac12\right)$ \underline{\hspace{3cm}}。
  \item 设 \(\displaystyle L:\frac{x^2}{a^2}+\frac{y^2}{b^2}=1\) 取逆时针方向，则
  \[
  \int_L(x-y)\,dx+(y-x)\,dy=
  \]
  \underline{\hspace{3cm}}。
\end{enumerate}

\subsubsection*{二、计算题（每小题 10 分，共 40 分）}
\begin{enumerate}
  \item 设
  \[
  u=yf\!\left(\frac{x}{y}\right)+xg\!\left(\frac{y}{x}\right),
  \]
  其中 $f,g$ 具有二阶连续偏导数，求
  \[
  x\frac{\partial^2u}{\partial x^2}+y\frac{\partial^2u}{\partial x\partial y}.
  \]
  \item 计算三重积分
  \[
  I=\iiint_\Omega (x^2+y^2)\,dx\,dy\,dz,
  \]
  其中 $\Omega$ 是由圆锥面 $z=\sqrt{x^2+y^2}$ 与平面 $z=1$ 所围成的区域。
  \item 计算曲线积分
  \[
  I=\int_L \cos(x+y^2)\,dx+\left[2y\cos(x+y^2)-\frac1{\sqrt{1+y^4}}\right]\,dy,
  \]
  其中
  \[
  L:\begin{cases}
  x=a(t-\sin t),\\
  y=a(1-\cos t),
  \end{cases}
  \]
  上从 $O(0,0)$ 到 $A(2\pi a,0)$ 的有向曲线弧段。
  \item 设 $\Sigma$ 是抛物面 $z=x^2+y^2$ 被平面 $z=0$ 及 $z=1$ 所截部分的下侧，计算曲面积分
  \[
  I=\iint_\Sigma xy\,dy\,dz+2y\,dz\,dx+(z^2-3z)\,dx\,dy.
  \]
\end{enumerate}

\subsubsection*{三、解答题（每小题 10 分，共 20 分）}
\begin{enumerate}
  \item 求幂级数
  \[
  \sum_{n=1}^{\infty}\frac{(-1)^n n}{(2n+1)!}x^{2n+1}
  \]
  的和函数。
  \item 设
  \[
  f(0)=0,\qquad f'(x)=1+\int_0^x(6\sin^2t-f(t))\,dt,
  \]
  其中 $f(x)$ 二次可微，求 $f(x)$。
\end{enumerate}

\subsubsection*{四、证明题（本题 10 分）}
设正项级数 $\sum_{n=1}^{\infty}a_n$ 收敛，证明
\[
f(x)=\sum_{n=1}^{\infty}a_nx^n
\]
在 $(-1,1)$ 上连续。

\subsubsection*{五、应用题（本题 10 分）}
求 $a,b$ 的值，使得包含圆周 $(x-1)^2+y^2=1$ 在其内部的椭圆
\[
\frac{x^2}{a^2}+\frac{y^2}{b^2}=1\quad(a>0,\ b>0,\ a\ne b)
\]
有最小的面积。

\subsection{答案与解析}

\subsubsection*{一、填空题}
\begin{enumerate}
  \item 特征方程为 $r^2-4r=0$，$r=4$ 是一重特征根，右端与齐次解共振，故特解形式为
  \[
  y^*=Axe^{4x}.
  \]
  \item 令 $F=xyz+\sqrt{x^2+y^2+z^2}-\sqrt2$。在 $(1,0,-1)$ 处，
  \[
  F_x=\frac1{\sqrt2},\qquad F_y=-1,\qquad F_z=-\frac1{\sqrt2}.
  \]
  由 $F_x\,dx+F_y\,dy+F_z\,dz=0$ 得
  \[
  dz=dx-\sqrt2\,dy.
  \]
  \item 区域满足 $0\le y\le1,\ y\le x\le\sqrt{2y-y^2}$，即
  \[
  x^2+(y-1)^2\le1,\qquad y\le x.
  \]
  交换顺序得
  \[
  \int_0^1dx\int_{1-\sqrt{1-x^2}}^x f(x,y)\,dy.
  \]
  \item 正弦级数对应 $[0,1]$ 上的奇延拓且周期为 $2$，故
  \[
  S\!\left(-\frac12\right)=-f\!\left(\frac12\right)=-\frac14.
  \]
  \item 由 Green 公式，
  \[
  \int_L(x-y)\,dx+(y-x)\,dy
  =\iint_D\bigl((-1)-(-1)\bigr)\,dA=0.
  \]
\end{enumerate}

\subsubsection*{二、计算题}
\begin{enumerate}
  \item 设 $s=x/y,\ t=y/x$。直接求导得
  \[
  u_x=f'(s)+g(t)-tg'(t),
  \]
  \[
  u_{xx}=\frac1y f''(s)+\frac{t^3}{y}g''(t),\qquad
  u_{xy}=-\frac{x}{y^2}f''(s)-\frac{t^2}{y}g''(t).
  \]
  因此
  \[
  xu_{xx}+yu_{xy}=0.
  \]
  \item 用柱坐标，区域为 $0\le r\le z,\ 0\le z\le1,\ 0\le\theta\le2\pi$。于是
  \[
  I=\int_0^{2\pi}\int_0^1\int_0^z r^2\cdot r\,dr\,dz\,d\theta
  =\frac{\pi}{10}.
  \]
  \item 注意
  \[
  d\sin(x+y^2)=\cos(x+y^2)\,dx+2y\cos(x+y^2)\,dy.
  \]
  故
  \[
  I=\left.\sin(x+y^2)\right|_{O}^{A}-\int_L\frac{dy}{\sqrt{1+y^4}}.
  \]
  端点处正弦项均为 $0$。沿摆线一拱，$y$ 从 $0$ 到 $2a$ 再回到 $0$，第二项为同一单变量函数关于 $y$ 的全微分积分，故也为 $0$。因此
  \[
  I=0.
  \]
  \item 下侧取向对应法向量
  \[
  (z_x,z_y,-1)=(2x,2y,-1).
  \]
  对 $z=x^2+y^2$，有
  \[
  dy\,dz=2x\,dx\,dy,\qquad dz\,dx=2y\,dx\,dy,\qquad dx\,dy=-dx\,dy.
  \]
  投影区域为 $D:x^2+y^2\le1$，故
  \[
  I=\iint_D\bigl(2x^2y+4y^2-(z^2-3z)\bigr)\,dx\,dy.
  \]
  奇函数项积分为 $0$，令 $z=r^2$，得
  \[
  I=\iint_D(4y^2-r^4+3r^2)\,dA
  =\pi+\left(-\frac{\pi}{3}+\frac{3\pi}{2}\right)
  =\frac{13\pi}{6}.
  \]
\end{enumerate}

\subsubsection*{三、解答题}
\begin{enumerate}
  \item 令
  \[
  S(x)=\sum_{n=1}^{\infty}\frac{(-1)^n n}{(2n+1)!}x^{2n+1}.
  \]
  由
  \[
  \sin x=\sum_{n=0}^{\infty}\frac{(-1)^n x^{2n+1}}{(2n+1)!}
  \]
  及
  \[
  x\cos x-\sin x=\sum_{n=1}^{\infty}\frac{(-1)^n 2n}{(2n+1)!}x^{2n+1}
  \]
  得
  \[
  S(x)=\frac{x\cos x-\sin x}{2}.
  \]
  \item 对已知等式求导，得
  \[
  f''(x)=6\sin^2x-f(x).
  \]
  即
  \[
  f''+f=3-3\cos2x.
  \]
  又由原式 $f'(0)=1$，并给出 $f(0)=0$。通解为
  \[
  f=C_1\cos x+C_2\sin x+3+\cos2x.
  \]
  代入初值 $f(0)=0,\ f'(0)=1$，得 $C_1=-4,\ C_2=1$，所以
  \[
  f(x)=-4\cos x+\sin x+3+\cos2x.
  \]
\end{enumerate}

\subsubsection*{四、证明题}
因 $\sum a_n$ 为正项收敛级数，对任意 $x\in(-1,1)$ 有
\[
|a_nx^n|\le a_n.
\]
由 Weierstrass 判别法，$\sum a_nx^n$ 在 $[-r,r]$ 上一致收敛，其中任取 $0<r<1$。每项 $a_nx^n$ 连续，所以和函数在任意 $[-r,r]\subset(-1,1)$ 上连续，从而在 $(-1,1)$ 上连续。

\subsubsection*{五、应用题}
椭圆包含圆周 $(x-1)^2+y^2=1$ 等价于对圆盘内所有点有
\[
\frac{x^2}{a^2}+\frac{y^2}{b^2}\le1.
\]
令圆周参数为 $x=1+\cos t,\ y=\sin t$。要求
\[
\max_t\left(\frac{(1+\cos t)^2}{a^2}+\frac{\sin^2t}{b^2}\right)\le1.
\]
最小面积时有相切的临界情形。用支持函数或 Lagrange 乘子求解，可得最优椭圆为圆盘的最小面积中心椭圆：
\[
a=2,\qquad b=\sqrt2.
\]
此时椭圆面积为 $2\sqrt2\pi$。

\section{2014级工科数学分析下 B卷}

\subsection{资料来源}
\begin{itemize}
  \item 试题：2014级软件工科数学分析下 B 卷，DOC 正文已抽取；公式对象已按 Word 公式图片逐项恢复。
  \item 公式图片审计：\texttt{tmp/past-exam-equations/2014-B/manifest.csv}。
\end{itemize}

\subsection{试题}
诚信应考，考试作弊将带来严重后果！

华南理工大学本科生期末考试

《工科数学分析》2014--2015 学年第二学期期末考试试卷（B）卷

注意事项：
\begin{enumerate}
  \item 开考前请将密封线内各项信息填写清楚；
  \item 所有答案请直接答在试卷上（或答题纸上）；
  \item 考试形式：闭卷；
  \item 本试卷共 5 个大题，满分 100 分，考试时间 120 分钟。
\end{enumerate}

\subsubsection*{一、填空题（每小题 4 分，共 20 分）}
\begin{enumerate}
  \item 微分方程
  \[
  y''-y=2xe^x
  \]
  的特解形式为 \underline{\hspace{4cm}}。
  \item 设由方程
  \[
  f(x-y,y-z,z-x)=0
  \]
  确定了 $z=z(x,y)$，其中 $f$ 有连续的偏导数，则 $dz=$ \underline{\hspace{3cm}}。
  \item 交换积分的顺序，则
  \[
  \int_0^1dy\int_{y^2}^{\sqrt y}f(x,y)\,dx=
  \]
  \underline{\hspace{3cm}}。
  \item 设 $f(x)$ 是周期为 $2$ 的周期函数，它在区间 $(-1,1]$ 上的表达式为
  \[
  f(x)=
  \begin{cases}
  3,&-1<x\le0,\\
  x^3,&0<x\le1,
  \end{cases}
  \]
  则 $f(x)$ 的 Fourier 级数在 $x=1$ 处收敛于 \underline{\hspace{3cm}}。
  \item 设
  \[
  L:|x|+|y|=1
  \]
  取逆时针方向，则
  \[
  \int_L xy^2\,dx+x^2y\,dy=
  \]
  \underline{\hspace{3cm}}。
\end{enumerate}

\subsubsection*{二、计算题（每小题 10 分，共 40 分）}
\begin{enumerate}
  \item 设
  \[
  z=f\!\left(2x,\frac{y}{x}\right),
  \]
  其中 $f$ 具有二阶连续的偏导数，求 $\dfrac{\partial^2z}{\partial x\partial y}$。
  \item 计算三重积分
  \[
  I=\iiint_\Omega z\,dx\,dy\,dz,
  \]
  其中 $\Omega$ 是由球面 $x^2+y^2+z^2=4$ 与抛物面 $x^2+y^2=3z$ 所围成的区域。
  \item 计算曲线积分
  \[
  I=\int_L\frac{-y\,dx+x\,dy}{x^2+4y^2},
  \]
  其中 $L$ 是从 $A(-\pi,0)$ 沿曲线 $y=\cos\dfrac{x}{2}$ 到 $B(\pi,0)$ 的有向曲线弧段。
  \item 设 $\Sigma$ 是锥面 $z=\sqrt{x^2+y^2}$ 被平面 $z=0$ 及 $z=1$ 所截部分的下侧，计算曲面积分
  \[
  I=\iint_\Sigma 2x\,dy\,dz+3y\,dz\,dx+(z^2-5z)\,dx\,dy.
  \]
\end{enumerate}

\subsubsection*{三、解答题（每小题 10 分，共 20 分）}
\begin{enumerate}
  \item 求幂级数 \(\displaystyle \sum_{n=1}^{\infty}\frac{2n-1}{n!}x^{2n}\) 的和函数。
  \item 设 \(\displaystyle f(x)=\sin x-\int_0^x(x-t)f(t)\,dt\)，其中 $f(x)$ 为连续函数，求 $f(x)$。
\end{enumerate}

\subsubsection*{四、证明题（本题 10 分）}
证明：级数 \(\displaystyle \sum_{n=1}^{\infty}\frac{\sin(nx)}{n^2}\) 在 $(-\infty,+\infty)$ 上一致收敛。

\subsubsection*{五、应用题（本题 10 分）}
求椭球面 \(\displaystyle x^2+y^2+\frac{z^2}{4}=1\) 在第一卦限的一点，使该点处的切平面在三个坐标轴上的截距平方和最小。

\subsection{答案与解析}

\subsubsection*{一、填空题}
\begin{enumerate}
  \item 特征方程 $r^2-1=0$，$r=1$ 是一重特征根。右端为一次多项式乘 $e^x$，故特解形式为 \(\displaystyle y^*=xe^x(Ax+B)\)。
  \item 记 $f_1,f_2,f_3$ 为 $f$ 对三个自变量的偏导数。由
  \[
  f_1(dx-dy)+f_2(dy-dz)+f_3(dz-dx)=0
  \]
  得
  \[
  dz=\frac{(f_3-f_1)\,dx+(f_1-f_2)\,dy}{f_3-f_2}.
  \]
  \item 区域为 $0\le y\le1,\ y^2\le x\le\sqrt y$。等价于 $0\le x\le1,\ x^2\le y\le\sqrt x$，故
  \[
  \int_0^1dx\int_{x^2}^{\sqrt x}f(x,y)\,dy.
  \]
  \item 周期为 $2$，在 $x=1$ 处左右极限分别为 $x^3\to1$ 与周期延拓后的 $3$，故 Fourier 级数收敛于 \(\displaystyle \frac{1+3}{2}=2\)。
  \item 由 Green 公式，\(\displaystyle \int_Lxy^2\,dx+x^2y\,dy=\iint_D(2xy-2xy)\,dA=0\)。
\end{enumerate}

\subsubsection*{二、计算题}
\begin{enumerate}
  \item 令 $\xi=2x,\ \eta=y/x$。先有 \(\displaystyle z_y=\frac1x f_\eta\)。
  再对 $x$ 求偏导，得
  \[
  z_{xy}=\frac2x f_{\xi\eta}-\frac{y}{x^3}f_{\eta\eta}-\frac1{x^2}f_\eta,
  \]
  其中各偏导在 $\left(2x,\dfrac{y}{x}\right)$ 处取值。
  \item 用柱坐标。两曲面交于 $r^2+z^2=4,\ r^2=3z$，得 $z=1,\ r=\sqrt3$。所围区域为
  \[
  0\le r\le\sqrt3,\qquad \frac{r^2}{3}\le z\le\sqrt{4-r^2},\qquad 0\le\theta\le2\pi .
  \]
  因此
  \[
  I=\int_0^{2\pi}\int_0^{\sqrt3}\int_{r^2/3}^{\sqrt{4-r^2}}z r\,dz\,dr\,d\theta
  =\frac{13\pi}{4}.
  \]
  \item 令 $u=x,\ v=2y$，则 \(\displaystyle \frac{-y\,dx+x\,dy}{x^2+4y^2}=\frac12\frac{-v\,du+u\,dv}{u^2+v^2}\)。
  沿题设方向从 $A(-\pi,0)$ 到 $B(\pi,0)$ 时，向量 $(u,v)$ 的极角从 $\pi$ 变到 $0$，故
  \[
  I=\frac12(0-\pi)=-\frac{\pi}{2}.
  \]
  \item 下侧锥面 $z=r$ 对应法向取 \(\displaystyle (z_x,z_y,-1)=\left(\frac{x}{r},\frac{y}{r},-1\right)\)。
  投影区域 $D:x^2+y^2\le1$，且 $z=r$。于是
  \[
  I=\iint_D\left(2x\frac{x}{r}+3y\frac{y}{r}-(r^2-5r)\right)\,dx\,dy.
  \]
  用极坐标积分：
  \[
  I=\int_0^{2\pi}\int_0^1\left(2r\cos^2\theta+3r\sin^2\theta-r^2+5r\right)r\,dr\,d\theta
  =\frac{9\pi}{2}.
  \]
\end{enumerate}

\subsubsection*{三、解答题}
\begin{enumerate}
  \item 令 $t=x^2$。则
  \[
  \sum_{n=1}^{\infty}\frac{2n-1}{n!}x^{2n}
  =2\sum_{n=1}^{\infty}\frac{n t^n}{n!}-\sum_{n=1}^{\infty}\frac{t^n}{n!}
  =2te^t-(e^t-1).
  \]
  因此和函数为 \(\displaystyle (2x^2-1)e^{x^2}+1\)。
  \item 对积分方程求两次导数。先有 \(\displaystyle f'(x)=\cos x-\int_0^x f(t)\,dt\)，
  再得 \(\displaystyle f''(x)=-\sin x-f(x)\)，即 \(\displaystyle f''+f=-\sin x\)。由原式得 $f(0)=0,\ f'(0)=1$。通解为 \(\displaystyle f=C_1\cos x+C_2\sin x+\frac12 x\cos x\)。代入初值，得 \(\displaystyle f(x)=\sin x+\frac12 x\cos x\)。
\end{enumerate}

\subsubsection*{四、证明题}
对任意 $x\in(-\infty,+\infty)$，\(\displaystyle \left|\frac{\sin(nx)}{n^2}\right|\le\frac1{n^2}\)。
而 $\sum_{n=1}^{\infty}\dfrac1{n^2}$ 收敛，由 Weierstrass 判别法，原级数在 $(-\infty,+\infty)$ 上一致收敛。

\subsubsection*{五、应用题}
椭球面可写为 $F=x^2+y^2+\dfrac{z^2}{4}-1=0$。在点 $(x_0,y_0,z_0)$ 处切平面为
\(\displaystyle x_0x+y_0y+\frac{z_0z}{4}=1\)。三个截距为 \(\displaystyle \frac1{x_0},\ \frac1{y_0},\ \frac4{z_0}\)。
需在约束 $x_0^2+y_0^2+\dfrac{z_0^2}{4}=1$ 下最小化
\(\displaystyle \frac1{x_0^2}+\frac1{y_0^2}+\frac{16}{z_0^2}\)。
令 $u=x_0^2,\ v=y_0^2,\ w=z_0^2/4$，则 $u+v+w=1$，目标为
\(\displaystyle \frac1u+\frac1v+\frac4w\)。
Lagrange 乘子法给出 $u=v,\ w=2u$，故 $4u=1$。第一卦限内所求点为
\(\displaystyle \left(\frac12,\frac12,\sqrt2\right)\)。

\section{2015级工科数学分析下 A卷}

\subsection{资料来源}
\begin{itemize}
  \item 试题：2015级软件工科数学分析下 A 卷，DOC 正文已抽取；公式对象已按 Word 公式图片逐项恢复。
  \item 公式图片审计：\texttt{tmp/past-exam-equations/2015-A/manifest.csv}。
\end{itemize}

\subsection{试题}
诚信应考，考试作弊将带来严重后果！

华南理工大学本科生期末考试

《工科数学分析》2015--2016 学年第二学期期末考试试卷（A）卷

注意事项：
\begin{enumerate}
  \item 开考前请将密封线内各项信息填写清楚；
  \item 所有答案请直接答在试卷上（或答题纸上）；
  \item 考试形式：闭卷；
  \item 本试卷共 5 个大题，满分 100 分，考试时间 120 分钟。
\end{enumerate}

\subsubsection*{一、填空题（每小题 3 分，共 15 分）}
\begin{enumerate}
  \item 设 $y_1=3,\ y_2=3+x^2,\ y_3=3+x^2+e^x$ 都是方程
  \[
  (x^2-2x)y''-(x^2-2)y'+(2x-2)y=6x-6
  \]
  的解，则该方程的通解为 \underline{\hspace{4cm}}。
  \item 设由方程 $F(x-y,y-z,z-x)=0$ 确定了 $z=z(x,y)$，其中 $F$ 可微，则 $dz=$ \underline{\hspace{3cm}}。
  \item 函数
  \[
    f(x,y,z)=\ln(x^2+y^2+z^2)
  \]
  在点 $P(2,0,1)$ 处沿梯度方向的方向导数是 \underline{\hspace{4cm}}。
  \item 设 $L:\dfrac{x^2}{a^2}+\dfrac{y^2}{b^2}=1$ 取顺时针方向，则
  \[
  \int_L(x-y)\,dx+(x+y)\,dy=
  \]
  \underline{\hspace{3cm}}。
  \item 设 $f(x)=\sin(\pi x^2)\ (0\le x<1)$，而
  \[
  S(x)=\sum_{n=1}^{\infty}b_n\sin(n\pi x)\quad(-\infty<x<+\infty),
  \]
  其中 $b_n=2\int_0^1 f(x)\sin(n\pi x)\,dx\ (n=1,2,\cdots)$，则 $S\!\left(-\dfrac12\right)$ \underline{\hspace{3cm}}。
\end{enumerate}

\subsubsection*{二、计算题（每小题 8 分，共 40 分）}
\begin{enumerate}
  \item 设函数 $z=f(x^2-y^2,2xy)$，其中 $f(\xi,\eta)$ 具有连续的二阶偏导数，且满足
  \[
  \frac{\partial^2 f}{\partial \xi^2}+\frac{\partial^2 f}{\partial \eta^2}=0,
  \]
  求 $\dfrac{\partial^2 z}{\partial x^2}+\dfrac{\partial^2 z}{\partial y^2}$。
  \item 计算曲线积分
  \[
  I=\int_L\sqrt{x^2+y^2}\,ds,
  \]
  其中 $L$ 是圆周 $x^2+y^2=ax$。
  \item 设 $f(x)$ 在 $(-\infty,+\infty)$ 内有连续导数，$L$ 是从点 $A\!\left(3,\dfrac23\right)$ 到点 $B(1,2)$ 的直线段，计算曲线积分
  \[
  I=\int_L\frac{1+y^2f(x)}{y}\,dx+\frac{x}{y^2}\bigl[y^2f(x)-1\bigr]\,dy.
  \]
  \item 设 $\Sigma$ 为曲面 $x^2+y^2+z^2=2ax\ (a>0)$，计算曲面积分
  \[
  I=\iint_\Sigma (x^2+y^2+z^2)\,dS.
  \]
  \item 设 $\Sigma$ 是上半球面 $z=\sqrt{a^2-x^2-y^2}\ (a>0)$，取上侧，计算曲面积分
  \[
  I=\iint_\Sigma x^3\,dy\,dz+y^3\,dz\,dx+z^3\,dx\,dy.
  \]
\end{enumerate}

\subsubsection*{三、解答题（每小题 9 分，共 18 分）}
\begin{enumerate}
  \item 求幂级数
  \[
  \sum_{n=1}^{\infty}\frac{2n+1}{n!}x^{2n}
  \]
  的和函数。
  \item 求微分方程
  \[
  xy'-2y=x^3\sin x
  \]
  的通解。
\end{enumerate}

\subsubsection*{四、证明题（每小题 9 分，共 18 分）}
\begin{enumerate}
  \item 设 $f(x)\in C[0,1]$，求证：
  \[
  \int_0^1 dx\int_x^1 f(x)f(y)\,dy=\frac12\left[\int_0^1 f(x)\,dx\right]^2.
  \]
  \item 证明：函数项级数
  \[
  \sum_{n=1}^{\infty}\frac{\cos(nx)}{n^2}
  \]
  在 $(-\infty,+\infty)$ 上一致收敛。
\end{enumerate}

\subsubsection*{五、应用题（本题 9 分）}
求内接于半径为 $a$ 的球且有最大体积的长方体。

\subsection{答案与解析}

\subsubsection*{一、填空题}
\begin{enumerate}
  \item 因 $y_2-y_1=x^2,\ y_3-y_2=e^x$ 是对应齐次方程的两个线性无关解，且 $y_1=3$ 是一个特解，故通解为
  \[
  y=3+C_1x^2+C_2e^x.
  \]
  \item 记 $F_1,F_2,F_3$ 分别为 $F$ 对三个自变量的偏导数。由
  \[
  F_1(dx-dy)+F_2(dy-dz)+F_3(dz-dx)=0
  \]
  得
  \[
  dz=\frac{(F_3-F_1)\,dx+(F_1-F_2)\,dy}{F_3-F_2}.
  \]
  \item $\nabla f=\dfrac{2(x,y,z)}{x^2+y^2+z^2}$，故
  \[
  |\nabla f(2,0,1)|=\frac{2\sqrt5}{5}.
  \]
  \item 由 Green 公式，逆时针时积分为
  \[
  \iint_D\left(\frac{\partial}{\partial x}(x+y)-\frac{\partial}{\partial y}(x-y)\right)\,dA
  =2\pi ab.
  \]
  本题 $L$ 取顺时针方向，故结果为 $-2\pi ab$。
  \item 该正弦级数对应 $[0,1]$ 上的奇延拓且周期为 $2$，所以
  \[
  S\!\left(-\frac12\right)=-f\!\left(\frac12\right)=-\sin\frac{\pi}{4}=-\frac{\sqrt2}{2}.
  \]
\end{enumerate}

\subsubsection*{二、计算题}
\begin{enumerate}
  \item 令 $\xi=x^2-y^2,\ \eta=2xy$。该变换满足
  \[
  \xi_x^2+\eta_x^2=\xi_y^2+\eta_y^2=4(x^2+y^2),\qquad
  \xi_x\xi_y+\eta_x\eta_y=0,
  \]
  且 $\xi_{xx}+\xi_{yy}=\eta_{xx}+\eta_{yy}=0$。因此
  \[
  z_{xx}+z_{yy}=4(x^2+y^2)(f_{\xi\xi}+f_{\eta\eta})=0.
  \]
  \item 圆 $x^2+y^2=ax$ 可参数化为
  \[
  x=\frac a2(1+\cos t),\qquad y=\frac a2\sin t,\qquad 0\le t\le 2\pi.
  \]
  此时 $\sqrt{x^2+y^2}=a|\cos(t/2)|,\ ds=\dfrac a2\,dt$，故
  \[
  I=\frac{a^2}{2}\int_0^{2\pi}\left|\cos\frac t2\right|dt=2a^2.
  \]
  \item 直线段 $AB$ 上
  \[
  y=\frac{8-2x}{3},\qquad x:3\to 1.
  \]
  被积表达式可写为
  \[
  d\!\left(\frac{x}{y}\right)+f(x)\,d(xy).
  \]
  于是
  \[
  I=\left.\frac{x}{y}\right|_A^B+\int_3^1 f(x)\frac{8-4x}{3}\,dx
  =-4+\frac13\int_3^1(8-4x)f(x)\,dx.
  \]
  \item 曲面是以 $(a,0,0)$ 为球心、半径为 $a$ 的球面。在该球面上
  \[
  x^2+y^2+z^2=2ax.
  \]
  由球面对称性 $\iint_\Sigma x\,dS=a\cdot 4\pi a^2$，故
  \[
  I=2a\iint_\Sigma x\,dS=8\pi a^4.
  \]
  \item 把上半球面与底面圆盘合成上半球体的闭曲面。底面上 $z=0$，第三项通量为 $0$。由 Gauss 公式，
  \[
  I=\iiint_\Omega 3(x^2+y^2+z^2)\,dV
  =3\int_0^a r^2\cdot 2\pi r^2\,dr
  =\frac{6\pi a^5}{5}.
  \]
\end{enumerate}

\subsubsection*{三、解答题}
\begin{enumerate}
  \item 令 $t=x^2$，则
  \[
  \sum_{n=1}^{\infty}\frac{2n+1}{n!}x^{2n}
  =2\sum_{n=1}^{\infty}\frac{n t^n}{n!}+\sum_{n=1}^{\infty}\frac{t^n}{n!}
  =2te^t+(e^t-1).
  \]
  因此和函数为
  \[
  (2x^2+1)e^{x^2}-1.
  \]
  \item 当 $x\ne0$ 时，方程化为
  \[
  y'-\frac2x y=x^2\sin x.
  \]
  取积分因子 $x^{-2}$，得
  \[
  \left(\frac{y}{x^2}\right)'=\sin x.
  \]
  因而
  \[
  \frac{y}{x^2}=-\cos x+C,\qquad y=Cx^2-x^2\cos x.
  \]
\end{enumerate}

\subsubsection*{四、证明题}
\begin{enumerate}
  \item 记 $I=\int_0^1 dx\int_x^1 f(x)f(y)\,dy$。交换 $x,y$ 的角色可得
  \[
  I=\int_0^1 dy\int_0^y f(x)f(y)\,dx.
  \]
  两式相加覆盖单位正方形除对角线外的两个三角区域，而对角线面积为零，所以
  \[
  2I=\int_0^1\int_0^1 f(x)f(y)\,dx\,dy
  =\left[\int_0^1 f(x)\,dx\right]^2.
  \]
  结论得证。
  \item 对任意 $x\in(-\infty,+\infty)$，有
  \[
  \left|\frac{\cos(nx)}{n^2}\right|\le \frac1{n^2}.
  \]
  又 $\sum_{n=1}^{\infty}\dfrac1{n^2}$ 收敛，由 Weierstrass 判别法，原函数项级数在 $(-\infty,+\infty)$ 上一致收敛。
\end{enumerate}

\subsubsection*{五、应用题}
设长方体中心与球心重合，三个半棱长为 $x,y,z>0$。约束为
\[
x^2+y^2+z^2=a^2,\qquad V=8xyz.
\]
由对称性或 Lagrange 乘子法，最大值在 $x=y=z=\dfrac a{\sqrt3}$ 处取得。因此最大体积长方体为正方体，棱长为
\[
\frac{2a}{\sqrt3},
\]
最大体积为
\[
V_{\max}=\left(\frac{2a}{\sqrt3}\right)^3=\frac{8a^3}{3\sqrt3}.
\]

\section{2015级工科数学分析下 B卷}

\subsection{资料来源}
\begin{itemize}
  \item 试题：2015级软件工科数学分析下 B 卷，DOC 正文已抽取；公式对象已按 Word 公式图片逐项恢复。
  \item 公式图片审计：\texttt{tmp/past-exam-equations/2015-B/manifest.csv}。
\end{itemize}

\subsection{试题}
诚信应考，考试作弊将带来严重后果！

华南理工大学本科生期末考试

《工科数学分析》2015--2016 学年第二学期期末考试试卷（B）卷

注意事项：
\begin{enumerate}
  \item 开考前请将密封线内各项信息填写清楚；
  \item 所有答案请直接答在试卷上（或答题纸上）；
  \item 考试形式：闭卷；
  \item 本试卷共 5 个大题，满分 100 分，考试时间 120 分钟。
\end{enumerate}

\subsubsection*{一、填空题（每小题 3 分，共 15 分）}
\begin{enumerate}
  \item 微分方程
  \[
  y''-4y=xe^{2x}
  \]
  的特解形式为 \underline{\hspace{4cm}}。
  \item 设 $z=z(x,y)$ 满足方程
  \[
  x-az=\varphi(y-bz),
  \]
  其中 $\varphi$ 可微，$a,b$ 为常数，则
  \[
  a\frac{\partial z}{\partial x}+b\frac{\partial z}{\partial y}=
  \]
  \underline{\hspace{3cm}}。
  \item 函数
  \[
    f(x,y,z)=\cos(xyz)
  \]
  在点 $P\!\left(\dfrac13,1,\pi\right)$ 处使方向导数取得最大的方向是 \underline{\hspace{4cm}}。
  \item 设 $L:x^2+y^2=a^2$ 取逆时针方向，则
  \[
  \int_L 2xy\,dx+(x^2+8y^2)\,dy=
  \]
  \underline{\hspace{3cm}}。
  \item 设幂级数
  \[
  \sum_{n=0}^{\infty}a_nx^n
  \]
  在 $x=-3$ 条件收敛，则该幂级数的收敛半径为 $R=$ \underline{\hspace{2cm}}。
\end{enumerate}

\subsubsection*{二、计算题（每小题 8 分，共 40 分）}
\begin{enumerate}
  \item 设函数
  \[
  z=f\!\left(xy,\frac{x}{y}\right),
  \]
  其中 $f(\xi,\eta)$ 具有连续的二阶偏导数，求 $\dfrac{\partial^2z}{\partial x\partial y}$。
  \item 计算曲线积分
  \[
  I=\int_\Gamma x^2\,ds,
  \]
  其中 $\Gamma$ 是球面 $x^2+y^2+z^2=a^2$ 与平面 $x+y+z=0$ 的交线。
  \item 设曲线积分
  \[
  \int_L\bigl(f(x)-e^x\bigr)\sin y\,dx-f'(x)\cos y\,dy
  \]
  与路径无关，其中 $f'(x)$ 有一阶的连续导数，且 $f'(0)=0$。

  （1）求 $f'(x)$；（2）计算曲线积分
  \[
  I=\int_{(0,0)}^{(1,1)}\bigl(f(x)-e^x\bigr)\sin y\,dx-f'(x)\cos y\,dy.
  \]

  \item 设 $\Sigma$ 为平面 $y+z=5$ 被柱面 $x^2+y^2=25$ 所截得的部分，计算曲面积分
  \[
  I=\iint_\Sigma (x+y+z)\,dS.
  \]
  \item 设 $\Sigma$ 是上半球面 $z=\sqrt{a^2-x^2-y^2}\ (a>0)$，取上侧，计算曲面积分
  \[
  I=\iint_\Sigma (x^3-xy^2)\,dy\,dz+(y^3-yz^2)\,dz\,dx+(z^3-zx^2)\,dx\,dy.
  \]
\end{enumerate}

\subsubsection*{三、解答题（每小题 9 分，共 18 分）}
\begin{enumerate}
  \item 将函数
  \[
  \frac{d}{dx}\left(\frac{e^x-1}{x}\right)
  \]
  展开成 $x$ 的幂级数，并求数项级数
  \[
  \sum_{n=1}^{\infty}\frac{n}{(n+1)!}
  \]
  的和。
  \item 求微分方程
  \[
  y'\tan x=y\ln y
  \]
  的通解。
\end{enumerate}

\subsubsection*{四、证明题（每小题 9 分，共 18 分）}
\begin{enumerate}
  \item 设 $f(x)\in C[a,b]$，证明：
  \[
  \int_a^b dx\int_a^x f(y)\,dy=\int_a^b f(y)(b-y)\,dy.
  \]
  \item 证明：函数项级数
  \[
  \sum_{n=1}^{\infty}\frac{(-1)^n(1-e^{-nx})}{n^2+x^2}
  \]
  在 $[0,+\infty)$ 上一致收敛。
\end{enumerate}

\subsubsection*{五、应用题（本题 9 分）}
从斜边之长为 $a$ 的一切直角三角形中，求有最大周长的直角三角形。

\subsection{答案与解析}

\subsubsection*{一、填空题}
\begin{enumerate}
  \item 特征方程 $r^2-4=0$ 有根 $r=\pm2$，右端为一次多项式乘 $e^{2x}$，且 $2$ 是一重特征根，故特解形式为
  \[
  y^*=xe^{2x}(Ax+B).
  \]
  \item 记 $u=y-bz$。对 $x,y$ 分别求偏导：
  \[
  1-az_x=-b\varphi'(u)z_x,\qquad -az_y=\varphi'(u)(1-bz_y).
  \]
  化简得
  \[
  a z_x+b z_y=1.
  \]
  \item 最大方向为梯度方向。因
  \[
  \nabla f=-(yz,xz,xy)\sin(xyz),
  \]
  在 $P\!\left(\dfrac13,1,\pi\right)$ 处有
  \[
  \nabla f(P)=-\frac{\sqrt3}{6}(3\pi,\pi,1).
  \]
  故所求单位方向可取
  \[
  -\frac{(3\pi,\pi,1)}{\sqrt{10\pi^2+1}}.
  \]
  \item 由 Green 公式，
  \[
  \int_L 2xy\,dx+(x^2+8y^2)\,dy
  =\iint_D(2x-2x)\,dA=0.
  \]
  \item 幂级数在 $x=-3$ 条件收敛，说明 $x=-3$ 是收敛区间端点，故收敛半径为
  \[
  R=3.
  \]
\end{enumerate}

\subsubsection*{二、计算题}
\begin{enumerate}
  \item 令 $\xi=xy,\ \eta=\dfrac{x}{y}$。先求
  \[
  z_x=yf_\xi+\frac1y f_\eta.
  \]
  再对 $y$ 求偏导，混合项相消，得
  \[
  z_{xy}=f_\xi+xyf_{\xi\xi}-\frac{x}{y^3}f_{\eta\eta}-\frac1{y^2}f_\eta,
  \]
  其中各偏导均在 $\left(xy,\dfrac{x}{y}\right)$ 处取值。
  \item $\Gamma$ 是半径为 $a$ 的大圆，所在平面的单位法向量为
  \[
  \mathbf n=\frac{(1,1,1)}{\sqrt3}.
  \]
  单位向量 $\mathbf e_x$ 在该平面内的投影长度平方为 $1-\dfrac13=\dfrac23$。故圆周上 $x^2$ 的平均值为
  \[
  \frac{a^2}{2}\cdot\frac23=\frac{a^2}{3}.
  \]
  因此
  \[
  I=2\pi a\cdot\frac{a^2}{3}=\frac{2\pi a^3}{3}.
  \]
  \item 记
  \[
  P=(f(x)-e^x)\sin y,\qquad Q=-f'(x)\cos y.
  \]
  路径无关给出 $P_y=Q_x$，即
  \[
  f''(x)+f(x)=e^x.
  \]
  通解为
  \[
  f(x)=C_1\cos x+C_2\sin x+\frac12e^x.
  \]
  又 $f'(0)=0$，得 $C_2=-\dfrac12$，于是
  \[
  f'(x)=-C_1\sin x-\frac12\cos x+\frac12e^x.
  \]
  因为
  \[
  (f-e^x)\sin y\,dx-f'(x)\cos y\,dy=d\bigl[-f'(x)\sin y\bigr],
  \]
  所以
  \[
  I=-f'(1)\sin1
  =C_1\sin^2 1+\frac12(\cos1-e)\sin1.
  \]
  原题只给出 $f'(0)=0$，只能确定 $C_2=-\dfrac12$，不能唯一确定 $C_1$。因此在现有题面条件下，积分值应保留为上述含任意常数 $C_1$ 的表达式；若要得到唯一数值，还需另给一个关于 $f$ 或 $f'$ 的独立条件。
  \item 将平面写为 $z=5-y$，则
  \[
  dS=\sqrt2\,dx\,dy,\qquad x+y+z=x+5.
  \]
  投影区域为圆盘 $D:x^2+y^2\le25$，故
  \[
  I=\sqrt2\iint_D(x+5)\,dx\,dy
  =125\sqrt2\,\pi.
  \]
  \item 令
  \[
  \mathbf F=(x^3-xy^2,\ y^3-yz^2,\ z^3-zx^2).
  \]
  则
  \[
  \nabla\cdot\mathbf F
  =(3x^2-y^2)+(3y^2-z^2)+(3z^2-x^2)
  =2(x^2+y^2+z^2).
  \]
  上半球面与底面圆盘围成上半球体，底面上 $z=0$，第三分量通量为 $0$。于是
  \[
  I=\iiint_\Omega 2(x^2+y^2+z^2)\,dV
  =2\cdot 2\pi\int_0^a r^4\,dr
  =\frac{4\pi a^5}{5}.
  \]
\end{enumerate}

\subsubsection*{三、解答题}
\begin{enumerate}
  \item 先写
  \[
  \frac{e^x-1}{x}=\sum_{n=0}^{\infty}\frac{x^n}{(n+1)!}.
  \]
  逐项求导得
  \[
  \frac{d}{dx}\left(\frac{e^x-1}{x}\right)
  =\sum_{n=1}^{\infty}\frac{n}{(n+1)!}x^{n-1}.
  \]
  令 $x=1$，也可由闭式
  \[
  \frac{d}{dx}\left(\frac{e^x-1}{x}\right)
  =\frac{(x-1)e^x+1}{x^2}
  \]
  得
  \[
  \sum_{n=1}^{\infty}\frac{n}{(n+1)!}=1.
  \]
  \item 设 $y>0$，方程化为
  \[
  \frac{dy}{y\ln y}=\cot x\,dx.
  \]
  积分得
  \[
  \ln|\ln y|=\ln|\sin x|+C,
  \]
  即
  \[
  \ln y=C\sin x,\qquad y=e^{C\sin x}.
  \]
  其中 $C=0$ 时包含常值解 $y=1$。
\end{enumerate}

\subsubsection*{四、证明题}
\begin{enumerate}
  \item 左端积分区域为 $a\le y\le x\le b$。交换积分次序得
  \[
  \int_a^b dx\int_a^x f(y)\,dy
  =\int_a^b dy\int_y^b f(y)\,dx
  =\int_a^b f(y)(b-y)\,dy.
  \]
  \item 对 $x\in[0,+\infty)$，有
  \[
  0\le 1-e^{-nx}\le1,\qquad n^2+x^2\ge n^2.
  \]
  因而
  \[
  \left|\frac{(-1)^n(1-e^{-nx})}{n^2+x^2}\right|\le\frac1{n^2}.
  \]
  由 $\sum_{n=1}^{\infty}\dfrac1{n^2}$ 收敛及 Weierstrass 判别法，原级数在 $[0,+\infty)$ 上一致收敛。
\end{enumerate}

\subsubsection*{五、应用题}
设直角三角形两直角边为 $x,y>0$，则
\[
x^2+y^2=a^2,\qquad P=a+x+y.
\]
在圆弧约束下，$x+y$ 最大时 $x=y$，故
\[
x=y=\frac{a}{\sqrt2}.
\]
因此所求三角形为等腰直角三角形，最大周长为
\[
P_{\max}=a+\sqrt2\,a=(1+\sqrt2)a.
\]

\section{2016级工科数学分析下 A卷}

\subsection{资料来源}
\begin{itemize}
  \item 试题：2016 级工科数学分析下 A 卷，已导出页面图校正公式。
  \item 未发现配套 A 卷答案或解析文本；下列答案与解析为据题面补写。
\end{itemize}

\subsection{试题}

\paragraph*{试卷信息}
诚信应考，考试作弊将带来严重后果！

华南理工大学本科生期末考试《工科数学分析（二）》2016--2017学年第二学期期末考试试卷 A 卷。

注意事项：开考前请将密封线内各项信息填写清楚；所有答案请直接答在试卷上；考试形式：闭卷；本试卷共5大题，满分100分，考试时间120分钟。

\paragraph*{一、填空题（共5题，每题2分，共10分）}
\begin{enumerate}
  \item 函数
  \[
    f(x,y)=2x^2+y^2
  \]
  在点 \((1,1)\) 处沿该点的梯度方向的方向导数为 \underline{\hspace{5cm}}。
  \item 向量场 \((2x-3y)\mathbf i+(3x-z)\mathbf j+(y-2x)\mathbf k\) 的旋度向量为 \underline{\hspace{5cm}}。
  \item 设 \(\Sigma\) 表示球面 \(x^2+y^2+z^2=R^2\) 的外侧，则第二类曲面积分
  \[
    \iint_\Sigma x\,\mathrm{d}y\mathrm{d}z+y\,\mathrm{d}z\mathrm{d}x+z\,\mathrm{d}x\mathrm{d}y=
  \]
  \underline{\hspace{5cm}}。
  \item 初值问题
  \[
    \begin{cases}
      y'+y\cos x=\mathrm{e}^{-\sin x},\\
      y|_{x=0}=1
    \end{cases}
  \]
  的解为 \underline{\hspace{5cm}}。
  \item 设 \(f(x)\) 是周期为 \(2\) 的周期函数，它在 \((-1,1]\) 上的表达式为
  \[
    f(x)=
    \begin{cases}
      2, & -1<x\le 0,\\
      x^3+1, & 0<x\le 1,
    \end{cases}
  \]
  则 \(f(x)\) 的傅里叶级数在 \(x=0\) 处收敛于 \underline{\hspace{5cm}}。
\end{enumerate}

\paragraph*{二、单选题（每题只有一个正确选项，共5题，每题2分，共10分）}
\begin{enumerate}
  \item 二元函数
  \[
    f(x,y)=
    \begin{cases}
      \dfrac{xy}{x^2+y^2}, & (x,y)\ne(0,0),\\
      0, & (x,y)=(0,0)
    \end{cases}
  \]
  在点 \((0,0)\) 处（\quad）。
  \begin{enumerate}
    \item[A.] 连续，偏导数存在。
    \item[B.] 不连续，偏导数存在。
    \item[C.] 连续，偏导数不存在。
    \item[D.] 不连续，偏导数不存在。
  \end{enumerate}

  \item 曲线
  \[
    \begin{cases}
      x^2+y^2+z^2=9,\\
      x-y+z=1
    \end{cases}
  \]
  在点 \(M(1,2,2)\) 处的切线一定平行于（\quad）。
  \begin{enumerate}
    \item[A.] \(Oxy\) 平面。
    \item[B.] \(Oyz\) 平面。
    \item[C.] \(Ozx\) 平面。
    \item[D.] 平面 \(x-y+z=1\)。
  \end{enumerate}

  \item 设 \(D\) 是一个有界的平面闭区域，其边界曲线 \(\Gamma\) 分段光滑，则下列积分值不等于区域 \(D\) 的面积的是（\quad）。
  \begin{enumerate}
    \item[A.] \(\displaystyle \oint_\Gamma y\,\mathrm{d}x\)
    \item[B.] \(\displaystyle \frac12\oint_\Gamma x\,\mathrm{d}y-y\,\mathrm{d}x\)
    \item[C.] \(\displaystyle \oint_\Gamma x\,\mathrm{d}y\)
    \item[D.] \(\displaystyle \iint_D 1\,\mathrm{d}x\mathrm{d}y\)
  \end{enumerate}

  \item 关于未知函数 \(y\) 的微分方程 \((y-\sin x)\,\mathrm{d}x+x\,\mathrm{d}y=0\) 是（\quad）。
  \begin{enumerate}
    \item[A.] 可分离变量方程。
    \item[B.] 一阶非齐次线性方程。
    \item[C.] 一阶齐次线性方程。
    \item[D.] 非线性方程。
  \end{enumerate}

  \item 使得级数 \(\displaystyle \sum_{n=1}^{\infty}\frac{(-1)^n}{n^p}\) 条件收敛的常数 \(p\) 的取值范围是（\quad）。
  \begin{enumerate}
    \item[A.] \(p\le 0\)
    \item[B.] \(0<p<1\)
    \item[C.] \(0<p\le 1\)
    \item[D.] \(p>1\)
  \end{enumerate}
\end{enumerate}

\paragraph*{三、计算题（共4题，每题7分，共28分）}
\begin{enumerate}
  \item 设函数 \(u=f(x,y,z)\) 有连续的偏导数，函数 \(y=y(x)\) 和 \(z=z(x)\) 分别由方程
  \(\mathrm{e}^{xy}=y\) 和 \(\mathrm{e}^z=xz\) 确定，计算 \(\dfrac{\mathrm{d}u}{\mathrm{d}x}\)。

  \item 计算累次积分
  \[
    \int_1^2\mathrm{d}x\int_{\sqrt{x}}^x\sin\frac{\pi x}{2y}\,\mathrm{d}y
    +\int_2^4\mathrm{d}x\int_{\sqrt{x}}^2\sin\frac{\pi x}{2y}\,\mathrm{d}y.
  \]

  \item 计算三重积分 \(\displaystyle \iiint_\Omega z\,\mathrm{d}x\mathrm{d}y\mathrm{d}z\)，其中 \(\Omega\) 是由上半球面
  \(x^2+y^2+z^2=R^2\)，\(z\ge 0\)，\(R>0\) 和锥面 \(z=\sqrt{x^2+y^2}\) 所围成的闭区域。

  \item 计算第二类曲线积分
  \[
    \oint_\Gamma xy^2\,\mathrm{d}y-x^2y\,\mathrm{d}x,
  \]
  其中 \(\Gamma\) 为椭圆 \(\dfrac{x^2}{4}+y^2=1\)，取逆时针方向。
\end{enumerate}

\paragraph*{四、解答题（共4题，每题8分，共32分）}
\begin{enumerate}
  \item 设二阶常系数线性微分方程
  \[
    y''+\alpha y'+\beta y=-\mathrm{e}^x
  \]
  的一个特解为 \(y=(1+x)\mathrm{e}^x\)，试确定常数 \(\alpha,\beta\)，并求该方程的通解。

  \item 已知螺旋形弹簧一圈的方程为
  \[
    x=a\cos t,\qquad y=a\sin t,\qquad z=bt,\qquad 0\le t\le 2\pi,
  \]
  其中 \(a,b\) 为大于零的常数，且弹簧上各点处的线密度等于该点到 \(Oxy\) 平面的距离，求此弹簧的质心坐标。

  \item 求上半球面 \(x^2+y^2+z^2=a^2,\ z\ge 0\) 被柱面 \(x^2+y^2=ax\) 截下的部分的面积。

  \item 将函数 \(f(x)=\arctan x\) 展开为 \(x\) 的幂级数。
\end{enumerate}

\paragraph*{五、证明题（共2题，每题6分，共12分）}
\begin{enumerate}
  \item 设函数 \(f(\xi,\eta)\) 具有连续的二阶偏导数，且满足
  \[
    \frac{\partial^2f}{\partial \xi^2}+\frac{\partial^2f}{\partial \eta^2}=0.
  \]
  证明：函数 \(z=f(x^2-y^2,2xy)\) 满足
  \[
    \frac{\partial^2z}{\partial x^2}+\frac{\partial^2z}{\partial y^2}=0.
  \]

  \item 证明：函数项级数 \(\displaystyle \sum_{n=0}^{\infty}(1-x)x^n\) 在区间 \((0,1)\) 上点态收敛，但不一致收敛。
\end{enumerate}

\paragraph*{六、应用题（共1题，共8分）}
求曲线
\[
  \begin{cases}
    z=x^2+y^2,\\
    y=\dfrac1x
  \end{cases}
\]
上到 \(Oxy\) 平面距离最近的点。

\subsection{答案与解析}
\paragraph*{一、填空题}
\begin{enumerate}
  \item \(\nabla f(1,1)=(4,2)\)，沿梯度方向的方向导数为 \(\|\nabla f(1,1)\|=2\sqrt5\)。
  \item
  \[
    \nabla\times\mathbf F=
    \begin{vmatrix}
      \mathbf i&\mathbf j&\mathbf k\\
      \partial_x&\partial_y&\partial_z\\
      2x-3y&3x-z&y-2x
    \end{vmatrix}
    =2\mathbf i+2\mathbf j+6\mathbf k.
  \]
  \item 由 Gauss 公式，
  \[
    \iint_\Sigma x\,\mathrm{d}y\mathrm{d}z+y\,\mathrm{d}z\mathrm{d}x+z\,\mathrm{d}x\mathrm{d}y
    =\iiint_V3\,\mathrm{d}V=4\pi R^3.
  \]
  \item 积分因子为 \(\mathrm{e}^{\sin x}\)，故
  \[
    (y\mathrm{e}^{\sin x})'=1,\qquad y=(x+1)\mathrm{e}^{-\sin x}.
  \]
  \item 傅里叶级数在跳跃点收敛于左右极限平均值：
  \[
    \frac{f(0-0)+f(0+0)}2=\frac{2+1}{2}=\frac32.
  \]
\end{enumerate}

\paragraph*{二、单选题}
\begin{enumerate}
  \item 选 B。沿 \(y=x\) 趋于原点时 \(f(x,x)=1/2\)，故不连续；沿坐标轴计算偏导数均存在。
  \item 选 D。切向量垂直于两个曲面的法向量
  \[
    (2,4,4)\times(1,-1,1)=(8,2,-6),
  \]
  它与平面 \(x-y+z=1\) 的法向量 \((1,-1,1)\) 垂直，故切线平行于该平面。
  \item 选 A。正向边界下 \(\oint_\Gamma y\,\mathrm{d}x=-S_D\)，而其余三项均等于区域面积。
  \item 选 B。方程化为
  \[
    y'+\frac1x y=\frac{\sin x}{x},
  \]
  是一阶非齐次线性方程。
  \item 选 C。交错级数收敛要求 \(p>0\)，绝对收敛要求 \(p>1\)，故条件收敛为 \(0<p\le 1\)。
\end{enumerate}

\paragraph*{三、计算题}
\begin{enumerate}
  \item 由 \(\mathrm{e}^{xy}=y\) 得
  \[
    y'(1-xy)=y^2,\qquad y'=\frac{y^2}{1-xy}.
  \]
  由 \(\mathrm{e}^z=xz\) 得
  \[
    z'(\mathrm{e}^z-x)=z,\qquad z'=\frac{z}{x(z-1)}.
  \]
  因此
  \[
    \frac{\mathrm{d}u}{\mathrm{d}x}
    =f_x+f_y\frac{y^2}{1-xy}+f_z\frac{z}{x(z-1)}.
  \]

  \item 积分区域合并后为
  \[
    1\le y\le 2,\qquad y\le x\le y^2.
  \]
  所以
  \[
    I=\int_1^2\mathrm{d}y\int_y^{y^2}\sin\frac{\pi x}{2y}\,\mathrm{d}x
    =-\frac2\pi\int_1^2 y\cos\frac{\pi y}{2}\,\mathrm{d}y
    =\frac4{\pi^2}+\frac8{\pi^3}.
  \]

  \item 用球坐标，区域为
  \[
    0\le r\le R,\qquad 0\le\varphi\le\frac{\pi}{4},\qquad 0\le\theta\le2\pi.
  \]
  故
  \[
    \iiint_\Omega z\,\mathrm{d}V
    =\int_0^{2\pi}\!\!\int_0^{\pi/4}\!\!\int_0^R r\cos\varphi\cdot r^2\sin\varphi\,\mathrm{d}r\,\mathrm{d}\varphi\,\mathrm{d}\theta
    =\frac{\pi R^4}{8}.
  \]

  \item 令 \(P=-x^2y\)，\(Q=xy^2\)。由 Green 公式，
  \[
    \oint_\Gamma P\,\mathrm{d}x+Q\,\mathrm{d}y
    =\iint_D(x^2+y^2)\,\mathrm{d}A.
  \]
  对椭圆作变换 \(x=2r\cos\theta,\ y=r\sin\theta\)，得
  \[
    \iint_D(x^2+y^2)\,\mathrm{d}A
    =\int_0^{2\pi}\!\!\int_0^1(4r^2\cos^2\theta+r^2\sin^2\theta)\,2r\,\mathrm{d}r\,\mathrm{d}\theta
    =\frac{5\pi}{2}.
  \]
\end{enumerate}

\paragraph*{四、解答题}
\begin{enumerate}
  \item 将 \(y=(1+x)\mathrm{e}^x\) 代入方程，得
  \[
    \bigl[(1+\alpha+\beta)x+(3+2\alpha+\beta)\bigr]\mathrm{e}^x=-\mathrm{e}^x.
  \]
  所以
  \[
    1+\alpha+\beta=0,\qquad 3+2\alpha+\beta=-1,
  \]
  解得 \(\alpha=-3,\ \beta=2\)。方程为
  \[
    y''-3y'+2y=-\mathrm{e}^x,
  \]
  齐次解为 \(C_1\mathrm{e}^x+C_2\mathrm{e}^{2x}\)，故通解可写为
  \[
    y=C_1\mathrm{e}^x+C_2\mathrm{e}^{2x}+x\mathrm{e}^x.
  \]

  \item 有
  \[
    \mathrm{d}s=\sqrt{a^2+b^2}\,\mathrm{d}t,\qquad \rho=bt.
  \]
  质量
  \[
    m=\int_0^{2\pi}bt\sqrt{a^2+b^2}\,\mathrm{d}t
    =2\pi^2b\sqrt{a^2+b^2}.
  \]
  因而
  \[
    \bar x=0,\qquad
    \bar y=\frac{\int_0^{2\pi}a\sin t\cdot bt\sqrt{a^2+b^2}\,\mathrm{d}t}{m}
    =-\frac a\pi,
  \]
  \[
    \bar z=\frac{\int_0^{2\pi}bt\cdot bt\sqrt{a^2+b^2}\,\mathrm{d}t}{m}
    =\frac{4\pi b}{3}.
  \]
  所以质心为
  \[
    \left(0,-\frac a\pi,\frac{4\pi b}{3}\right).
  \]

  \item 投影区域为柱面内部
  \[
    D:\ 0\le r\le a\cos\theta,\qquad -\frac{\pi}{2}\le\theta\le\frac{\pi}{2}.
  \]
  上半球面 \(z=\sqrt{a^2-r^2}\) 上
  \[
    \mathrm{d}S=\frac{a}{\sqrt{a^2-r^2}}r\,\mathrm{d}r\,\mathrm{d}\theta.
  \]
  故面积
  \[
    S=\int_{-\pi/2}^{\pi/2}\int_0^{a\cos\theta}
    \frac{ar}{\sqrt{a^2-r^2}}\,\mathrm{d}r\,\mathrm{d}\theta
    =a^2(\pi-2).
  \]

  \item 由
  \[
    \frac1{1+x^2}=\sum_{n=0}^{\infty}(-1)^n x^{2n}\qquad (|x|<1)
  \]
  逐项积分得
  \[
    \arctan x=\sum_{n=0}^{\infty}\frac{(-1)^n x^{2n+1}}{2n+1},
    \qquad |x|\le 1.
  \]
\end{enumerate}

\paragraph*{五、证明题}
\begin{enumerate}
  \item 记 \(\xi=x^2-y^2,\ \eta=2xy\)。链式法则给出
  \[
    \frac{\partial^2z}{\partial x^2}+\frac{\partial^2z}{\partial y^2}
    =4(x^2+y^2)\left(\frac{\partial^2f}{\partial \xi^2}+\frac{\partial^2f}{\partial \eta^2}\right)=0.
  \]

  \item 对固定 \(x\in(0,1)\)，
  \[
    \sum_{n=0}^{\infty}(1-x)x^n=1,
  \]
  因而点态收敛。第 \(N\) 项部分和为
  \[
    S_N(x)=1-x^{N+1},
  \]
  余项为 \(x^{N+1}\)。由于
  \[
    \sup_{0<x<1}|1-S_N(x)|=\sup_{0<x<1}x^{N+1}=1,
  \]
  余项不一致趋于 \(0\)，故不一致收敛。
\end{enumerate}

\paragraph*{六、应用题}
曲线上点到 \(Oxy\) 平面的距离为 \(|z|\)。由题设
\[
  z=x^2+\frac1{x^2}\qquad (x\ne0),
\]
故只需求 \(z\) 的最小值。令
\[
  z'=2x-\frac2{x^3}=0,
\]
得 \(x^4=1\)，即 \(x=\pm1\)。此时 \(y=1/x\)，\(z=2\)。所以最近点为
\[
  (1,1,2),\qquad (-1,-1,2).
\]

\section{2016级工科数学分析下 B卷}

\subsection{资料来源}
\begin{itemize}
  \item 2016级工科数学分析试卷B.doc（试题骨架，DOC 文本抽取中大量公式对象缺失，仅作题型与题号参考）。
  \item 2016级工科数学分析期末试卷B解答.pdf（试题与答案，已按页面图和文本层人工校对）。
\end{itemize}

\subsection{试题}

\paragraph*{一、填空题（共 5 题，每题 2 分，共 10 分）}
\begin{enumerate}
  \item 由方程
  \[
    xyz+\sqrt{x^2+y^2+z^2}=\sqrt2
  \]
  所确定的函数 $z=z(x,y)$ 在点 $(1,0,-1)$ 处的全微分 $dz=$\underline{\hspace{4cm}}。

  \item 设 $\Gamma$ 为椭圆
  \[
    \frac{x^2}{9}+\frac{y^2}{4}=1
  \]
  取逆时针方向，则第二类曲线积分
  \[
    \oint_\Gamma x\,dy-y\,dx=
  \]
  \underline{\hspace{4cm}}。

  \item 向量场
  \[
    x^2yz\,\mathbf i+xy^2z\,\mathbf j+xyz^2\,\mathbf k
  \]
  在点 $(1,1,1)$ 的散度为\underline{\hspace{4cm}}。

  \item 初值问题
  \[
    \begin{cases}
      y'+\dfrac{3}{x}y=\dfrac{2}{x^3},\\
      y|_{x=1}=1
    \end{cases}
  \]
  的解为\underline{\hspace{4cm}}。

  \item 设 $f(x)=\pi x+x^2,\ -\pi<x<\pi$ 的傅里叶级数为
  \[
    \frac{a_0}{2}+\sum_{n=1}^{\infty}(a_n\cos nx+b_n\sin nx),
  \]
  则其中系数 $b_2=$\underline{\hspace{4cm}}。
\end{enumerate}

\paragraph*{二、单选题（共 5 题，每题 2 分，只有一个正确选项，共 10 分）}
\begin{enumerate}
  \item 二元函数 $f(x,y)$ 在点 $(a,b)$ 处的两个偏导数
  \[
    \frac{\partial f}{\partial x}(a,b),\qquad
    \frac{\partial f}{\partial y}(a,b)
  \]
  存在，则（\quad）。
  \begin{enumerate}
    \item[A.] $f(x,y)$ 在点 $(a,b)$ 处连续；
    \item[B.] $f(x,y)$ 在点 $(a,b)$ 处可微；
    \item[C.] $\displaystyle \lim_{x\to a}f(x,b)$ 和 $\displaystyle \lim_{y\to b}f(a,y)$ 都存在；
    \item[D.] $\displaystyle \lim_{(x,y)\to(a,b)}f(x,y)$ 存在。
  \end{enumerate}

  \item 已知曲面 $z=4-x^2-y^2$ 上点 $P$ 处的切平面平行于平面 $2x+2y+z=1$，则 $P$ 点的坐标是（\quad）。
  \begin{enumerate}
    \item[A.] $(1,-1,2)$；
    \item[B.] $(-1,1,2)$；
    \item[C.] $(1,1,2)$；
    \item[D.] $(-1,-1,2)$。
  \end{enumerate}

  \item 一个均匀物体由曲面 $z=x^2+y^2$ 及 $z=1$ 围成，则该物体的质心坐标为（\quad）。
  \begin{enumerate}
    \item[A.] $\left(\dfrac23,\dfrac23,\dfrac23\right)$；
    \item[B.] $\left(0,0,\dfrac23\right)$；
    \item[C.] $(0,0,1)$；
    \item[D.] $\left(\dfrac23,0,0\right)$。
  \end{enumerate}

  \item 微分方程
  \[
    (y-\ln x)\,dx+x\,dy=0
  \]
  是（\quad）。
  \begin{enumerate}
    \item[A.] 可分离变量方程；
    \item[B.] 一阶非齐次线性方程；
    \item[C.] 一阶齐次线性方程；
    \item[D.] 非线性方程。
  \end{enumerate}

  \item 下列级数条件收敛的是（\quad）。
  \begin{enumerate}
    \item[A.] $\displaystyle \sum_{n=1}^{\infty}\frac{(-1)^n}{n^2}$；
    \item[B.] $\displaystyle \sum_{n=1}^{\infty}\left(1-\cos\frac1n\right)$；
    \item[C.] $\displaystyle \sum_{n=1}^{\infty}\frac{n!}{n^n}$；
    \item[D.] $\displaystyle \sum_{n=1}^{\infty}\frac{(-1)^n}{n}$。
  \end{enumerate}
\end{enumerate}

\paragraph*{三、计算题（共 4 题，每题 7 分，共 28 分）}
\begin{enumerate}
  \item 设函数
  \[
    u(x,y)=y f\left(\frac{x}{y}\right)+x g\left(\frac{y}{x}\right),
  \]
  其中 $f$ 和 $g$ 具有连续的二阶导数，计算
  \[
    x\frac{\partial^2u}{\partial x^2}+y\frac{\partial^2u}{\partial x\partial y}.
  \]

  \item 计算二重积分
  \[
    \iint_D \frac{1-x^2-y^2}{1+x^2+y^2}\,dx\,dy,
  \]
  其中 $D$ 是区域 $x^2+y^2\le1,\ x\ge0,\ y\ge0$。

  \item 计算三重积分
  \[
    \iiint_\Omega z^2\,dx\,dy\,dz,
  \]
  其中 $\Omega$ 是平面 $x+y+z=1$ 与三个坐标面围成的区域。

  \item 计算第一类曲面积分
  \[
    \iint_\Sigma (x+y+z)\,dS,
  \]
  其中 $\Sigma$ 为平面 $y+z=5$ 被柱面 $x^2+y^2=25$ 所截得的部分。
\end{enumerate}

\paragraph*{四、解答题（共 4 题，每题 8 分，共 32 分）}
\begin{enumerate}
  \item 求方程
  \[
    y''+4y'+3y=0
  \]
  的一个解 $y=y(x)$，使其在点 $(0,2)$ 处与直线 $x-y+2=0$ 相切。

  \item 求圆柱面 $x^2+y^2=4,\ x\ge0,\ y\ge0$ 介于 $xOy$ 平面和曲面 $z=xy$ 之间的部分的面积。

  \item 设 $\Sigma$ 是锥面 $z=\sqrt{x^2+y^2}$ 被平面 $z=0$ 及 $z=1$ 截下的部分的下侧，计算第二类曲面积分
  \[
    \iint_\Sigma x\,dy\,dz+y\,dz\,dx+(z^2-2z)\,dx\,dy.
  \]

  \item 求幂级数
  \[
    \sum_{n=1}^{\infty}\frac{x^n}{n}
  \]
  的收敛域，并在收敛域上求其和函数。
\end{enumerate}

\paragraph*{五、证明题（共 2 题，每题 6 分，共 12 分）}
\begin{enumerate}
  \item 证明：曲面
  \[
    z=xe^{y/x}
  \]
  上所有点处的切平面都过一定点。

  \item 证明：函数项级数
  \[
    \sum_{n=0}^{\infty}x^n
  \]
  在 $(0,1)$ 上点态收敛，但不一致收敛。
\end{enumerate}

\paragraph*{六、应用题（共 1 题，共 8 分）}
从斜边长为 $l$ 的直角三角形中，求周长最大的直角三角形。

\subsection{答案与解析}

\paragraph*{一、填空题}
\begin{enumerate}
  \item 设
  \[
    F(x,y,z)=xyz+\sqrt{x^2+y^2+z^2}-\sqrt2.
  \]
  在点 $(1,0,-1)$ 处，
  \[
    F_x=\frac1{\sqrt2},\qquad F_y=-1,\qquad F_z=-\frac1{\sqrt2}.
  \]
  因此
  \[
    z_x=-\frac{F_x}{F_z}=1,\qquad z_y=-\frac{F_y}{F_z}=-\sqrt2,
  \]
  所以 $dz=dx-\sqrt2\,dy$。

  \item 由格林公式，
  \[
    \oint_\Gamma x\,dy-y\,dx=2S,
  \]
  其中椭圆面积 $S=\pi\cdot3\cdot2=6\pi$，故结果为 $12\pi$。

  \item
  \[
    \operatorname{div}(x^2yz,xy^2z,xyz^2)
    =2xyz+2xyz+2xyz=6xyz.
  \]
  在 $(1,1,1)$ 处取值为 $6$。

  \item 原方程乘以积分因子 $x^3$，得
  \[
    (x^3y)'=2.
  \]
  故 $x^3y=2x+C$。由 $y(1)=1$ 得 $C=-1$，所以
  \[
    y=\frac{2}{x^2}-\frac{1}{x^3}.
  \]

  \item
  \[
    b_n=\frac1\pi\int_{-\pi}^{\pi}(\pi x+x^2)\sin nx\,dx.
  \]
  由于 $x^2\sin nx$ 为奇函数，其积分为 $0$，故
  \[
    b_n=2\int_0^\pi x\sin nx\,dx=-\frac{2\pi(-1)^n}{n}.
  \]
  取 $n=2$，得 $b_2=-\pi$。
\end{enumerate}

\paragraph*{二、单选题}
\begin{enumerate}
  \item 两个偏导数存在只能保证沿坐标轴方向的一元极限存在，不保证二元连续或可微。选 C。
  \item 曲面法向量可取 $(2x,2y,1)$，与平面 $2x+2y+z=1$ 的法向量 $(2,2,1)$ 平行，得 $x=1,\ y=1,\ z=2$。选 C。
  \item 区域关于 $z$ 轴旋转对称，质心的 $x,y$ 坐标为 $0$；用柱坐标算得 $\bar z=\dfrac23$。选 B。
  \item 方程可化为 \(y'+\dfrac1x y=\dfrac{\ln x}{x}\)，是一阶非齐次线性方程。选 B。
  \item A 绝对收敛，B、C 收敛，D 是交错调和级数，只条件收敛。选 D。
\end{enumerate}

\paragraph*{三、计算题}
\begin{enumerate}
  \item 先求一阶偏导：
  \[
    u_x=f'\left(\frac{x}{y}\right)-\frac{y}{x}g'\left(\frac{y}{x}\right),
  \]
  \[
    u_y=-\frac{x}{y}f'\left(\frac{x}{y}\right)+g'\left(\frac{y}{x}\right).
  \]
  继续求导得
  \[
    u_{xx}=\frac1y f''\left(\frac{x}{y}\right)
    +\frac{y^2}{x^3}g''\left(\frac{y}{x}\right),
  \]
  \[
    u_{xy}=-\frac{x}{y^2}f''\left(\frac{x}{y}\right)
    -\frac{y}{x^2}g''\left(\frac{y}{x}\right).
  \]
  因此
  \[
    x u_{xx}+y u_{xy}=0.
  \]

  \item 作极坐标变换
  \[
    x=r\cos\theta,\qquad y=r\sin\theta,\qquad 0\le r\le1,\quad 0\le\theta\le\frac{\pi}{2}.
  \]
  则
  \[
    \iint_D \frac{1-x^2-y^2}{1+x^2+y^2}\,dx\,dy
    =
    \int_0^{\pi/2}d\theta\int_0^1\frac{1-r^2}{1+r^2}r\,dr.
  \]
  因为
  \[
    \int_0^1\frac{1-r^2}{1+r^2}r\,dr
    =\int_0^1\left(\frac{2r}{1+r^2}-r\right)\,dr
    =\ln2-\frac12,
  \]
  所以原积分为
  \[
    \frac{\pi}{4}(2\ln2-1).
  \]

  \item 对 $z$ 先积分外层，得
  \[
    \iiint_\Omega z^2\,dx\,dy\,dz
    =
    \int_0^1 z^2
    \left(\iint_{x+y\le1-z,\ x\ge0,\ y\ge0}dx\,dy\right)dz.
  \]
  截面面积为 $\dfrac12(1-z)^2$，因此
  \[
    \iiint_\Omega z^2\,dx\,dy\,dz
    =
    \frac12\int_0^1 z^2(1-z)^2\,dz
    =
    \frac1{60}.
  \]

  \item 曲面有显式方程 $z=5-y$，故
  \[
    dS=\sqrt{1+z_x^2+z_y^2}\,dx\,dy=\sqrt2\,dx\,dy.
  \]
  投影区域为 $x^2+y^2\le25$。于是
  \[
    \iint_\Sigma (x+y+z)\,dS
    =
    \iint_{x^2+y^2\le25}(x+5)\sqrt2\,dx\,dy.
  \]
  其中 $\iint x\,dx\,dy=0$，圆盘面积为 $25\pi$，故结果为
  \[
    125\sqrt2\,\pi.
  \]
\end{enumerate}

\paragraph*{四、解答题}
\begin{enumerate}
  \item 特征方程为
  \[
    \lambda^2+4\lambda+3=0,
  \]
  根为 $\lambda_1=-3,\ \lambda_2=-1$，故通解为
  \[
    y=C_1e^{-3x}+C_2e^{-x}.
  \]
  直线 $x-y+2=0$ 的斜率为 $1$，因此
  \[
    y(0)=2,\qquad y'(0)=1.
  \]
  即
  \[
    C_1+C_2=2,\qquad -3C_1-C_2=1.
  \]
  解得
  \[
    C_1=-\frac32,\qquad C_2=\frac72.
  \]
  所求解为
  \[
    y=-\frac32e^{-3x}+\frac72e^{-x}.
  \]

  \item 该柱面片可以看成以四分之一圆弧
  \[
    \Gamma:\ x^2+y^2=4,\quad x\ge0,\quad y\ge0
  \]
  为底线、以高度 $z=xy$ 竖直生成的曲面，故面积为
  \[
    S=\int_\Gamma xy\,ds.
  \]
  令
  \[
    x=2\cos t,\qquad y=2\sin t,\qquad 0\le t\le\frac{\pi}{2},
  \]
  则 $ds=2\,dt$，所以
  \[
    S=\int_0^{\pi/2}4\cos t\sin t\cdot2\,dt=4.
  \]

  \item 令 $\Sigma_1$ 为平面 $z=1$ 被锥面截下的圆盘，并取上侧，则 $\Sigma$ 与 $\Sigma_1$ 围成区域
  \[
    \Omega=\{(x,y,z):x^2+y^2\le z^2,\ 0\le z\le1\}.
  \]
  对向量场 $(P,Q,R)=(x,y,z^2-2z)$ 用高斯公式，有
  \[
    \iint_\Sigma P\,dy\,dz+Q\,dz\,dx+R\,dx\,dy
    +\iint_{\Sigma_1} P\,dy\,dz+Q\,dz\,dx+R\,dx\,dy
    =
    \iiint_\Omega 2z\,dx\,dy\,dz.
  \]
  右端为
  \[
    \int_0^1 2z\cdot \pi z^2\,dz=\frac{\pi}{2}.
  \]
  在 $\Sigma_1$ 上 $z=1$ 且取上侧，故
  \[
    \iint_{\Sigma_1} P\,dy\,dz+Q\,dz\,dx+R\,dx\,dy
    =
    \iint_{x^2+y^2\le1}(-1)\,dx\,dy=-\pi.
  \]
  因此所求积分为
  \[
    \frac{\pi}{2}-(-\pi)=\frac{3\pi}{2}.
  \]

  \item 幂级数的收敛半径为 $1$。当 $x=1$ 时，级数 $\sum 1/n$ 发散；当 $x=-1$ 时，级数
  $\sum (-1)^n/n$ 收敛。因此收敛域为 $[-1,1)$。

  记
  \[
    f(x)=\sum_{n=1}^{\infty}\frac{x^n}{n},\qquad x\in[-1,1).
  \]
  当 $x\in(-1,1)$ 时，
  \[
    f'(x)=\sum_{n=1}^{\infty}x^{n-1}=\frac1{1-x}.
  \]
  由 $f(0)=0$ 得
  \[
    f(x)=\int_0^x\frac{dt}{1-t}=-\ln(1-x),\qquad -1<x<1.
  \]
  在端点 $x=-1$ 由连续性延拓，故
  \[
    f(x)=-\ln(1-x),\qquad x\in[-1,1).
  \]
\end{enumerate}

\paragraph*{五、证明题}
\begin{enumerate}
  \item 设 $(x_0,y_0,z_0)$ 为曲面上一点，则
  \[
    z_0=x_0e^{y_0/x_0}.
  \]
  曲面的一个法向量为
  \[
    \left(-z_x,-z_y,1\right)\bigg|_{(x_0,y_0)}
    =
    \left(
      -\left(1-\frac{y_0}{x_0}\right)e^{y_0/x_0},
      -e^{y_0/x_0},
      1
    \right).
  \]
  于是切平面方程为
  \[
    -\left(1-\frac{y_0}{x_0}\right)e^{y_0/x_0}(x-x_0)
    -e^{y_0/x_0}(y-y_0)
    +(z-z_0)=0.
  \]
  将点 $(0,0,0)$ 代入左端，得到
  \[
    \left(1-\frac{y_0}{x_0}\right)x_0e^{y_0/x_0}
    +y_0e^{y_0/x_0}-z_0
    =
    x_0e^{y_0/x_0}-z_0=0.
  \]
  因此所有切平面都过定点 $(0,0,0)$。

  \item 对任意 $x\in(0,1)$，几何级数
  \[
    \sum_{n=0}^{\infty}x^n
  \]
  收敛，且和函数为
  \[
    S(x)=\frac1{1-x}.
  \]
  其部分和
  \[
    S_n(x)=\sum_{k=0}^{n}x^k=\frac{1-x^{n+1}}{1-x}.
  \]
  取 $x_n=1-\dfrac1n\in(0,1)$，则
  \[
    |S_n(x_n)-S(x_n)|
    =
    \frac{x_n^{n+1}}{1-x_n}
    =
    n\left(1-\frac1n\right)^{n+1}\not\to0.
  \]
  因此余项不能在 $(0,1)$ 上一致趋于 $0$，原函数项级数在 $(0,1)$ 上不一致收敛。
\end{enumerate}

\paragraph*{六、应用题}
设直角三角形两条直角边长分别为 $x,y$，则
\[
  x^2+y^2=l^2,\qquad x>0,\quad y>0,
\]
周长为
\[
  P=x+y+l.
\]
作拉格朗日函数
\[
  L(x,y,\lambda)=x+y+l-\lambda(x^2+y^2-l^2).
\]
驻点满足
\[
  1-2\lambda x=0,\qquad 1-2\lambda y=0,\qquad x^2+y^2=l^2.
\]
由前两式得 $x=y$，再由约束得
\[
  x=y=\frac{l}{\sqrt2}.
\]
此时周长最大，最大周长为
\[
  P_{\max}=(1+\sqrt2)l.
\]

\section{2017级工科数学分析下 A卷}

\subsection{试题}

\paragraph*{一、填空题：共 5 题，每题 2 分，共 10 分}
\begin{enumerate}
  \item 微分方程 $y''+y'-2y=1-2x$ 的通解为 \underline{\hspace{5cm}}。

  \item 设函数 $u=\ln(x^2+y^2+z^2)$，求 $\operatorname{div}(\operatorname{grad}u)=$ \underline{\hspace{4cm}}。

  \item 设 $\Gamma$ 是球面 $x^2+y^2+z^2=R^2$ 与平面 $x+y+z=0$ 的交线，则第一类曲线积分
  \[
    \oint_{\Gamma} y^2\,ds=
  \]
  \underline{\hspace{4cm}}。

  \item 参数曲线
  \[
    \begin{cases}
      x=t-\cos t,\\
      y=1-\sin t,\\
      z=t,
    \end{cases}
  \]
  在 $t=0$ 对应的点处的切线方程为 \underline{\hspace{6cm}}。

  \item 设周期为 $2\pi$ 的函数
  \[
    f(x)=
    \begin{cases}
      -1, & -\pi<x\le 0,\\
      1, & 0<x\le \pi,
    \end{cases}
  \]
  则 $f(x)$ 的傅里叶（Fourier）级数在 $x=\pi$ 处收敛于 \underline{\hspace{4cm}}。
\end{enumerate}

\paragraph*{二、单选题：共 5 题，每题 2 分，共 10 分}
\begin{enumerate}
  \item 关于未知函数 $y$ 的微分方程
  \[
    (y-\sin x)\,dx+\ln x\,dy=0
  \]
  是（\quad）。
  \[
  \begin{array}{ll}
    \text{A. 可分离变量方程；} & \text{B. 一阶线性非齐次方程；}\\
    \text{C. 一阶线性齐次方程；} & \text{D. 非线性方程。}
  \end{array}
  \]

  \item 二元函数
  \[
    f(x,y)=
    \begin{cases}
      x\sin\dfrac{1}{y}, & xy\ne 0,\\
      0, & xy=0,
    \end{cases}
  \]
  则 $\displaystyle \lim_{(x,y)\to(0,0)}f(x,y)$（\quad）。
  \[
  \begin{array}{llll}
    \text{A. 不存在；} & \text{B. 等于 1；} &
    \text{C. 等于 0；} & \text{D. 等于 2。}
  \end{array}
  \]

  \item 设函数 $z=f(x,y)$ 在 $(a,b)$ 的某个邻域内有直到二阶的连续偏导数，且
  \[
    \frac{\partial z}{\partial x}(a,b)=0,\qquad
    \frac{\partial z}{\partial y}(a,b)=0.
  \]
  记
  \[
    A=\frac{\partial^2 z}{\partial x^2}(a,b),\qquad
    B=\frac{\partial^2 z}{\partial x\partial y}(a,b),\qquad
    C=\frac{\partial^2 z}{\partial y^2}(a,b).
  \]
  则函数 $z=f(x,y)$ 在点 $(a,b)$ 取极大值的充分条件是（\quad）。
  \[
  \begin{array}{ll}
    \text{A. } A>0,\ AC>B^2; & \text{B. } A<0,\ AC>B^2;\\
    \text{C. } A>0,\ AC<B^2; & \text{D. } A<0,\ AC<B^2.
  \end{array}
  \]

  \item 函数 $\ln(1-x)$ 在 $x=0$ 处的泰勒（Taylor）展开式正确的是（\quad）。
  \[
  \begin{array}{ll}
    \text{A. } \displaystyle \ln(1-x)=\sum_{n=1}^{\infty}\frac{(-1)^n}{n}x^n,\quad x\in(-1,1];&
    \text{B. } \displaystyle \ln(1-x)=\sum_{n=1}^{\infty}\frac{(-1)^{n-1}}{n}x^n,\quad x\in(-1,1];\\[1.2em]
    \text{C. } \displaystyle \ln(1-x)=-\sum_{n=1}^{\infty}\frac{1}{n}x^n,\quad x\in[-1,1);&
    \text{D. } \displaystyle \ln(1-x)=\sum_{n=1}^{\infty}\frac{1}{n}x^n,\quad x\in[-1,1).
  \end{array}
  \]

  \item 使得级数
  \[
    \sum_{n=1}^{\infty}\frac{(-1)^n\ln n}{n^p}
  \]
  条件收敛的常数 $p$ 的取值范围是（\quad）。
  \[
  \begin{array}{llll}
    \text{A. }p\le 0; & \text{B. }0<p<1; &
    \text{C. }0<p\le 1; & \text{D. }p>1.
  \end{array}
  \]
\end{enumerate}

\paragraph*{三、计算题：共 3 题，每题 10 分，共 30 分}
\begin{enumerate}
  \item 设 $u=f(x,y,z)$，其中函数 $f$ 有二阶连续的偏导数，且 $z=z(x,y)$ 由方程
  \[
    z^5-5xy+5z=1
  \]
  所确定，求 $\dfrac{\partial u}{\partial x}$ 和 $\dfrac{\partial^2u}{\partial x^2}$。

  \item 计算累次积分
  \[
    \int_{\frac14}^{\frac12} dx\int_{\frac12}^{\sqrt{x}} e^{x/y}\,dy
    +\int_{\frac12}^{1} dx\int_x^{\sqrt{x}} e^{x/y}\,dy.
  \]

  \item 计算球面 $x^2+y^2+z^2=4z$ 含在抛物面 $z=x^2+y^2$ 内的部分的面积。
\end{enumerate}

\paragraph*{四、解答题：共 3 题，每题 10 分，共 30 分}
\begin{enumerate}
  \item 设 $\Sigma$ 是锥面 $z=\sqrt{x^2+y^2}$ 被平面 $z=0$ 及 $z=1$ 截下的部分的下侧，计算第二类曲面积分
  \[
    \iint_{\Sigma} xye^z\,dy\,dz+yz^2\,dz\,dx-ye^z\,dx\,dy.
  \]

  \item 设曲线积分
  \[
    \int_{\Gamma}\left(\sin x-f(x)\right)\frac{y}{x}\,dx+f(x)\,dy
  \]
  与路径无关，其中 $f(x)$ 有一阶连续导数且 $f(\pi)=0$，求 $f(x)$ 并计算曲线积分
  \[
    \int_{(1,0)}^{(\pi,\pi)}
    \left(\sin x-f(x)\right)\frac{y}{x}\,dx+f(x)\,dy.
  \]

  \item 求幂级数
  \[
    \sum_{n=0}^{\infty}\frac{1}{2n+1}x^{2n+1}
  \]
  的收敛域及和函数。
\end{enumerate}

\paragraph*{五、证明题：共 1 题，每题 10 分，共 10 分}
证明函数项级数
\[
  \sum_{n=0}^{\infty}x^2e^{-nx}
\]
在 $[0,+\infty)$ 一致收敛。

\paragraph*{六、应用题：共 1 题，每题 10 分，共 10 分}
将长度为 $2a$ 的铁丝分成三段，分别围成一个正方形、一个圆形和一个正三角形，求三个图形面积之和的最小值。

\subsection{答案与解析}

\paragraph*{一、填空题}
\begin{enumerate}
  \item
  \[
    y=C_1e^{-2x}+C_2e^x+x.
  \]

  \item
  \[
    \operatorname{div}(\operatorname{grad}u)
    =\Delta u=\frac{2}{x^2+y^2+z^2}.
  \]

  \item
  \[
    \oint_{\Gamma}y^2\,ds=\frac{2}{3}\pi R^3.
  \]

  \item 当 $t=0$ 时，点为 $(-1,1,0)$，切向量为 $(1,-1,1)$，故切线方程为
  \[
    \frac{x+1}{1}=\frac{y-1}{-1}=\frac{z}{1}.
  \]

  \item 傅里叶级数在跳跃点收敛于左右极限的平均值。由周期延拓知
  \[
    \frac{f(\pi-0)+f(\pi+0)}{2}=\frac{1+(-1)}{2}=0.
  \]
\end{enumerate}

\paragraph*{二、单选题}
\begin{enumerate}
  \item 方程可化为 \(y'+\dfrac1{\ln x}y=\dfrac{\sin x}{\ln x}\)，是一阶线性非齐次方程。选 B。
  \item 因 \(|x\sin(1/y)|\le |x|\)，故 $(x,y)\to(0,0)$ 时极限为 $0$。选 C。
  \item 二元函数极大值的充分条件为 \(A<0,\ AC-B^2>0\)。选 B。
  \item 按题面函数应有 \(\displaystyle \ln(1-x)=-\sum_{n=1}^{\infty}\frac{x^n}{n}\)，收敛区间为 $[-1,1)$，故选 C。\textit{（校勘：答案页标为 A，与题面函数及端点范围不符。）}
  \item \(\displaystyle \sum (-1)^n\frac{\ln n}{n^p}\) 在 $p>0$ 时由交错判别法收敛，在 $p>1$ 时绝对收敛，故条件收敛范围为 $0<p\le1$。选 C。
\end{enumerate}

\paragraph*{三、计算题}
\begin{enumerate}
  \item 对方程 $z^5-5xy+5z=1$ 两边求全微分，得
  \[
    5z^4\,dz-5x\,dy-5y\,dx+5\,dz=0.
  \]
  因此
  \[
    dz=\frac{y\,dx+x\,dy}{z^4+1},
  \]
  即
  \[
    \frac{\partial z}{\partial x}=\frac{y}{z^4+1},
    \qquad
    \frac{\partial z}{\partial y}=\frac{x}{z^4+1}.
  \]
  由复合函数求导法则，
  \[
    \frac{\partial u}{\partial x}
    =f'_1(x,y,z)+f'_3(x,y,z)\frac{\partial z}{\partial x}
    =f'_1(x,y,z)+f'_3(x,y,z)\frac{y}{z^4+1}.
  \]
  进一步，
  \[
    \frac{\partial^2u}{\partial x^2}
    =f''_{11}(x,y,z)+2f''_{13}(x,y,z)\frac{y}{z^4+1}
    +f''_{33}(x,y,z)\left(\frac{y}{z^4+1}\right)^2
    -f'_3(x,y,z)\frac{4y^2z^3}{(z^4+1)^3}.
  \]

  \item 记 $D$ 为由 $y=\dfrac12$、$y=\sqrt{x}$ 和 $y=x$ 围成的区域。原积分化为
  \[
    \iint_D e^{x/y}\,dx\,dy
    =\int_{\frac12}^{1}dy\int_{y^2}^{y}e^{x/y}\,dx.
  \]
  于是
  \[
    \int_{\frac12}^{1}dy\int_{y^2}^{y}e^{x/y}\,dx
    =\int_{\frac12}^{1}\left(ey-ye^y\right)\,dy
    =\frac{3}{8}e-\frac12e^{1/2}.
  \]

  \item 球面可写为
  \[
    z=2+\sqrt{4-x^2-y^2}.
  \]
  其含在抛物面 $z=x^2+y^2$ 内的部分在 $xOy$ 平面上的投影区域为
  \[
    D=\{(x,y)\mid x^2+y^2\le 3\}.
  \]
  曲面面积微元为
  \[
    dS=\sqrt{1+\left(\frac{\partial z}{\partial x}\right)^2
    +\left(\frac{\partial z}{\partial y}\right)^2}\,dx\,dy
    =\frac{2}{\sqrt{4-x^2-y^2}}\,dx\,dy.
  \]
  因此所求面积为
  \[
    A=\iint_D\frac{2}{\sqrt{4-x^2-y^2}}\,dx\,dy
    =\int_0^{2\pi}d\theta\int_0^{\sqrt3}\frac{2r}{\sqrt{4-r^2}}\,dr
    =4\pi.
  \]
\end{enumerate}

\paragraph*{四、解答题}
\begin{enumerate}
  \item 考虑平面 $\Sigma_1:z=1$ 在 $x^2+y^2\le 1$ 的部分，取上侧，则 $\Sigma\cup\Sigma_1$ 形成闭曲面，取外侧。记其围成的闭区域为 $\Omega$。由 Gauss 公式，
  \[
    \iint_{\Sigma\cup\Sigma_1} xye^z\,dy\,dz+yz^2\,dz\,dx-ye^z\,dx\,dy
    =\iiint_{\Omega}z^2\,dx\,dy\,dz.
  \]
  而
  \[
    \iiint_{\Omega}z^2\,dx\,dy\,dz
    =\int_0^1 dz\iint_{x^2+y^2\le z^2}z^2\,dx\,dy
    =\int_0^1\pi z^4\,dz=\frac{\pi}{5}.
  \]
  又 $\Sigma_1$ 的单位正法向量为 $\mathbf n=(0,0,1)$，故
  \[
    \iint_{\Sigma_1} xye^z\,dy\,dz+yz^2\,dz\,dx-ye^z\,dx\,dy
    =\iint_{x^2+y^2\le 1}(-ye)\,dx\,dy=0.
  \]
  因此
  \[
    \iint_{\Sigma} xye^z\,dy\,dz+yz^2\,dz\,dx-ye^z\,dx\,dy
    =\frac{\pi}{5}.
  \]

  \item 记
  \[
    P=\left(\sin x-f(x)\right)\frac{y}{x},\qquad Q=f(x).
  \]
  由曲线积分与路径无关，有 $\dfrac{\partial P}{\partial y}=\dfrac{\partial Q}{\partial x}$，即
  \[
    \frac{\sin x}{x}-\frac{f(x)}{x}=f'(x).
  \]
  这是一个一阶线性非齐次微分方程，积分因子为 $x$，从而
  \[
    (xf(x))'=xf'(x)+f(x)=\sin x.
  \]
  其通解为
  \[
    f(x)=-\frac{\cos x}{x}+\frac{C}{x}.
  \]
  由 $f(\pi)=0$ 得 $C=-1$。因此，
  \[
    f(x)=\frac{-\cos x-1}{x}.
  \]
  此时
  \[
    P\,dx+Q\,dy
    =d\left(\frac{(-\cos x-1)y}{x}\right).
  \]
  因此
  \[
    \int_{(1,0)}^{(\pi,\pi)}P\,dx+Q\,dy=0.
  \]

  \item 考虑幂级数
  \[
    \sum_{n=0}^{\infty}\frac{1}{2n+1}t^n.
  \]
  由于
  \[
    \lim_{n\to\infty}
    \frac{\frac{1}{2(n+1)+1}}{\frac{1}{2n+1}}
    =\lim_{n\to\infty}\frac{2n+1}{2n+3}=1,
  \]
  其收敛半径为 $1$。故原幂级数
  \[
    \sum_{n=0}^{\infty}\frac{1}{2n+1}x^{2n+1}
    =x\sum_{n=0}^{\infty}\frac{1}{2n+1}(x^2)^n
  \]
  的收敛半径也为 $1$。当 $x=\pm1$ 时，级数为 $\pm\sum_{n=0}^{\infty}\dfrac{1}{2n+1}$，发散。因此收敛域为 $(-1,1)$。

  设其和函数为
  \[
    f(x)=\sum_{n=0}^{\infty}\frac{1}{2n+1}x^{2n+1},\qquad x\in(-1,1).
  \]
  则
  \[
    f'(x)=\sum_{n=0}^{\infty}x^{2n}=\frac{1}{1-x^2}.
  \]
  因为 $f(0)=0$，所以
  \[
    f(x)=\int_0^x\frac{1}{1-t^2}\,dt
    =\frac12\ln\frac{1+x}{1-x},\qquad x\in(-1,1).
  \]
\end{enumerate}

\paragraph*{五、证明题}
证明一：对任意 $x\in[0,+\infty)$，当 $n\ge 1$ 时，
\[
  e^{nx}>1+nx+\frac12(nx)^2>\frac12(nx)^2.
\]
因此
\[
  |x^2e^{-nx}|=\frac{x^2}{e^{nx}}\le \frac{2}{n^2}.
\]
而数项级数 $\displaystyle \sum_{n=1}^{\infty}\frac{1}{n^2}$ 收敛。由优级数判别法，
\[
  \sum_{n=0}^{\infty}x^2e^{-nx}
\]
在 $[0,+\infty)$ 一致收敛。

证明二：该函数项级数的部分和为
\[
  S_n(x)=
  \begin{cases}
    \dfrac{x^2(1-e^{-(n+1)x})}{1-e^{-x}}, & x>0,\\
    0, & x=0.
  \end{cases}
\]
因此
\[
  S(x)=\lim_{n\to\infty}S_n(x)=
  \begin{cases}
    \dfrac{x^2}{1-e^{-x}}, & x>0,\\
    0, & x=0.
  \end{cases}
\]
下证 $S_n(x)$ 一致收敛于 $S(x)$。对任意 $\varepsilon>0$，取 $N>\dfrac1\varepsilon$，当 $n>N$ 时，对任意 $x\in[0,+\infty)$，
\[
  |S_n(x)-S(x)|
  \le \left|\frac{x^2e^{-(n+1)x}}{1-e^{-x}}\right|
  =\left|\frac{x}{e^x-1}\right|\left|\frac{x}{e^{nx}}\right|
  <\frac1n<\varepsilon.
\]
其中用到不等式 $e^x\ge x+1\ge x$。故该函数项级数在 $[0,+\infty)$ 一致收敛。

\paragraph*{六、应用题}
设正方形的边长为 $x$，圆形的半径为 $y$，正三角形的边长为 $z$，则三个图形的周长之和为
\[
  4x+2\pi y+3z=2a.
\]
它们的面积之和为
\[
  f(x,y,z)=x^2+\pi y^2+\frac{\sqrt3}{4}z^2.
\]
构造 Lagrange 函数
\[
  L(x,y,z,\lambda)
  =x^2+\pi y^2+\frac{\sqrt3}{4}z^2
  -\lambda(4x+2\pi y+3z-2a).
\]
极值点满足
\[
  \begin{cases}
    \dfrac{\partial L}{\partial x}=2x-4\lambda=0,\\
    \dfrac{\partial L}{\partial y}=2\pi y-2\pi\lambda=0,\\
    \dfrac{\partial L}{\partial z}=\dfrac{\sqrt3}{2}z-3\lambda=0,\\
    \dfrac{\partial L}{\partial \lambda}=-(4x+2\pi y+3z-2a)=0.
  \end{cases}
  \]
  解得
  \[
    \lambda=\frac{a}{4+\pi+3\sqrt3},\qquad
    x_0=\frac{2a}{4+\pi+3\sqrt3},\qquad
    y_0=\frac{a}{4+\pi+3\sqrt3},\qquad
    z_0=\frac{2\sqrt3a}{4+\pi+3\sqrt3}.
  \]
  所求最小面积为
  \[
    f(x_0,y_0,z_0)=\frac{a^2}{4+\pi+3\sqrt3}.
  \]

\section{2017级工科数学分析下 B卷}

\subsection{试题}

\paragraph*{一、填空题：共 5 题，每题 2 分，共 10 分。}
\begin{enumerate}
  \item 微分方程 $y''+2y'-3y=2-3x$ 的通解为\underline{\hspace{4cm}}。
  \item 向量场 $2xy\vec{i}+e^z\sin y\vec{j}+(x^2+y^2+z^2)\vec{k}$ 的旋度为\underline{\hspace{4cm}}。
  \item 设 $\Gamma$ 是平面上的下半圆周
  \[
    y=-\sqrt{1-x^2},\qquad -1\le x\le 1,
  \]
  则第一类曲线积分
  \[
    \int_\Gamma (x^2+y^2)\,ds=\underline{\hspace{4cm}}.
  \]
  \item 幂级数 $\displaystyle \sum_{n=1}^{\infty}(-1)^n\dfrac{1}{n\cdot 2^n}x^n$ 的收敛域为\underline{\hspace{4cm}}。
  \item 设周期为 $2\pi$ 的函数 $f(x)=|x|,\ -\pi<x\le\pi$，则 $f(x)$ 的傅里叶（Fourier）级数在 $x=\pi$ 处收敛于\underline{\hspace{4cm}}。
\end{enumerate}

\paragraph*{二、选择题：共 5 题，每题 2 分，共 10 分。}
\begin{enumerate}
  \item 下列微分方程中，属于二阶线性常微分方程的是（\quad）
  \begin{enumerate}
    \item[A.] $(x+y^2)\,dy+e^x\,dx=0$；
    \item[B.] $xy''+(\sin x)y'+(\ln x)y=\tan x$；
    \item[C.] $(y'')^2+y=2$；
    \item[D.] $y''+\ln y=2$。
  \end{enumerate}

  \item 二元函数 $f(x,y)$ 在点 $(a,b)$ 的两个偏导数
  \[
    \frac{\partial f}{\partial x}(a,b),\qquad
    \frac{\partial f}{\partial y}(a,b)
  \]
  存在是 $f(x,y)$ 在该点连续的（\quad）
  \begin{enumerate}
    \item[A.] 充分而非必要条件；
    \item[B.] 必要而非充分条件；
    \item[C.] 充分必要条件；
    \item[D.] 既非充分也非必要条件。
  \end{enumerate}

  \item 曲面 $3x^2+y^2+z^2=12$ 上点 $M(-1,0,3)$ 处的切平面与平面 $z=0$ 的夹角是（\quad）
  \begin{enumerate}
    \item[A.] $\dfrac{\pi}{6}$；
    \item[B.] $\dfrac{\pi}{4}$；
    \item[C.] $\dfrac{\pi}{3}$；
    \item[D.] $\dfrac{\pi}{2}$。
  \end{enumerate}

  \item 设 $\Gamma$ 是平面曲线
  \[
    \frac{x^2}{a^2}+\frac{y^2}{b^2}=1,\qquad a,b>0,
  \]
  取逆时针方向，则第二类曲线积分
  \[
    \int_\Gamma x\,dy-y\,dx
  \]
  等于（\quad）
  \begin{enumerate}
    \item[A.] $\pi ab$；
    \item[B.] $2\pi ab$；
    \item[C.] $0$；
    \item[D.] $-2\pi ab$。
  \end{enumerate}

  \item 下列级数发散的是（\quad）
  \begin{enumerate}
    \item[A.] $\displaystyle \sum_{n=1}^{\infty}2^n\sin\frac{\pi}{3^n}$；
    \item[B.] $\displaystyle \sum_{n=2}^{\infty}\frac{1}{\ln n}$；
    \item[C.] $\displaystyle \sum_{n=2018}^{\infty}\frac{2^n}{n!}$；
    \item[D.] $\displaystyle \sum_{n=1}^{\infty}\left(1-\frac1n\right)^{n^2}$。
  \end{enumerate}
\end{enumerate}

\paragraph*{三、计算题：共 3 题，每题 10 分，共 30 分。}
\begin{enumerate}
  \item 设
  \[
    z=xf\left(\frac{y}{x}\right)+g(x,xy),
  \]
  其中 $f(t)$ 具有连续的二阶导数，$g(u,v)$ 具有连续的二阶偏导数，计算二阶偏导数
  \[
    \frac{\partial^2 z}{\partial x\partial y}.
  \]

  \item 计算累次积分
  \[
    \int_{\frac14}^{\frac12}dx\int_{\frac12}^{\sqrt{x}}\cos\left(\pi\frac{x}{y}\right)\,dy
    +\int_{\frac12}^{1}dx\int_x^{\sqrt{x}}\cos\left(\pi\frac{x}{y}\right)\,dy.
  \]

  \item 计算曲面 $z=\sqrt{x^2+y^2}$ 夹在两柱面
  \[
    x^2+y^2=x,\qquad x^2+y^2=2x
  \]
  之间的那一部分的面积。
\end{enumerate}

\paragraph*{四、解答题：共 3 题，每题 10 分，共 30 分。}
\begin{enumerate}
  \item 设 $\Sigma$ 是曲线
  \[
    \begin{cases}
      z=\sqrt{y-1},\\
      x=0,
    \end{cases}
    \qquad 1\le y\le 3
  \]
  绕 $y$ 轴旋转一周所成的曲面，其法向量与 $y$ 轴正向夹角恒大于 $\dfrac{\pi}{2}$。求曲面积分
  \[
    \iint_\Sigma (8y+1)x\,dy\,dz+2(1-y^2)\,dz\,dx+(z^2-2z)\,dx\,dy.
  \]

  \item 设曲线积分
  \[
    \int_\Gamma (f(x)-e^x)\sin\frac{y}{2}\,dx-\frac12 f(x)\cos\frac{y}{2}\,dy
  \]
  与路径无关，其中 $f(x)$ 有一阶连续导数且 $f(0)=0$。求 $f(x)$，并计算曲线积分
  \[
    \int_{(0,0)}^{(1,1)}(f(x)-e^x)\sin\frac{y}{2}\,dx-\frac12 f(x)\cos\frac{y}{2}\,dy.
  \]

  \item 将函数 $f(x)=\cos^2x$ 展开成 $x$ 的幂级数。
\end{enumerate}

\paragraph*{五、证明题：共 1 题，每题 10 分，共 10 分。}
设 $0<c<1$，证明函数项级数
\[
  \sum_{n=0}^{\infty}(1-x)^n
\]
在 $[c,1]$ 一致收敛，但在 $(0,1)$ 不一致收敛。

\paragraph*{六、应用题：共 1 题，每题 10 分，共 10 分。}
设椭圆
\[
  x^2+3y^2=12
\]
的内接等腰三角形之底边平行于椭圆的长轴，求这样的三角形的最大面积。

\subsection{答案与解析}

\paragraph*{一、填空题答案}
\begin{enumerate}
  \item
  \[
    y=C_1e^{-3x}+C_2e^x+x,\qquad C_1,C_2\text{ 为任意常数}.
  \]

  \item 由旋度公式
  \[
    \nabla\times\vec{F}
    =
    \begin{vmatrix}
      \vec{i} & \vec{j} & \vec{k}\\
      \partial_x & \partial_y & \partial_z\\
      2xy & e^z\sin y & x^2+y^2+z^2
    \end{vmatrix}
    =(2y-e^z\sin y)\vec{i}-2x\vec{j}-2x\vec{k}.
  \]

  \item 因 $\Gamma$ 为单位圆下半圆周，故 $x^2+y^2=1$，弧长为 $\pi$，所以
  \[
    \int_\Gamma (x^2+y^2)\,ds=\int_\Gamma 1\,ds=\pi.
  \]

  \item 答案页给出的收敛域为
  \[
    (-2,2].
  \]
  若按有意义的幂级数下标 $n=1$ 起算，则半径为 $2$，在 $x=2$ 为交错调和级数收敛，在 $x=-2$ 为调和级数发散，故收敛域为 $(-2,2]$。


  \item 傅里叶级数在 $x=\pi$ 处收敛于左右极限平均值。由周期延拓得
  \[
    \frac{f(\pi-0)+f(\pi+0)}{2}=\frac{\pi+\pi}{2}=\pi.
  \]
\end{enumerate}

\paragraph*{二、选择题}
\begin{enumerate}
  \item 只有 \(xy''+(\sin x)y'+(\ln x)y=\tan x\) 对未知函数 $y$ 及其导数均为一次，是二阶线性常微分方程。选 B。
  \item 偏导数存在不能推出连续，连续也不能推出偏导数存在，因此二者既非充分也非必要。选 D。
  \item 曲面法向量为 \((6x,2y,2z)\)，在 $(-1,0,3)$ 处为 $(-6,0,6)$；它与平面 $z=0$ 的法向量 $(0,0,1)$ 夹角为 $\pi/4$。选 B。
  \item 由 Green 公式，\(\displaystyle \oint_\Gamma x\,dy-y\,dx=2S=2\pi ab\)。选 B。
  \item A、C、D 均收敛；\(\displaystyle \sum_{n=2}^{\infty}\frac1{\ln n}\) 发散。选 B。
\end{enumerate}

\paragraph*{三、计算题答案与解析}
\begin{enumerate}
  \item 先对 $y$ 求偏导：
  \[
    \frac{\partial z}{\partial y}
    =f'\left(\frac{y}{x}\right)+xg_2'(x,xy).
  \]
  再对 $x$ 求偏导，得
  \[
    \frac{\partial^2z}{\partial x\partial y}
    =
    -\frac{y}{x^2}f''\left(\frac{y}{x}\right)
    +g_2'(x,xy)+xg_{21}''(x,xy)+xyg_{22}''(x,xy).
  \]

  \item 记 $D$ 为由 $y=\dfrac12,\ y=\sqrt{x}$ 和 $y=x$ 围成的区域。原积分等于
  \begin{align*}
    &\int_{\frac14}^{\frac12}dx\int_{\frac12}^{\sqrt{x}}\cos\left(\pi\frac{x}{y}\right)\,dy
    +\int_{\frac12}^{1}dx\int_x^{\sqrt{x}}\cos\left(\pi\frac{x}{y}\right)\,dy\\
    &\qquad
    =\iint_D\cos\left(\pi\frac{x}{y}\right)\,dx\,dy
    =\int_{\frac12}^{1}dy\int_{y^2}^{y}\cos\left(\pi\frac{x}{y}\right)\,dx.
  \end{align*}
  因此
  \begin{align*}
    \int_{\frac12}^{1}dy\int_{y^2}^{y}\cos\left(\pi\frac{x}{y}\right)\,dx
    &=-\frac1\pi\int_{\frac12}^{1}y\sin(\pi y)\,dy\\
    &=\frac{1-\pi}{\pi^3}.
  \end{align*}

  \item 曲面的面积微元为
  \begin{align*}
    dS
    &=\sqrt{1+\left(\frac{\partial z}{\partial x}\right)^2+\left(\frac{\partial z}{\partial y}\right)^2}\,dx\,dy\\
    &=\sqrt{1+\left(\frac{x}{\sqrt{x^2+y^2}}\right)^2+\left(\frac{y}{\sqrt{x^2+y^2}}\right)^2}\,dx\,dy\\
    &=\sqrt2\,dx\,dy.
  \end{align*}
  投影区域为 $x\le x^2+y^2\le 2x$，其面积为
  \[
    \pi-\frac14\pi=\frac34\pi.
  \]
  故所求面积为
  \[
    A=\iint_{x\le x^2+y^2\le 2x}\sqrt2\,dx\,dy
    =\sqrt2\left(\pi-\frac14\pi\right)
    =\frac{3\sqrt2}{4}\pi.
  \]
\end{enumerate}

\paragraph*{四、解答题答案与解析}
\begin{enumerate}
  \item 设 $\Sigma_1$ 为平面 $y=3$ 在 $x^2+z^2\le 2$ 的部分，取法向量 $(0,1,0)$。则 $\Sigma$ 与 $\Sigma_1$ 形成闭曲面，外侧为正向，所围区域为 $\Omega$。由 Gauss 公式，
  \begin{align*}
    &\iint_{\Sigma\cup\Sigma_1}(8y+1)x\,dy\,dz+2(1-y^2)\,dz\,dx+(z^2-2z)\,dx\,dy\\
    &\qquad
    =\iiint_\Omega (4y+2z-1)\,dx\,dy\,dz.
  \end{align*}
  由关于 $z$ 的对称性，
  \[
    \iiint_\Omega z\,dx\,dy\,dz=0.
  \]
  因此
  \begin{align*}
    \iiint_\Omega (4y+2z-1)\,dx\,dy\,dz
    &=\iiint_\Omega (4y-1)\,dx\,dy\,dz\\
    &=\int_1^3dy\iint_{x^2+z^2\le y-1}(4y-1)\,dx\,dz\\
    &=\int_1^3(4y-1)\pi(y-1)\,dy
    =\frac{50}{3}\pi.
  \end{align*}
  在 $\Sigma_1$ 上，$Q=2(1-y^2)=-16$，故
  \[
    \iint_{\Sigma_1}(8y+1)x\,dy\,dz+2(1-y^2)\,dz\,dx+(z^2-2z)\,dx\,dy
    =\iint_{x^2+z^2\le2}(-16)\,dx\,dz=-32\pi.
  \]
  所以
  \[
    \iint_\Sigma (8y+1)x\,dy\,dz+2(1-y^2)\,dz\,dx+(z^2-2z)\,dx\,dy
    =\frac{50}{3}\pi+32\pi
    =\frac{146}{3}\pi.
  \]

  \item 记
  \[
    P=(f(x)-e^x)\sin\frac{y}{2},\qquad
    Q=-\frac12f(x)\cos\frac{y}{2}.
  \]
  由于曲线积分与路径无关，故
  \[
    \frac{\partial P}{\partial y}=\frac{\partial Q}{\partial x}.
  \]
  即
  \[
    f(x)-e^x=-f'(x).
  \]
  解一阶线性方程 $f'(x)+f(x)=e^x$，得
  \[
    f(x)=\frac12e^x+Ce^{-x}.
  \]
  由 $f(0)=0$，得 $C=-\dfrac12$，所以
  \[
    f(x)=\frac12e^x-\frac12e^{-x}.
  \]
  此时
  \begin{align*}
    P\,dx+Q\,dy
    &=-\frac12(e^x+e^{-x})\sin\frac{y}{2}\,dx
      -\frac14(e^x-e^{-x})\cos\frac{y}{2}\,dy\\
    &=d\left[-\frac12(e^x-e^{-x})\sin\frac{y}{2}\right].
  \end{align*}
  因此
  \begin{align*}
    \int_{(0,0)}^{(1,1)}P\,dx+Q\,dy
    &=
    \left[-\frac12(e^x-e^{-x})\sin\frac{y}{2}\right]_{(0,0)}^{(1,1)}\\
    &=-\frac12(e-e^{-1})\sin\frac12.
  \end{align*}

  \item 由
  \[
    \cos t=\sum_{n=0}^{\infty}\frac{(-1)^n}{(2n)!}t^{2n},
  \]
  以及
  \[
    \cos^2x=\frac12+\frac12\cos2x,
  \]
  得
  \begin{align*}
    \cos^2x
    &=\frac12+\frac12\sum_{n=0}^{\infty}\frac{(-1)^n}{(2n)!}(2x)^{2n}\\
    &=\frac12+\frac12\sum_{n=0}^{\infty}\frac{(-1)^n4^n}{(2n)!}x^{2n}.
  \end{align*}
  其收敛域为 $(-\infty,+\infty)$。
\end{enumerate}

\paragraph*{五、证明题答案与解析}
证明：对任意 $x\in[c,1]$，有
\[
  (1-x)^n\le (1-c)^n.
\]
由于级数
\[
  \sum_{n=0}^{\infty}(1-c)^n
\]
收敛，由 Weierstrass 优级数判别法，函数项级数
\[
  \sum_{n=0}^{\infty}(1-x)^n
\]
在 $[c,1]$ 上一致收敛。

另一方面，该级数在 $(0,1]$ 上逐点收敛。其前 $n+1$ 项部分和为
\[
  S_n(x)=\sum_{k=0}^{n}(1-x)^k
  =\frac{1-(1-x)^{n+1}}{x},
\]
和函数为
\[
  S(x)=\frac1x.
\]
取 $x_n=\dfrac1n\in(0,1)$，则
\[
  \left|S_n(x_n)-S(x_n)\right|
  =\left(1-\frac1n\right)^n(n-1)\to+\infty.
\]
因此函数项级数
\[
  \sum_{n=0}^{\infty}(1-x)^n
\]
在 $(0,1)$ 不一致收敛。

\paragraph*{六、应用题答案与解析}
设等腰三角形的三个顶点为 $(-x,y)$、$(x,y)$ 和 $(0,2)$，其中 $x>0$。则三角形面积为
\[
  f(x,y)=x(2-y).
\]
只需求函数 $f(x,y)$ 在条件 $x^2+3y^2=12$ 下的条件极值。设 Lagrange 函数
\[
  L(x,y,\lambda)=x(2-y)-\lambda(x^2+3y^2-12).
\]
极值点满足
\[
  \begin{cases}
    \dfrac{\partial L}{\partial x}=2-y-2\lambda x=0,\\
    \dfrac{\partial L}{\partial y}=-x-6\lambda y=0,\\
    \dfrac{\partial L}{\partial \lambda}=-(x^2+3y^2-12)=0.
  \end{cases}
\]
解得
\[
  (x,y,\lambda)=\left(3,-1,\frac12\right),\quad
  \left(-3,-1,-\frac12\right),\quad
  (0,2,0).
\]
因 $x>0$，极大值点为 $(x,y)=(3,-1)$。因此三角形最大面积为
\[
  f(3,-1)=9.
\]

\section{2018级工科数学分析下 A卷}

\subsection{资料来源}
\begin{itemize}
  \item 试题：2018 级工科数学分析（二）A 卷，已导出页面图校正公式。
  \item 答案：A 卷附解答文件；抽取缺页处据题面补写解析。
\end{itemize}

\subsection{试题}

\paragraph*{试卷信息}
诚信应考，考试作弊将带来严重后果！

华南理工大学本科生期末考试《工科数学分析（二）》A卷，2018--2019学年第二学期。

注意事项：开考前请将密封线内各项信息填写清楚；所有答案请直接答在试卷上；考试形式：闭卷；本试卷共5大题，满分100分，考试时间120分钟。评阅教师请在试卷袋上评阅栏签名。

\paragraph*{一、填空题（每小题3分，共15分）}
\begin{enumerate}
  \item 微分方程 \(y''-y=2x\mathrm{e}^x\) 的特解形式为 \(y^*=x(ax+b)\mathrm{e}^x\)。
  \item 设函数 \(u(x,y,z)=\mathrm{e}^{x^2+y^2+z^2}\)，则 \(\operatorname{div}(\operatorname{grad}u)=\) \underline{\hspace{5cm}}。
  \item 设 \(L\) 是任意一条光滑的闭曲线，取逆时针方向，则
  \[
    \oint_L(2xy+\cos x)\,\mathrm{d}x+(x^2+\sin y)\,\mathrm{d}y=
  \]
  \underline{\hspace{5cm}}。
  \item 若级数 \(\displaystyle \sum_{n=1}^{\infty}\frac{(-1)^n(n+1)}{n^p}\) 条件收敛，则常数 \(p\) 的取值范围为 \underline{\hspace{5cm}}。
  \item 设 \(f(x)=\dfrac{x^2}{2}\)，\(0\le x<\pi\)，\(S(x)\) 为 \(f(x)\) 的展成以 \(2\pi\) 为周期的正弦级数的和函数，则
  \(S\left(-\dfrac{\pi}{2}\right)=\) \underline{\hspace{5cm}}。
\end{enumerate}

\paragraph*{二、计算题（每小题8分，共32分）}
\begin{enumerate}
  \item 设 \(z=f\left(x,\dfrac{y^2}{x}\right)\)，其中函数 \(f\) 有二阶连续的偏导数，求
  \(\dfrac{\partial^2 z}{\partial x\partial y}\)。

  \item 计算三重积分
  \[
    \iiint_\Omega \left(x+\frac y2-z\right)^2\,\mathrm{d}x\mathrm{d}y\mathrm{d}z,
  \]
  其中 \(\Omega\) 是曲面 \(x^2+\dfrac{y^2}{4}+z^2=1\) 所围成的区域。

  \item 求密度均匀的曲面 \(\Sigma:z=\sqrt{a^2-x^2-y^2}\) 的质心坐标。

  \item 求微分方程 \(y'\cot x=y\ln y\) 的通解。
\end{enumerate}

\paragraph*{三、解答题（每小题10分，共30分）}
\begin{enumerate}
  \item 设曲线积分
  \[
    \int_L xy^2\,\mathrm{d}x+y\varphi(x)\,\mathrm{d}y
  \]
  与积分路径无关，其中 \(\varphi(x)\) 有一阶连续导数且 \(\varphi(0)=0\)。求 \(\varphi(x)\)，并计算曲线积分
  \[
    \int_{(0,0)}^{(1,1)}xy^2\,\mathrm{d}x+y\varphi(x)\,\mathrm{d}y.
  \]

  \item 计算
  \[
    I=\oiint_\Sigma
    \frac{x\,\mathrm{d}y\mathrm{d}z+y\,\mathrm{d}z\mathrm{d}x+z\,\mathrm{d}x\mathrm{d}y}
    {(a^2x^2+b^2y^2+c^2z^2)^{3/2}},
  \]
  其中 \(a,b,c>0\)，\(\Sigma:x^2+y^2+z^2=1\)，取外侧。

  \item 求幂级数 \(\displaystyle \sum_{n=1}^{\infty}\frac{2n-1}{n!}x^{2n}\) 的收敛域及和函数。
\end{enumerate}

\paragraph*{四、证明题（每小题7分，共14分）}
\begin{enumerate}
  \item 设 \(f(x)\in C[0,a]\)，证明：
  \[
    2\int_0^a f(x)\,\mathrm{d}x\int_x^a f(y)\,\mathrm{d}y
    =\left[\int_0^a f(x)\,\mathrm{d}x\right]^2.
  \]
  \item 证明函数项级数
  \[
    \sum_{n=1}^{\infty}\frac{\arctan(nx)}{n^2}
  \]
  在 \((-\infty,+\infty)\) 上一致收敛。
\end{enumerate}

\paragraph*{五、应用题（本题9分）}
在椭圆 \(x^2+4y^2=4\) 的第一象限部分上求一点，使其到直线 \(2x+3y-6=0\) 的距离最短。

\subsection{答案与解析}

\paragraph*{一、填空题}
\begin{enumerate}
  \item 由于 \(1\) 是特征方程 \(r^2-1=0\) 的单根，故特解形式为
  \[
    y^*=x(ax+b)\mathrm{e}^x.
  \]
  \item
  \[
    \operatorname{div}(\operatorname{grad}u)=\Delta u
    =\left[6+4(x^2+y^2+z^2)\right]\mathrm{e}^{x^2+y^2+z^2}.
  \]
  \item 由 Green 公式，\(Q_x-P_y=2x-2x=0\)，故积分为 \(0\)。
  \item 交错收敛要求 \((n+1)/n^p\to0\)，即 \(p>1\)；绝对收敛要求 \(\sum n^{1-p}\) 收敛，即 \(p>2\)。故条件收敛范围为
  \[
    1<p\le2.
  \]
  \item 正弦级数对应奇延拓，故
  \[
    S\left(-\frac{\pi}{2}\right)=-f\left(\frac{\pi}{2}\right)=-\frac{\pi^2}{8}.
  \]
\end{enumerate}

\paragraph*{二、计算题}
\begin{enumerate}
  \item 记 \(u=x,\ v=y^2/x\)，则 \(z=f(u,v)\)。先求
  \[
    z_y=f_2\frac{2y}{x}.
  \]
  再对 \(x\) 求导，得到
  \[
    \frac{\partial^2z}{\partial x\partial y}
    =\frac{2y}{x}f_{12}-\frac{2y^3}{x^3}f_{22}-\frac{2y}{x^2}f_2,
  \]
  其中 \(f_2,f_{12},f_{22}\) 均在 \(\left(x,\dfrac{y^2}{x}\right)\) 处取值。

  \item 展开平方后，由椭球关于三个坐标面对称，交叉项积分均为 \(0\)，所以原积分等于
  \[
    \iiint_\Omega\left(x^2+\frac{y^2}{4}+z^2\right)\,\mathrm{d}V.
  \]
  作变换
  \[
    x=r\sin\varphi\cos\theta,\quad
    y=2r\sin\varphi\sin\theta,\quad
    z=r\cos\varphi,\quad J=2r^2\sin\varphi.
  \]
  因此
  \[
    I=\int_0^{2\pi}\!\!\int_0^\pi\!\!\int_0^1 r^2\cdot2r^2\sin\varphi\,\mathrm{d}r\,\mathrm{d}\varphi\,\mathrm{d}\theta
    =\frac{8\pi}{5}.
  \]

  \item 由对称性，\(\bar x=\bar y=0\)。上半球面面积为 \(2\pi a^2\)，且
  \[
    \iint_\Sigma z\,\mathrm{d}S
    =\int_0^{2\pi}\!\!\int_0^{\pi/2}a\cos\varphi\cdot a^2\sin\varphi\,\mathrm{d}\varphi\,\mathrm{d}\theta
    =\pi a^3.
  \]
  故质心为
  \[
    \left(0,0,\frac a2\right).
  \]

  \item 当 \(y>0\) 时，方程化为
  \[
    \frac{y'}{y\ln y}=\tan x.
  \]
  令 \(u=\ln y\)，则
  \[
    \frac{u'}u=\tan x,\qquad \ln|u|=-\ln|\cos x|+C.
  \]
  因此通解为
  \[
    \ln y=C\sec x,\qquad y=\mathrm{e}^{C\sec x}.
  \]
  其中 \(C=0\) 给出常值解 \(y=1\)。
\end{enumerate}

\paragraph*{三、解答题}
\begin{enumerate}
  \item 令 \(P=xy^2\)，\(Q=y\varphi(x)\)。路径无关要求
  \[
    P_y=Q_x,\qquad 2xy=y\varphi'(x),
  \]
  故 \(\varphi'(x)=2x\)。由 \(\varphi(0)=0\)，得
  \[
    \varphi(x)=x^2.
  \]
  势函数可取
  \[
    F(x,y)=\frac12x^2y^2,
  \]
  所以
  \[
    \int_{(0,0)}^{(1,1)}xy^2\,\mathrm{d}x+yx^2\,\mathrm{d}y
    =F(1,1)-F(0,0)=\frac12.
  \]

  \item 在单位球面上，外法向量为 \((x,y,z)\)，故
  \[
    x\,\mathrm{d}y\mathrm{d}z+y\,\mathrm{d}z\mathrm{d}x+z\,\mathrm{d}x\mathrm{d}y=\mathrm{d}S.
  \]
  于是
  \[
    I=\iint_{x^2+y^2+z^2=1}
    \frac{\mathrm{d}S}{(a^2x^2+b^2y^2+c^2z^2)^{3/2}}.
  \]
  由线性变换 \(X=ax,\ Y=by,\ Z=cz\) 对单位球面面积元的标准公式可得
  \[
    I=\frac{4\pi}{abc}.
  \]

  \item 令 \(t=x^2\)。因为
  \[
    \sum_{n=1}^{\infty}\frac{nt^n}{n!}=t\mathrm{e}^t,\qquad
    \sum_{n=1}^{\infty}\frac{t^n}{n!}=\mathrm{e}^t-1,
  \]
  所以
  \[
    \sum_{n=1}^{\infty}\frac{2n-1}{n!}x^{2n}
    =2x^2\mathrm{e}^{x^2}-(\mathrm{e}^{x^2}-1)
    =(2x^2-1)\mathrm{e}^{x^2}+1.
  \]
  收敛域为 \((-\infty,+\infty)\)。
\end{enumerate}

\paragraph*{四、证明题}
\begin{enumerate}
  \item 令
  \[
    F(x)=\int_x^a f(y)\,\mathrm{d}y.
  \]
  则 \(F'(x)=-f(x)\)，从而
  \[
    2\int_0^a f(x)\,\mathrm{d}x\int_x^a f(y)\,\mathrm{d}y
    =2\int_0^a f(x)F(x)\,\mathrm{d}x
    =-2\int_0^aF(x)F'(x)\,\mathrm{d}x.
  \]
  因此
  \[
    -2\int_0^aF(x)F'(x)\,\mathrm{d}x
    =F(0)^2-F(a)^2
    =\left[\int_0^a f(x)\,\mathrm{d}x\right]^2.
  \]

  \item 对任意 \(x\in(-\infty,+\infty)\)，
  \[
    \left|\frac{\arctan(nx)}{n^2}\right|\le \frac{\pi}{2n^2}.
  \]
  由于 \(\sum_{n=1}^{\infty}\dfrac{\pi}{2n^2}\) 收敛，由 Weierstrass 判别法，原函数项级数在 \((-\infty,+\infty)\) 上一致收敛。
\end{enumerate}

\paragraph*{五、应用题}
第一象限椭圆上有 \(2x+3y\le6\)，故到直线 \(2x+3y-6=0\) 的距离最短等价于使 \(2x+3y\) 最大。设
\[
  F(x,y)=2x+3y,\qquad x^2+4y^2=4.
\]
由 Lagrange 乘子法，
\[
  2=2\lambda x,\qquad 3=8\lambda y.
\]
解得
\[
  x=\frac85,\qquad y=\frac35.
\]
因此所求点为
\[
  \left(\frac85,\frac35\right).
\]

\section{2018级工科数学分析下 B卷}

\subsection{资料来源}
\begin{itemize}
  \item 试题：2018 级工科数学分析（二）B 卷，已导出页面图校正公式。
  \item 答案：B 卷附解答文件，已导出页面图校正公式。
\end{itemize}

\subsection{试题}

\paragraph*{试卷信息}
诚信应考，考试作弊将带来严重后果！

华南理工大学本科生期末考试《工科数学分析（二）》B卷，2018--2019学年第二学期。

注意事项：开考前请将密封线内各项信息填写清楚；所有答案请直接答在试卷上；考试形式：闭卷；本试卷共5大题，满分100分，考试时间120分钟。评阅教师请在试卷袋上评阅栏签名。

\paragraph*{一、填空题（每小题3分，共15分）}
\begin{enumerate}
  \item 设函数 \(u(x,y,z)=\sqrt{x^2+y^2+z^2}\)，则 \(\operatorname{div}(\operatorname{grad}u)=\) \underline{\hspace{5cm}}。
  \item 交换积分次序，则
  \[
    \int_0^2\mathrm{d}y\int_{y^2}^{4}f(x,y)\,\mathrm{d}x=
  \]
  \underline{\hspace{5cm}}。
  \item 设曲线 \(L:\dfrac{x^2}{a^2}+\dfrac{y^2}{b^2}=1\)，取顺时针方向，则
  \[
    \oint_L(x-y)\,\mathrm{d}x+(x+y)\,\mathrm{d}y=
  \]
  \underline{\hspace{5cm}}。
  \item 若级数 \(\displaystyle \sum_{n=1}^{\infty}\frac{(-1)^n(n+2)}{n^p}\) 条件收敛，则常数 \(p\) 的取值范围为 \underline{\hspace{5cm}}。
  \item 设
  \[
    f(x)=
    \begin{cases}
      -1, & -1<x\le0,\\
      x^2+2, & 0<x\le1,
    \end{cases}
  \]
  设 \(S(x)\) 为 \(f(x)\) 展成以 \(2\) 为周期的傅里叶级数的和函数。
  则 \(S(0)=\) \underline{\hspace{5cm}}。
\end{enumerate}

\paragraph*{二、计算题（每小题9分，共36分）}
\begin{enumerate}
  \item 设 \(u=f(x+y+z,x^2+y^2+z^2)\)，其中函数 \(f\) 有二阶连续的偏导数，求
  \[
    \frac{\partial^2u}{\partial x^2}+
    \frac{\partial^2u}{\partial y^2}+
    \frac{\partial^2u}{\partial z^2}.
  \]

  \item 计算三重积分
  \[
    \iiint_\Omega\left(\frac x2-y-z\right)^2\,\mathrm{d}x\mathrm{d}y\mathrm{d}z,
  \]
  其中 \(\Omega\) 是曲面 \(\dfrac{x^2}{4}+y^2+z^2=1\) 所围成的区域。

  \item 求抛物面 \(\Sigma:z=2-(x^2+y^2)\) 的质量（\(z\ge0\) 的部分），其密度函数为
  \(\rho(x,y,z)=x^2+y^2\)。

  \item 求微分方程
  \[
    x(2+y)\,\mathrm{d}x+y(1-x)\,\mathrm{d}y=0
  \]
  的通解。
\end{enumerate}

\paragraph*{三、解答题（每小题10分，共30分）}
\begin{enumerate}
  \item 设曲线积分
  \[
    \int_L\left(\cos x-\varphi(x)\right)\frac yx\,\mathrm{d}x+\varphi(x)\,\mathrm{d}y
  \]
  与积分路径无关，其中 \(\varphi(x)\) 有一阶连续导数且
  \(\varphi\left(\dfrac{\pi}{2}\right)=0\)，求 \(\varphi(x)\)，并计算曲线积分
  \[
    \int_{(1,0)}^{(\pi,\pi)}
    \left(\cos x-\varphi(x)\right)\frac yx\,\mathrm{d}x+\varphi(x)\,\mathrm{d}y.
  \]

  \item 计算
  \[
    I=\oiint_\Sigma
    \frac{x\,\mathrm{d}y\mathrm{d}z+y\,\mathrm{d}z\mathrm{d}x+z\,\mathrm{d}x\mathrm{d}y}
    {(x^2+y^2+z^2)^{3/2}},
  \]
  其中 \(\Sigma\) 为有界单连通区域的闭边界曲面，取外侧。

  \item 求幂级数 \(\displaystyle \sum_{n=0}^{\infty}\frac{3n+1}{n!}x^{3n}\) 的收敛域及和函数。
\end{enumerate}

\paragraph*{四、证明题（本题10分）}
证明函数项级数
\[
  \sum_{n=1}^{\infty}\frac{\cos(nx)}{n^2}
\]
在 \((-\infty,+\infty)\) 上一致收敛。

\paragraph*{五、应用题（本题9分）}
欲造一无盖的长方体容器，已知底部造价为每平方米3元，侧面造价均为每平方米1元，现想用36元造一个容积最大的容器，求它的尺寸。

\subsection{答案与解析}

\paragraph*{一、填空题}
\begin{enumerate}
  \item 当 \((x,y,z)\ne(0,0,0)\) 时，
  \[
    \operatorname{div}(\operatorname{grad}u)=\Delta\sqrt{x^2+y^2+z^2}
    =\frac{2}{\sqrt{x^2+y^2+z^2}}.
  \]
  \item 区域为 \(0\le y\le2,\ y^2\le x\le4\)，换序得
  \[
    \int_0^2\mathrm{d}y\int_{y^2}^{4}f(x,y)\,\mathrm{d}x
    =\int_0^4\mathrm{d}x\int_0^{\sqrt x}f(x,y)\,\mathrm{d}y.
  \]
  \item 由 Green 公式，逆时针时
  \[
    \oint_L(x-y)\,\mathrm{d}x+(x+y)\,\mathrm{d}y
    =\iint_D 2\,\mathrm{d}A=2\pi ab.
  \]
  题中取顺时针方向，故结果为
  \[
    -2\pi ab.
  \]
  \item 交错收敛要求 \((n+2)/n^p\to0\)，即 \(p>1\)；绝对收敛要求 \(\sum n^{1-p}\) 收敛，即 \(p>2\)。故条件收敛范围为
  \[
    1<p\le2.
  \]
  \item 傅里叶级数在跳跃点取左右极限平均值：
  \[
    S(0)=\frac{f(0-0)+f(0+0)}2=\frac{-1+2}{2}=\frac12.
  \]
\end{enumerate}

\paragraph*{二、计算题}
\begin{enumerate}
  \item 记
  \[
    s=x+y+z,\qquad r=x^2+y^2+z^2.
  \]
  有
  \[
    u_x=f_1+2xf_2,
  \]
  从而
  \[
    u_{xx}=f_{11}+4xf_{12}+4x^2f_{22}+2f_2.
  \]
  同理可得 \(u_{yy}\)、\(u_{zz}\)，相加得到
  \[
    u_{xx}+u_{yy}+u_{zz}
    =3f_{11}+4(x+y+z)f_{12}+4(x^2+y^2+z^2)f_{22}+6f_2.
  \]
  各偏导数均在 \((x+y+z,x^2+y^2+z^2)\) 处取值。

  \item 展开平方后由椭球区域对称性，交叉项积分为 \(0\)，所以
  \[
    \iiint_\Omega\left(\frac x2-y-z\right)^2\,\mathrm{d}V
    =\iiint_\Omega\left(\frac{x^2}{4}+y^2+z^2\right)\,\mathrm{d}V.
  \]
  作广义球坐标变换
  \[
    x=2r\sin\varphi\cos\theta,\quad
    y=r\sin\varphi\sin\theta,\quad
    z=r\cos\varphi,\quad J=2r^2\sin\varphi.
  \]
  因此
  \[
    I=2\int_0^{2\pi}\int_0^\pi\int_0^1r^4\sin\varphi\,\mathrm{d}r\,\mathrm{d}\varphi\,\mathrm{d}\theta
    =\frac{8\pi}{5}.
  \]

  \item 投影区域为 \(D:x^2+y^2\le2\)，且
  \[
    \mathrm{d}S=\sqrt{1+4x^2+4y^2}\,\mathrm{d}x\mathrm{d}y.
  \]
  因而质量
  \[
    M=\iint_D(x^2+y^2)\sqrt{1+4x^2+4y^2}\,\mathrm{d}x\mathrm{d}y
    =2\pi\int_0^{\sqrt2}r^3\sqrt{1+4r^2}\,\mathrm{d}r
    =\frac{149\pi}{30}.
  \]

  \item 原方程分离为
  \[
    \frac{y\,\mathrm{d}y}{2+y}=\frac{x\,\mathrm{d}x}{x-1}.
  \]
  两边积分得
  \[
    y-2\ln|y+2|=x+\ln|x-1|+C,
  \]
  等价地可写为
  \[
    (x-1)(y+2)^2=C\mathrm{e}^{y-x}.
  \]
\end{enumerate}

\paragraph*{三、解答题}
\begin{enumerate}
  \item 取右半平面作为单连通区域。令
  \[
    P=\left(\cos x-\varphi(x)\right)\frac yx,\qquad Q=\varphi(x).
  \]
  路径无关要求 \(P_y=Q_x\)，即
  \[
    \varphi'(x)+\frac1x\varphi(x)=\frac{\cos x}{x}.
  \]
  因此
  \[
    (x\varphi(x))'=\cos x,\qquad
    x\varphi(x)=\sin x+C.
  \]
  由 \(\varphi(\pi/2)=0\)，得 \(C=-1\)，故
  \[
    \varphi(x)=\frac{\sin x-1}{x}.
  \]
  势函数可取 \(F(x,y)=y\varphi(x)\)，于是
  \[
    \int_{(1,0)}^{(\pi,\pi)}P\,\mathrm{d}x+Q\,\mathrm{d}y
    =F(\pi,\pi)-F(1,0)=-1.
  \]

  \item 记 \(\Sigma\) 围成的区域为 \(\Omega\)。若原点不在 \(\Omega\) 内，则向量场
  \[
    \mathbf F=\frac{(x,y,z)}{(x^2+y^2+z^2)^{3/2}}
  \]
  在 \(\Omega\) 内无奇点，且散度为 \(0\)，由 Gauss 公式得
  \[
    I=0.
  \]
  若原点在 \(\Omega\) 内，则挖去小球 \(x^2+y^2+z^2=\varepsilon^2\)。外边界和内侧小球面合用 Gauss 公式，得
  \[
    I=\frac1{\varepsilon^3}\oiint_{x^2+y^2+z^2=\varepsilon^2}
    x\,\mathrm{d}y\mathrm{d}z+y\,\mathrm{d}z\mathrm{d}x+z\,\mathrm{d}x\mathrm{d}y
    =4\pi.
  \]

  \item 记
  \[
    s(x)=\sum_{n=0}^{\infty}\frac{3n+1}{n!}x^{3n}.
  \]
  因为
  \[
    \lim_{n\to\infty}\frac{a_{n+1}}{a_n}=0,
  \]
  收敛半径为 \(+\infty\)，收敛域为 \((-\infty,+\infty)\)。又
  \[
    s(x)=3\sum_{n=0}^{\infty}\frac{n(x^3)^n}{n!}
    +\sum_{n=0}^{\infty}\frac{(x^3)^n}{n!}
    =(3x^3+1)\mathrm{e}^{x^3}.
  \]
\end{enumerate}

\paragraph*{四、证明题}
对任意 \(x\in(-\infty,+\infty)\)，有
\[
  \left|\frac{\cos(nx)}{n^2}\right|\le\frac1{n^2}.
\]
又 \(\sum_{n=1}^{\infty}1/n^2\) 收敛，由 Weierstrass 判别法，原函数项级数在 \((-\infty,+\infty)\) 上一致收敛。

\paragraph*{五、应用题}
设长方体的长、宽、高分别为 \(x,y,z\) 米。费用约束为
\[
  3xy+2xz+2yz=36.
\]
最大化 \(V=xyz\)。令
\[
  F=\ln x+\ln y+\ln z+\lambda(3xy+2xz+2yz-36).
\]
由
\[
  \frac1x+3\lambda y+2\lambda z=0,\quad
  \frac1y+3\lambda x+2\lambda z=0,\quad
  \frac1z+2\lambda x+2\lambda y=0
\]
解得唯一临界点
\[
  x=2,\qquad y=2,\qquad z=3.
\]
实际问题中最大体积存在，故长、宽、高分别为 \(2\) 米、\(2\) 米、\(3\) 米。

\section{2019级工科数学分析下 A卷}

\subsection{资料来源}
\begin{itemize}
  \item 试题：2019 级工科数学分析（二）A 卷无答案版，已导出页面图校正公式。
  \item 答案一：A 卷答案 1。
  \item 答案二：A 卷答案 2；两份答案有少量抽取差异，以下以题面和清晰公式为准。
\end{itemize}

\subsection{试题}

\paragraph*{试卷信息}
诚信应考，考试作弊将带来严重后果！

华南理工大学本科生期末考试 2019--2020--2 学期《工科数学分析》A卷。

注意事项：开考前请将密封线内各项信息填写清楚；所有答案请直接答在试卷上；考试形式：闭卷；本试卷共（五）大题，满分100分，考试时间120分钟。评阅教师请在试卷袋上评阅栏签名。

\paragraph*{一、填空题（共5小题，每题3分，共15分）}
\begin{enumerate}
  \item 函数 \(u=(2x^2+y^2)z^2\) 在点 \((0,1,1)\) 沿该点梯度方向的方向导数为 \underline{\hspace{5cm}}。
  \item 曲线 \(x=t+1,\ y=t^2+1,\ z=t^3+1\) 在点 \((1,1,1)\) 的切线方程为 \underline{\hspace{5cm}}。
  \item 函数 \(z=x^3+y^3-3xy\) 的极小值为 \underline{\hspace{5cm}}。
  \item 级数 \(\displaystyle \sum_{n=1}^{\infty}\frac{x^n}{1+x^{2n}}\) 的收敛域为 \underline{\hspace{5cm}}。
  \item 设 \(f(x)=\mathrm{e}^{x^2}\)，\(0\le x<\pi\)，\(S(x)\) 为 \(f(x)\) 的展成以 \(2\pi\) 为周期的正弦级数的和函数，则 \(S(0)=\) \underline{\hspace{5cm}}。
\end{enumerate}

\paragraph*{二、计算题（共5小题，每题9分，共45分）}
\begin{enumerate}
  \item 求微分方程
  \[
    y''+y'-2y=18x\mathrm{e}^x
  \]
  的通解。

  \item 设 \(z=f(x^2+y^2,x^2y^2)\)，其中 \(f\) 为任意阶可微函数，求
  \[
    \frac{\partial^2z}{\partial x^2},\qquad
    \frac{\partial^2z}{\partial y\partial x}.
  \]

  \item 计算
  \[
    \iiint_\Omega (x+y+z)\,\mathrm{d}x\mathrm{d}y\mathrm{d}z,
  \]
  其中 \(\Omega\) 是曲面 \(z=\sqrt{x^2+y^2}\) 与 \(z=\sqrt{1-x^2-y^2}\) 所围成的区域。

  \item 计算第二类曲线积分
  \[
    \oint_L xy^2\,\mathrm{d}y-x^2y\,\mathrm{d}x,
  \]
  其中 \(L:\dfrac{x^2}{a^2}+\dfrac{y^2}{b^2}=1\)，\(a,b>0\)，取逆时针方向。

  \item 计算
  \[
    I=\iint_\Sigma 2x(y-z)\,\mathrm{d}y\mathrm{d}z+(1-y^2)\,\mathrm{d}z\mathrm{d}x+(1+z^2)\,\mathrm{d}x\mathrm{d}y,
  \]
  其中 \(\Sigma\) 为柱面 \(x^2+y^2=1\) 夹在平面 \(z=0\) 和 \(z=3\) 部分的曲面，方向指向外侧。
\end{enumerate}

\paragraph*{三、解答下列各题（共2小题，每题8分，共16分）}
\begin{enumerate}
  \item 试确定 \(\alpha\) 的值，使得函数
  \[
    f(x,y)=
    \begin{cases}
      (x^2+y^2)^\alpha\sin\dfrac1{x^2+y^2}, & (x,y)\ne(0,0),\\
      0, & (x,y)=(0,0)
    \end{cases}
  \]
  在点 \((0,0)\) 处可微。
  \item 求正项级数 \(\displaystyle \sum_{n=0}^{\infty}\frac{2n+1}{n!}\) 的和。
\end{enumerate}

\paragraph*{四、证明题（共2小题，每题8分，共16分）}
\begin{enumerate}
  \item 设 \(x=x(y,z)\)，\(y=y(x,z)\)，\(z=z(x,y)\)，都是由方程 \(F(x,y,z)=0\) 所确定的具有连续偏导数的函数，证明
  \[
    \frac{\partial x}{\partial y}\cdot
    \frac{\partial y}{\partial z}\cdot
    \frac{\partial z}{\partial x}=-1.
  \]
  \item 证明函数项级数
  \[
    \sum_{n=2}^{\infty}\frac{(-1)^n(1-\mathrm{e}^{-nx})}{n^2+\sin x}
  \]
  在 \([0,+\infty)\) 上一致收敛。
\end{enumerate}

\paragraph*{五、应用题（共1题，每题8分，共8分）}
将正数 \(9\) 分成三个正数 \(x,y,z\) 之和，使得 \(u=x^2y^3z^4\) 最大。

\subsection{答案与解析}

\paragraph*{一、填空题}
\begin{enumerate}
  \item
  \[
    \nabla u=(4xz^2,2yz^2,2(2x^2+y^2)z),
  \]
  故 \(\nabla u(0,1,1)=(0,2,2)\)，沿梯度方向的方向导数为
  \[
    2\sqrt2.
  \]
  \item 该点对应 \(t=0\)，切向量为
  \[
    (x'(0),y'(0),z'(0))=(1,0,0).
  \]
  因此切线可写为
  \[
    y=1,\qquad z=1.
  \]
  \item 驻点为 \((0,0)\) 与 \((1,1)\)。在 \((1,1)\) 处 Hessian 矩阵正定，故极小值为
  \[
    1+1-3=-1.
  \]
  \item 当 \(|x|<1\) 时与 \(\sum x^n\) 同敛散；当 \(|x|>1\) 时与 \(\sum x^{-n}\) 同敛散；当 \(x=\pm1\) 时通项不趋于 \(0\)。故收敛域为
  \[
    (-\infty,-1)\cup(-1,1)\cup(1,+\infty).
  \]
  \item 正弦级数对应奇延拓，在端点 \(0\) 处取左右极限平均值：
  \[
    S(0)=0.
  \]
\end{enumerate}

\paragraph*{二、计算题}
\begin{enumerate}
  \item 特征方程为
  \[
    r^2+r-2=0,
  \]
  特征根为 \(r_1=-2,\ r_2=1\)，齐次解为 \(C_1\mathrm{e}^{-2x}+C_2\mathrm{e}^x\)。因 \(1\) 是特征根，设特解
  \[
    y^*=x\mathrm{e}^x(Ax+B).
  \]
  代入得
  \[
    6Ax+(3B+2A)=18x,
  \]
  所以 \(A=3,\ B=-2\)。通解为
  \[
    y=C_1\mathrm{e}^{-2x}+C_2\mathrm{e}^x+x(3x-2)\mathrm{e}^x.
  \]

  \item 记
  \[
    u=x^2+y^2,\qquad v=x^2y^2.
  \]
  先有
  \[
    z_x=2xf_1+2xy^2f_2.
  \]
  继续求导得
  \[
    z_{xx}=2f_1+2y^2f_2+4x^2f_{11}+8x^2y^2f_{12}+4x^2y^4f_{22},
  \]
  \[
    z_{yx}=4xyf_{11}+4xy(x^2+y^2)f_{12}+4xyf_2+4x^3y^3f_{22}.
  \]
  其中各偏导数均在 \((x^2+y^2,x^2y^2)\) 处取值。

  \item 用球坐标
  \[
    x=\rho\sin\varphi\cos\theta,\quad
    y=\rho\sin\varphi\sin\theta,\quad
    z=\rho\cos\varphi.
  \]
  两个边界分别为 \(\varphi=\pi/4\) 与 \(\rho=1\)，故
  \[
    0\le\theta\le2\pi,\qquad 0\le\varphi\le\frac{\pi}{4},\qquad 0\le\rho\le1.
  \]
  由于 \(x,y\) 关于区域对称积分为 \(0\)，只需积 \(z\)：
  \[
    I=2\pi\int_0^{\pi/4}\int_0^1 \rho\cos\varphi\cdot\rho^2\sin\varphi\,\mathrm{d}\rho\,\mathrm{d}\varphi
    =\frac{\pi}{8}.
  \]

  \item 令
  \[
    x=a\cos t,\qquad y=b\sin t,\qquad 0\le t\le2\pi.
  \]
  代入得
  \[
    \oint_L xy^2\,\mathrm{d}y-x^2y\,\mathrm{d}x
    =\frac{\pi}{4}\left(ab^3+a^3b\right).
  \]
  也可由 Green 公式化为 \(\iint_D(x^2+y^2)\,\mathrm{d}A\) 得到同一结果。

  \item 添加底面 \(\Sigma_1:z=0\) 和顶面 \(\Sigma_2:z=3\)，构成闭曲面。令
  \[
    P=2x(y-z),\quad Q=1-y^2,\quad R=1+z^2.
  \]
  因为
  \[
    P_x+Q_y+R_z=2(y-z)-2y+2z=0,
  \]
  闭曲面总通量为 \(0\)。底面外法向向下，通量为 \(-\pi\)；顶面外法向向上，通量为 \(10\pi\)。所以侧面通量为
  \[
    I=0-(-\pi)-10\pi=-9\pi.
  \]
\end{enumerate}

\paragraph*{三、解答题}
\begin{enumerate}
  \item 由
  \[
    |f(x,y)|\le (x^2+y^2)^\alpha=r^{2\alpha}
  \]
  可知若 \(\alpha>\dfrac12\)，则
  \[
    \frac{|f(x,y)-f(0,0)|}{\sqrt{x^2+y^2}}\le r^{2\alpha-1}\to0.
  \]
  此时 \(f_x(0,0)=f_y(0,0)=0\)，且函数在原点可微。若 \(\alpha\le\dfrac12\)，上述商不趋于 \(0\)，不可微。因此
  \[
    \alpha>\frac12.
  \]

  \item
  \[
    \sum_{n=0}^{\infty}\frac{2n+1}{n!}
    =2\sum_{n=0}^{\infty}\frac{n}{n!}+\sum_{n=0}^{\infty}\frac1{n!}
    =2\mathrm{e}+\mathrm{e}=3\mathrm{e}.
  \]
\end{enumerate}

\paragraph*{四、证明题}
\begin{enumerate}
  \item 由隐函数求导公式，
  \[
    \frac{\partial x}{\partial y}=-\frac{F_y}{F_x},\qquad
    \frac{\partial y}{\partial z}=-\frac{F_z}{F_y},\qquad
    \frac{\partial z}{\partial x}=-\frac{F_x}{F_z}.
  \]
  三式相乘即得
  \[
    \frac{\partial x}{\partial y}\cdot
    \frac{\partial y}{\partial z}\cdot
    \frac{\partial z}{\partial x}=-1.
  \]

  \item 对任意 \(x\in[0,+\infty)\) 且 \(n\ge2\)，有
  \[
    \left|\frac{(-1)^n(1-\mathrm{e}^{-nx})}{n^2+\sin x}\right|
    \le\frac1{n^2-1}.
  \]
  而数项级数 \(\sum_{n=2}^{\infty}\dfrac1{n^2-1}\) 收敛，由 Weierstrass 判别法，原函数项级数在 \([0,+\infty)\) 上一致收敛。
\end{enumerate}

\paragraph*{五、应用题}
在约束 \(x+y+z=9\) 下最大化 \(u=x^2y^3z^4\)。等价于最大化
\[
  \ln u=2\ln x+3\ln y+4\ln z.
  \]
由 Lagrange 乘子法，
\[
  \frac2x=\lambda,\qquad \frac3y=\lambda,\qquad \frac4z=\lambda.
  \]
故
\[
  x:y:z=2:3:4.
  \]
又 \(x+y+z=9\)，所以
\[
  x=2,\qquad y=3,\qquad z=4.
  \]
最大值为
\[
  2^2\cdot3^3\cdot4^4=27648.
\]

\section{2019级工科数学分析下 B卷补考}

\subsection{资料来源}
\begin{itemize}
  \item 试题：2019 级工科数学分析（二）B 卷补考用，已导出页面图校正公式。
  \item 未发现配套答案或解析文本；下列答案与解析为据题面补写。
\end{itemize}

\subsection{试题}

\paragraph*{试卷信息}
诚信应考，考试作弊将带来严重后果！

华南理工大学本科生期末考试 2019--2020--2 学期《工科数学分析（二）》B卷。

注意事项：开考前请将密封线内各项信息填写清楚；所有答案请直接答在试卷上；考试形式：闭卷；本试卷共（五）大题，满分100分，考试时间120分钟。评阅教师请在试卷袋上评阅栏签名。

\paragraph*{一、填空题（共5小题，每题3分，共15分）}
\begin{enumerate}
  \item 函数 \(z=x^2+2y^2\) 在点 \((1,1)\) 沿该点梯度方向的方向导数为 \underline{\hspace{5cm}}。
  \item 曲线 \(x=t,\ y=t^2,\ z=t^3\) 在点 \((1,1,1)\) 的法平面方程为 \underline{\hspace{5cm}}。
  \item 函数 \(z=xy\) 在条件 \(x^2+y^2=1\) 下的极大值为 \underline{\hspace{5cm}}。
  \item 级数 \(\displaystyle \sum_{n=1}^{\infty}\frac{x^n}{n\cdot4^n}\) 的和函数为 \underline{\hspace{5cm}}。
  \item 设 \(S(x)\) 为 \(f(x)=\mathrm{e}^{\cos x}\)，\(x\in[0,\pi]\)，展成的以 \(2\pi\) 为周期的余弦级数的和函数，则 \(S(0)=\) \underline{\hspace{5cm}}。
\end{enumerate}

\paragraph*{二、计算题（共5小题，每题9分，共45分）}
\begin{enumerate}
  \item 求微分方程
  \[
    \frac{\mathrm{d}y}{\mathrm{d}x}=\frac{2x-y+2}{x+2y-4}
  \]
  的通解。

  \item 假设隐函数存在定理的条件都满足，且 \(y(t)\) 由方程组
  \[
    \begin{cases}
      y=f(x,t),\\
      x=g(y,t)
    \end{cases}
  \]
  确定，求 \(\dfrac{\mathrm{d}y}{\mathrm{d}t}\)。

  \item 计算
  \[
    \iiint_\Omega (x^2+y^2+z^2)\,\mathrm{d}x\mathrm{d}y\mathrm{d}z,
  \]
  其中 \(\Omega\) 是曲面 \(z=x^2+y^2\) 与 \(z=1-x^2-y^2\) 所围成的区域。

  \item 计算曲面积分
  \[
    \iint_\Sigma z\,\mathrm{d}S,
  \]
  其中 \(\Sigma\) 是球面 \(x^2+y^2+z^2=4\) 被平面 \(z=\sqrt2\) 截出的顶部。

  \item 计算
  \[
    I=\iint_{\Sigma^+}\frac{\cos(\vec r,\vec n)}{r^2}\,\mathrm{d}S,
  \]
  其中 \(\Sigma^+\) 为不经过 \(P(a,b,c)\) 的光滑闭曲面，方向向外，\(\Sigma\) 不包含点 \(P(a,b,c)\)。
  \(\vec n\) 为曲面 \(\Sigma\) 上任一点 \(M(x,y,z)\) 处的外法向量的单位向量，\(\cos(\vec r,\vec n)\) 表示两向量 \(\vec r,\vec n\) 夹角的余弦，且
  \[
    \vec r=\overrightarrow{MP}=(a-x,b-y,c-z),\qquad
    r=|\vec r|=\sqrt{(a-x)^2+(b-y)^2+(c-z)^2}.
  \]
\end{enumerate}

\paragraph*{三、解答下列各题（共2小题，每题8分，共16分）}
\begin{enumerate}
  \item 试确定 \(\alpha\) 的值，使得函数
  \[
    f(x,y)=
    \begin{cases}
      \dfrac{x^\alpha y^\alpha}{x^2+y^2}, & (x,y)\ne(0,0),\\
      0, & (x,y)=(0,0)
    \end{cases}
  \]
  在点 \((0,0)\) 处可微。
  \item 将 \(\ln(1-3x+2x^2)\) 展成麦克劳林级数。
\end{enumerate}

\paragraph*{四、证明题（共2小题，每题8分，共16分）}
\begin{enumerate}
  \item 设 \(f\) 为连续函数，\(C\) 为分段光滑的简单闭曲线，证明
  \[
    \oint_C f(x^2-y^2)(x\,\mathrm{d}x-y\,\mathrm{d}y)=0.
  \]
  \item 证明函数项级数
  \[
    \sum_{n=1}^{\infty}\frac{(-1)^n(1-\mathrm{e}^{-nx})}{n^2+x^2}
  \]
  在 \([0,+\infty)\) 上一致收敛。
\end{enumerate}

\paragraph*{五、应用题（共1题，每题8分，共8分）}
求点 \(\left(2,2,\dfrac12\right)\) 到抛物面 \(z=x^2+y^2\) 的最小距离。

\subsection{答案与解析}
\paragraph*{一、填空题}
\begin{enumerate}
  \item \(\nabla z(1,1)=(2,4)\)，沿梯度方向的方向导数为
  \[
    \|\nabla z(1,1)\|=2\sqrt5.
  \]
  \item 曲线切向量为 \((1,2t,3t^2)\)，在 \(t=1\) 处为 \((1,2,3)\)。法平面为
  \[
    (x-1)+2(y-1)+3(z-1)=0,
  \]
  即 \(x+2y+3z-6=0\)。
  \item 由 \(2xy\le x^2+y^2=1\)，得极大值为
  \[
    \frac12.
  \]
  \item
  \[
    \sum_{n=1}^{\infty}\frac{x^n}{n4^n}
    =\sum_{n=1}^{\infty}\frac{(x/4)^n}{n}
    =-\ln\left(1-\frac x4\right),\qquad -4\le x<4.
  \]
  \item 余弦级数在端点取偶延拓的函数值，故
  \[
    S(0)=f(0)=\mathrm{e}.
  \]
\end{enumerate}

\paragraph*{二、计算题}
\begin{enumerate}
  \item 作平移
  \[
    X=x,\qquad Y=y-2.
  \]
  方程化为
  \[
    \frac{\mathrm{d}Y}{\mathrm{d}X}=\frac{2X-Y}{X+2Y}.
  \]
  写成微分形式：
  \[
    (2X-Y)\,\mathrm{d}X-(X+2Y)\,\mathrm{d}Y=0.
  \]
  这是恰当方程，积分得
  \[
    X^2-XY-Y^2=C.
  \]
  因而原方程通解为
  \[
    x^2-x(y-2)-(y-2)^2=C.
  \]

  \item 对
  \[
    y=f(x,t),\qquad x=g(y,t)
  \]
  关于 \(t\) 求导，得
  \[
    y'=f_xx'+f_t,\qquad x'=g_yy'+g_t.
  \]
  消去 \(x'\)，得到
  \[
    \frac{\mathrm{d}y}{\mathrm{d}t}
    =\frac{f_t+f_xg_t}{1-f_xg_y}.
  \]

  \item 用柱坐标。区域为
  \[
    0\le r\le \frac1{\sqrt2},\qquad r^2\le z\le 1-r^2,\qquad 0\le\theta\le2\pi.
  \]
  所以
  \[
    I=2\pi\int_0^{1/\sqrt2}\int_{r^2}^{1-r^2}(r^2+z^2)r\,\mathrm{d}z\,\mathrm{d}r
    =\frac{11\pi}{96}.
  \]

  \item 在球面 \(x^2+y^2+z^2=4\) 上取球坐标 \(z=2\cos\varphi\)，顶部对应
  \(0\le\varphi\le\pi/4\)。于是
  \[
    \iint_\Sigma z\,\mathrm{d}S
    =\int_0^{2\pi}\int_0^{\pi/4}2\cos\varphi\cdot4\sin\varphi\,\mathrm{d}\varphi\,\mathrm{d}\theta
    =4\pi.
  \]

  \item 被积函数可看作向量场
  \[
    \mathbf F(M)=\frac{\overrightarrow{MP}}{|\overrightarrow{MP}|^3}
  \]
  通过 \(\Sigma\) 的通量。若点 \(P\) 在 \(\Sigma\) 所围区域外，则 \(\mathbf F\) 在该区域内无奇点且散度为 \(0\)，由 Gauss 公式可得
  \[
    I=0.
  \]
  若点 \(P\) 在 \(\Sigma\) 所围区域内，由于 \(\vec r=\overrightarrow{MP}=P-M\) 与标准径向场方向相反，通量为
  \[
    I=-4\pi.
  \]
\end{enumerate}

\paragraph*{三、解答题}
\begin{enumerate}
  \item 由极坐标估计，
  \[
    |f(x,y)|\le C r^{2\alpha-2}.
  \]
  若在原点可微且导数为零，需要
  \[
    \frac{|f(x,y)|}{r}\le C r^{2\alpha-3}\to0.
  \]
  因此
  \[
    \alpha>\frac32.
  \]
  此时 \(f_x(0,0)=f_y(0,0)=0\)，且 \(f(x,y)=o(\sqrt{x^2+y^2})\)，故函数在原点可微。

  \item 因为
  \[
    1-3x+2x^2=(1-x)(1-2x),
  \]
  所以
  \[
    \ln(1-3x+2x^2)=\ln(1-x)+\ln(1-2x)
    =-\sum_{n=1}^{\infty}\frac{1+2^n}{n}x^n,\qquad |x|<\frac12.
  \]
\end{enumerate}

\paragraph*{四、证明题}
\begin{enumerate}
  \item 取原函数 \(F'(u)=f(u)\)。由于
  \[
    \mathrm{d}(x^2-y^2)=2x\,\mathrm{d}x-2y\,\mathrm{d}y,
  \]
  故
  \[
    f(x^2-y^2)(x\,\mathrm{d}x-y\,\mathrm{d}y)
    =\frac12\,\mathrm{d}F(x^2-y^2).
  \]
  闭曲线上的全微分积分为 \(0\)，因此结论成立。

  \item 对 \(x\in[0,+\infty)\)，有
  \[
    \left|\frac{(-1)^n(1-\mathrm{e}^{-nx})}{n^2+x^2}\right|
    \le\frac1{n^2}.
  \]
  因为 \(\sum_{n=1}^{\infty}1/n^2\) 收敛，由 Weierstrass 判别法，原函数项级数在 \([0,+\infty)\) 上一致收敛。
\end{enumerate}

\paragraph*{五、应用题}
设抛物面上点为 \((x,y,x^2+y^2)\)。距离平方为
\[
  D=(x-2)^2+(y-2)^2+\left(x^2+y^2-\frac12\right)^2.
\]
由对称性，最近点满足 \(x=y=s\)。于是
\[
  D(s)=2(s-2)^2+\left(2s^2-\frac12\right)^2.
\]
令
\[
  D'(s)=8(2s^3-1)=0,
\]
得 \(s=2^{-1/3}\)。最小距离平方为
\[
  D_{\min}^2=\frac{33}{4}-3\sqrt[3]{4}.
\]
因此最小距离为
\[
  \sqrt{\frac{33}{4}-3\sqrt[3]{4}}.
\]

\section{2020级工科数学分析（二）期末考试 A 卷}

\subsection{试题}

\paragraph*{一、计算题（共 5 小题，每小题 8 分，共 40 分）}
\begin{enumerate}
  \item 设 $z=f(xy^2,x^2y)$，求 $\dfrac{\partial^2 z}{\partial y^2}$。
  \item 设
  \[
    f(x)=
    \begin{cases}
      -x, & |x|\le \dfrac{\pi}{2},\\
      x, & \dfrac{\pi}{2}<|x|\le \pi,
    \end{cases}
  \]
  写出 $f(x)$ 的以 $2\pi$ 为周期的 Fourier 级数的和函数 $S(x)$ 在 $[-\pi,\pi]$ 上的表达式。
  \item 计算二次积分
  \[
    \int_0^2 dx\int_{x/2}^1 e^{y^2}\,dy.
  \]
  \item 计算曲线积分
  \[
    I=\oint_L\frac{-y\,dx+x\,dy}{x^2+y^2},
  \]
  其中 $L:\dfrac{x^2}{a^2}+\dfrac{y^2}{b^2}=1$，取逆时针方向。
  \item 计算曲面积分
  \[
    I=\iint_\Sigma (4xz+x)\,dy\,dz-2yz\,dz\,dx+(1-z^2)\,dx\,dy,
  \]
  其中 $\Sigma:z=x^2+y^2$，$0\le z\le 1$，取下侧。
\end{enumerate}

\paragraph*{二、解答下列各题（共 5 小题，每小题 8 分，共 40 分）}
\begin{enumerate}
  \item 求曲线
  \[
    \Gamma:
    \begin{cases}
      x^2+y^2+z^2-3x=0,\\
      x+y-z=0
    \end{cases}
  \]
  在点 $(1,1,1)$ 处的单位切向量 $\vec t$，并求函数
  \[
    u(x,y,z)=\ln(x^2+y^2+z^2)
  \]
  在点 $(1,1,1)$ 处沿该单位切向量 $\vec t$ 的方向导数。
  \item 求曲面 $\Sigma:z=\sqrt{x^2+y^2}$ 夹在两曲面 $x^2+y^2=y$，$x^2+y^2=2y$ 之间的那部分面积。
  \item 求微分方程 $y''-6y'+9y=0$ 的通解，并写出 $y''-6y'+9y=(x+1)e^{3x}$ 的特解形式。
  \item 判断方程 $(x^2-y)\,dx-x\,dy=0$ 是否是全微分方程，如果是，求其通解。
  \item 求幂级数
  \[
    \sum_{n=0}^{\infty}(-1)^n\frac{2n^2+1}{(2n+1)!}x^{2n+1}
  \]
  的和函数，并指出其收敛域。
\end{enumerate}

\paragraph*{三、证明下列各题（共 2 小题，每小题 6 分，共 12 分）}
\begin{enumerate}
  \item 设幂级数 $\sum_{n=0}^\infty a_nx^n$ 的收敛半径为 $R_1$，幂级数 $\sum_{n=0}^\infty b_nx^n$ 的收敛半径为 $R_2$，且 $R_1<R_2$。证明幂级数
  \[
    \sum_{n=0}^\infty (a_n+b_n)x^n
  \]
  的收敛半径为 $R_1$。
  \item 证明函数项级数
  \[
    \sum_{n=0}^\infty \frac{\cos(nx)}{2^n}
  \]
  在 $(-\infty,+\infty)$ 上一致收敛。
\end{enumerate}

\paragraph*{四、应用题（本题 8 分）}
在第一象限内作椭球面
\[
  \frac{x^2}{a^2}+\frac{y^2}{b^2}+\frac{z^2}{c^2}=1
\]
的切平面，使得切平面与三个坐标面所围成的四面体体积最小，求切点坐标。

\subsection{答案与解析}

\paragraph*{一、计算题}
\begin{enumerate}
  \item 令 $u=xy^2$，$v=x^2y$。记 $f_1,f_2$ 为 $f$ 对第一、第二个变量的偏导数，则
  \[
    z_y=f_1u_y+f_2v_y=2xyf_1+x^2f_2.
  \]
  继续对 $y$ 求导，得
  \[
    z_{yy}=4x^2y^2f_{11}+4x^3yf_{12}+x^4f_{22}+2xf_1,
  \]
  其中各偏导数均在 $(xy^2,x^2y)$ 处取值。
  \item 函数 $f$ 是奇函数，因此只需考虑间断点处 Fourier 级数取左右极限平均值。故
  \[
    S(x)=
    \begin{cases}
      -x, & |x|<\dfrac{\pi}{2},\\
      x, & \dfrac{\pi}{2}<|x|<\pi,\\
      0, & x=\pm\dfrac{\pi}{2}\text{ 或 }x=\pm\pi.
    \end{cases}
  \]
  \item 积分区域为 $0\le x\le 2$，$x/2\le y\le 1$。交换积分次序得
  \[
    0\le y\le 1,\qquad 0\le x\le 2y.
  \]
  因此
  \[
    \int_0^2 dx\int_{x/2}^1 e^{y^2}\,dy
    =\int_0^1dy\int_0^{2y}e^{y^2}\,dx
    =\int_0^1 2y e^{y^2}\,dy=e-1.
  \]
  \item 被积表达式为极角微分 $d\theta$。椭圆 $L$ 逆时针绕原点一周，故极角总增量为 $2\pi$，所以
  \[
    I=2\pi.
  \]
  \item 令
  \[
    P=4xz+x,\qquad Q=-2yz,\qquad R=1-z^2.
  \]
  曲面 $\Sigma$ 与平面 $z=1$ 围成区域
  \[
    \Omega=\{(x,y,z):x^2+y^2\le z\le 1\}.
  \]
  $\Sigma$ 取下侧时正好是 $\Omega$ 的外侧。由高斯公式，
  \[
    \nabla\cdot(P,Q,R)=(4z+1)-2z-2z=1.
  \]
  上盖圆盘 $z=1$ 上 $R=1-z^2=0$，通量为 $0$。于是
  \[
    I=\iiint_\Omega 1\,dV
    =\iint_{x^2+y^2\le 1}(1-x^2-y^2)\,dx\,dy
    =\frac{\pi}{2}.
  \]
\end{enumerate}

\paragraph*{二、解答题}
\begin{enumerate}
  \item 设
  \[
    F=x^2+y^2+z^2-3x,\qquad G=x+y-z.
  \]
  在 $(1,1,1)$ 处，
  \[
    \nabla F=(-1,2,2),\qquad \nabla G=(1,1,-1).
  \]
  因此切向量可取
  \[
    \nabla F\times\nabla G=(-4,1,-3),
  \]
  单位切向量为
  \[
    \vec t=\pm\frac{(-4,1,-3)}{\sqrt{26}}.
  \]
  又
  \[
    \nabla u(1,1,1)=\left(\frac23,\frac23,\frac23\right).
  \]
  若取 $\vec t=(-4,1,-3)/\sqrt{26}$，方向导数为
  \[
    D_{\vec t}u=\nabla u\cdot \vec t=-\frac4{\sqrt{26}}.
  \]
  若取相反方向，则方向导数为 $4/\sqrt{26}$。
  \item 用柱坐标表示两圆柱为 $r=\sin\theta$ 与 $r=2\sin\theta$，其中 $0\le\theta\le\pi$。圆锥 $z=r$ 的面积元为
  \[
    dS=\sqrt{1+z_x^2+z_y^2}\,dx\,dy=\sqrt2\,r\,dr\,d\theta.
  \]
  故所求面积
  \[
    S=\sqrt2\int_0^\pi\int_{\sin\theta}^{2\sin\theta}r\,dr\,d\theta
    =\frac{3\sqrt2}{2}\int_0^\pi\sin^2\theta\,d\theta
    =\frac{3\pi\sqrt2}{4}.
  \]
  \item 齐次方程的特征方程为
  \[
    \lambda^2-6\lambda+9=(\lambda-3)^2=0.
  \]
  因此通解为
  \[
    y=(C_1+C_2x)e^{3x}.
  \]
  对方程 $y''-6y'+9y=(x+1)e^{3x}$，由于 $3$ 是二重特征根，特解可设为
  \[
    y_p=x^2(Ax+B)e^{3x}.
  \]
  \item 记 $P=x^2-y$，$Q=-x$。因为
  \[
    P_y=-1,\qquad Q_x=-1,
  \]
  所以这是全微分方程。设势函数为 $\Phi$，则
  \[
    \Phi_x=x^2-y,
  \]
  从而
  \[
    \Phi=\frac{x^3}{3}-xy+\varphi(y).
  \]
  由 $\Phi_y=-x+\varphi'(y)=Q=-x$，得 $\varphi'(y)=0$。通解为
  \[
    \frac{x^3}{3}-xy=C.
  \]
  \item 设
  \[
    S(x)=\sum_{n=0}^{\infty}(-1)^n\frac{2n^2+1}{(2n+1)!}x^{2n+1}.
  \]
  对 $\sin x=\sum_{n=0}^\infty(-1)^n\dfrac{x^{2n+1}}{(2n+1)!}$ 使用算子 $D=x\dfrac{d}{dx}$。因 $2n+1=m$ 时
  \[
    2n^2+1=\frac{m^2-2m+3}{2},
  \]
  所以
  \[
    S(x)=\frac12(D^2-2D+3)\sin x.
  \]
  又
  \[
    D\sin x=x\cos x,\qquad D^2\sin x=x\cos x-x^2\sin x.
  \]
  故
  \[
    S(x)=\frac{(3-x^2)\sin x-x\cos x}{2}.
  \]
  该幂级数由正弦级数逐项微分和线性组合得到，收敛域为 $(-\infty,+\infty)$。
\end{enumerate}

\paragraph*{三、证明题}
\begin{enumerate}
  \item 当 $|x|<R_1$ 时，两个幂级数都收敛，所以 $\sum(a_n+b_n)x^n$ 收敛，故其收敛半径至少为 $R_1$。若其收敛半径大于 $R_1$，取 $x_0$ 使
  \[
    R_1<|x_0|<R_2.
  \]
  则 $\sum b_nx_0^n$ 收敛，而 $\sum a_nx_0^n$ 发散。若 $\sum(a_n+b_n)x_0^n$ 收敛，则
  \[
    \sum a_nx_0^n=\sum(a_n+b_n)x_0^n-\sum b_nx_0^n
  \]
  也应收敛，矛盾。因此 $\sum(a_n+b_n)x^n$ 的收敛半径为 $R_1$。
  \item 对一切实数 $x$，
  \[
    \left|\frac{\cos(nx)}{2^n}\right|\le \frac1{2^n}.
  \]
  而
  \[
    \sum_{n=0}^{\infty}\frac1{2^n}
  \]
  收敛，所以由 Weierstrass 判别法，原函数项级数在 $(-\infty,+\infty)$ 上一致收敛。
\end{enumerate}

\paragraph*{四、应用题}
设切点为 $(x_0,y_0,z_0)$，其中 $x_0,y_0,z_0>0$。椭球面在该点的切平面为
\[
  \frac{x_0x}{a^2}+\frac{y_0y}{b^2}+\frac{z_0z}{c^2}=1.
\]
它在三个坐标轴上的截距分别为
\[
  \frac{a^2}{x_0},\qquad \frac{b^2}{y_0},\qquad \frac{c^2}{z_0}.
\]
故四面体体积为
\[
  V_T=\frac16\frac{a^2}{x_0}\frac{b^2}{y_0}\frac{c^2}{z_0}
  =\frac{a^2b^2c^2}{6x_0y_0z_0}.
\]
因此要使 $V_T$ 最小，只需在约束
\[
  \frac{x_0^2}{a^2}+\frac{y_0^2}{b^2}+\frac{z_0^2}{c^2}=1
\]
下使 $x_0y_0z_0$ 最大。令
\[
  X=\frac{x_0}{a},\qquad Y=\frac{y_0}{b},\qquad Z=\frac{z_0}{c},
\]
则 $X^2+Y^2+Z^2=1$。由对称性或拉格朗日乘数法，最大值在
\[
  X=Y=Z=\frac1{\sqrt3}
\]
处取得。所以所求切点为
\[
  \left(\frac{a}{\sqrt3},\frac{b}{\sqrt3},\frac{c}{\sqrt3}\right).
\]

\section{2020级工科数学分析（二）期末考试 B 卷}

\subsection{试题}

\paragraph*{一、计算题（共 5 小题，每小题 8 分，共 40 分）}
\begin{enumerate}
  \item 设 $z=f(x+y,xy)$，求 $\dfrac{\partial^2z}{\partial x\partial y}$。
  \item 设
  \[
    f(x)=
    \begin{cases}
      x+1, & -\pi\le x<0,\\
      x^2, & 0\le x<\pi,
    \end{cases}
  \]
  写出 $f(x)$ 的以 $2\pi$ 为周期的 Fourier 级数的和函数 $S(x)$ 在 $[-\pi,\pi]$ 上的表达式。
  \item 计算二次积分
  \[
    \int_0^1dy\int_{y^2}^{y}y\sin(x^2)\,dx.
  \]
  \item 计算曲线积分
  \[
    I=\oint_L\frac{-y\,dx+x\,dy}{4x^2+y^2},
  \]
  其中 $L:(x-1)^2+y^2=4$，取逆时针方向。
  \item 计算曲面积分
  \[
    I=\iint_\Sigma \frac{x^3\,dy\,dz+y^3\,dz\,dx+z^3\,dx\,dy}{\sqrt{x^2+y^2+z^2}},
  \]
  其中 $\Sigma:z=\sqrt{a^2-x^2-y^2}$，取下侧。
\end{enumerate}

\paragraph*{二、解答下列各题（共 5 小题，每小题 8 分，共 40 分）}
\begin{enumerate}
  \item 求曲线
  \[
    \Gamma:
    \begin{cases}
      x^2+y^2+z^2=6,\\
      x+y+z=0
    \end{cases}
  \]
  在点 $(1,-2,1)$ 处的单位切向量 $\vec t$，并求函数
  \[
    u(x,y,z)=x^2+y^2+z^2
  \]
  在点 $(1,-2,1)$ 处沿该单位切向量 $\vec t$ 的方向导数。
  \item 求密度 $\rho=1$ 的均匀上半球面 $\Sigma:z=\sqrt{a^2-x^2-y^2}$ 对 $z$ 轴的转动惯量。
  \item 求微分方程 $y'+y\tan x=\cos x$ 的通解。
  \item 求微分方程 $y''-y=0$ 的通解，并写出 $y''-y=4xe^{-x}$ 的特解形式。
  \item 求幂级数
  \[
    \sum_{n=1}^{\infty}\frac{n^2+1}{2^n n!}x^n
  \]
  的和函数，并指出其收敛域。
\end{enumerate}

\paragraph*{三、证明下列各题（共 2 小题，每小题 6 分，共 12 分）}
\begin{enumerate}
  \item 设正项级数 $\sum_{n=0}^{\infty}a_n$ 收敛，证明正项级数
  \[
    \sum_{n=1}^{\infty}\frac{\sqrt{a_n}}{n}
  \]
  也收敛。【原卷第二个求和下限印作 $n=0$，但含分母 $n$，本稿按 $n=1$ 校正。】
  \item 设正项级数 $\sum_{n=0}^{\infty}a_n$ 收敛，证明函数
  \[
    f(x)=\sum_{n=0}^{\infty}a_nx^n
  \]
  在 $(-1,1)$ 上连续。
\end{enumerate}

\paragraph*{四、应用题（本题 8 分）}
要造一容积为 $V$ 的长方体无盖铝盒，请问怎样设计尺寸，使它的表面积最小？

\subsection{答案与解析}

\paragraph*{一、计算题}
\begin{enumerate}
  \item 令 $u=x+y$，$v=xy$。则
  \[
    z_x=f_1+ yf_2.
  \]
  对 $y$ 求导，注意 $u_y=1$，$v_y=x$，得
  \[
    z_{xy}=f_{11}+xf_{12}+yf_{21}+xyf_{22}+f_2.
  \]
  因 $f_{12}=f_{21}$，可写为
  \[
    z_{xy}=f_{11}+(x+y)f_{12}+xyf_{22}+f_2,
  \]
  其中各偏导数均在 $(x+y,xy)$ 处取值。
  \item Fourier 级数在连续点处收敛到函数值，在跳跃间断点处收敛到左右极限的平均值。故
  \[
    S(x)=
    \begin{cases}
      x+1, & -\pi<x<0,\\
      x^2, & 0<x<\pi,\\
      \dfrac12, & x=0,\\
      \dfrac{\pi^2-\pi+1}{2}, & x=\pm\pi.
    \end{cases}
  \]
  \item 积分区域为 $0\le y\le 1$，$y^2\le x\le y$。交换积分次序得
  \[
    0\le x\le 1,\qquad x\le y\le \sqrt{x}.
  \]
  因此
  \[
    \int_0^1dy\int_{y^2}^{y}y\sin(x^2)\,dx
    =\int_0^1\sin(x^2)\,dx\int_x^{\sqrt{x}}y\,dy
    =\frac12\int_0^1(x-x^2)\sin(x^2)\,dx.
  \]
  进一步化简，
  \[
    \frac12\int_0^1x\sin(x^2)\,dx=\frac{1-\cos1}{4},
  \]
  且
  \[
    \int_0^1x^2\sin(x^2)\,dx=-\frac{\cos1}{2}+\frac12\int_0^1\cos(x^2)\,dx.
  \]
  所以
  \[
    I=\frac14\left(1-\int_0^1\cos(x^2)\,dx\right).
  \]
  其中 $\int_0^1\cos(x^2)\,dx$ 可用 Fresnel 积分表示。
  \item 作变量变换 $u=2x$，$v=y$，则
  \[
    \frac{-y\,dx+x\,dy}{4x^2+y^2}
    =\frac12\frac{-v\,du+u\,dv}{u^2+v^2}.
  \]
  曲线 $L$ 在 $(u,v)$ 平面中的像仍逆时针绕原点一周，因此
  \[
    I=\frac12\cdot 2\pi=\pi.
  \]
  \item 在上半球面上 $\sqrt{x^2+y^2+z^2}=a$，向外单位法向量为 $(x,y,z)/a$。题设取下侧，即取相反方向。因此
  \[
    I=-\iint_\Sigma \frac{x^4+y^4+z^4}{a^2}\,dS.
  \]
  由对称性，上半球面上
  \[
    \iint_\Sigma x^4\,dS=\iint_\Sigma y^4\,dS=\iint_\Sigma z^4\,dS=\frac{2\pi a^6}{5}.
  \]
  故
  \[
    I=-\frac1{a^2}\cdot 3\cdot\frac{2\pi a^6}{5}
    =-\frac{6\pi a^4}{5}.
  \]
\end{enumerate}

\paragraph*{二、解答题}
\begin{enumerate}
  \item 令
  \[
    F=x^2+y^2+z^2,\qquad G=x+y+z.
  \]
  在 $(1,-2,1)$ 处，
  \[
    \nabla F=(2,-4,2),\qquad \nabla G=(1,1,1).
  \]
  切向量可取
  \[
    \nabla F\times\nabla G=(-6,0,6),
  \]
  故单位切向量为
  \[
    \vec t=\pm\frac{(-1,0,1)}{\sqrt2}.
  \]
  又 $\nabla u(1,-2,1)=(2,-4,2)$，所以
  \[
    D_{\vec t}u=\nabla u\cdot\vec t=0.
  \]
  \item 对 $z$ 轴的转动惯量为
  \[
    J_z=\iint_\Sigma (x^2+y^2)\rho\,dS.
  \]
  用球坐标参数化上半球面，$x^2+y^2=a^2\sin^2\varphi$，$dS=a^2\sin\varphi\,d\varphi\,d\theta$，其中 $0\le\varphi\le\pi/2$，$0\le\theta\le2\pi$。于是
  \[
    J_z=\int_0^{2\pi}\int_0^{\pi/2}a^2\sin^2\varphi\cdot a^2\sin\varphi\,d\varphi\,d\theta
    =2\pi a^4\cdot\frac23
    =\frac{4\pi a^4}{3}.
  \]
  \item 这是线性方程。积分因子为
  \[
    \mu(x)=e^{\int\tan x\,dx}=\sec x.
  \]
  因而
  \[
    (y\sec x)'=1.
  \]
  积分得
  \[
    y\sec x=x+C,
  \]
  即
  \[
    y=(x+C)\cos x.
  \]
  \item 齐次方程 $y''-y=0$ 的特征根为 $\lambda=\pm1$，故通解为
  \[
    y=C_1e^x+C_2e^{-x}.
  \]
  对 $y''-y=4xe^{-x}$，因 $e^{-x}$ 对应特征根 $-1$，且重数为 $1$，特解可设为
  \[
    y_p=x(Ax+B)e^{-x}.
  \]
  \item 令 $t=x/2$。则
  \[
    \sum_{n=1}^{\infty}\frac{n^2+1}{2^n n!}x^n
    =\sum_{n=1}^{\infty}\frac{(n^2+1)t^n}{n!}.
  \]
  利用
  \[
    \sum_{n=1}^{\infty}\frac{n^2t^n}{n!}=(t^2+t)e^t,\qquad
    \sum_{n=1}^{\infty}\frac{t^n}{n!}=e^t-1,
  \]
  得和函数
  \[
    S(x)=\left(\frac{x^2}{4}+\frac{x}{2}+1\right)e^{x/2}-1.
  \]
  该幂级数收敛域为 $(-\infty,+\infty)$。
\end{enumerate}

\paragraph*{三、证明题}
\begin{enumerate}
  \item 对任意 $N$，由 Cauchy--Schwarz 不等式，
  \[
    \sum_{n=1}^{N}\frac{\sqrt{a_n}}{n}
    \le
    \left(\sum_{n=1}^{N}a_n\right)^{1/2}
    \left(\sum_{n=1}^{N}\frac1{n^2}\right)^{1/2}.
  \]
  由于 $\sum a_n$ 与 $\sum 1/n^2$ 均收敛，右侧对 $N$ 有一致上界。又原级数为正项级数，其部分和单调递增且有上界，故
  \[
    \sum_{n=1}^{\infty}\frac{\sqrt{a_n}}{n}
  \]
  收敛。
  \item 对 $x\in(-1,1)$，每一项 $a_nx^n$ 都连续。并且在 $[-1,1]$ 上
  \[
    |a_nx^n|\le a_n.
  \]
  因为 $\sum a_n$ 收敛，由 Weierstrass 判别法，$\sum a_nx^n$ 在 $[-1,1]$ 上一致收敛，从而其和函数在 $[-1,1]$ 上连续，特别地在 $(-1,1)$ 上连续。
\end{enumerate}

\paragraph*{四、应用题}
设无盖铝盒底面边长分别为 $x,y$，高为 $h$。约束条件为
\[
  xyh=V.
\]
表面积为
\[
  S=xy+2xh+2yh=xy+\frac{2V}{y}+\frac{2V}{x}.
\]
求偏导并令其为零：
\[
  S_x=y-\frac{2V}{x^2}=0,\qquad
  S_y=x-\frac{2V}{y^2}=0.
\]
由此得 $x=y$。设 $x=y=t$，则
\[
  t^3=2V,\qquad h=\frac{V}{t^2}=\frac{t}{2}.
\]
所以当底面为正方形，边长
\[
  t=\sqrt[3]{2V},
\]
高为
\[
  h=\frac12\sqrt[3]{2V}
\]
时，表面积最小。

\section{2021级工科数学分析（二）期末考试 A 卷}

\subsection{试题}

\paragraph*{一、叙述定义（共 5 小题，每小题 2 分，共 10 分）}
\begin{enumerate}
  \item 二元函数的偏导数。
  \item 高斯公式。
  \item 条件收敛级数。
  \item 函数列的一致收敛。
  \item 二阶线性齐次常微分方程 $L[y]=0$ 的基础解系。
\end{enumerate}

\paragraph*{二、计算题（共 4 小题，每小题 10 分，共 40 分）}
\begin{enumerate}
  \item 设 $u=z^2+2z$，且 $z=z(x,y)$ 由方程 $xe^x-ye^y=ze^z$（$z\ne 1$）所确定，求 $du$ 和 $\dfrac{\partial^2u}{\partial x^2}$。
  \item 计算
  \[
    \int_1^2 dx\int_{\sqrt{x}}^x \sin\frac{\pi x}{2y}\,dy
    +\int_2^4 dx\int_{\sqrt{x}}^2 \sin\frac{\pi x}{2y}\,dy .
  \]
  \item 计算 $\iint_\Sigma x\,dy\,dz$，其中 $\Sigma$ 为 $x^2+y^2=1$，$z=x+2$ 和 $z=0$ 所围立体的表面，取外侧。
  \item 计算
  \[
    \int_L e^{\sqrt{x^2+y^2}}\,ds,
  \]
  其中 $L$ 为圆周 $x^2+y^2=a^2$（$a>0$）、直线 $y=x$ 与 $x$ 轴在第一象限内所围成的扇形的整个边界。
\end{enumerate}

\paragraph*{三、解答题（共 3 小题，每小题 10 分，共 30 分）}
\begin{enumerate}
  \item 设曲线积分
  \[
    \int_L xy^2\,dx+y\varphi(x)\,dy
  \]
  与路径无关，其中 $\varphi(x)$ 具有连续的导数，且 $\varphi(0)=0$，计算
  \[
    \int_{(0,0)}^{(1,1)} xy^2\,dx+y\varphi(x)\,dy .
  \]
  \item 求级数 $\displaystyle \sum_{n=1}^\infty \frac{n(n+1)}{2^n}$ 的和。
  \item 求方程 $y''+4y'+4y=\cos 2x$ 的通解。
\end{enumerate}

\paragraph*{四、应用题（10 分）}
造一容积为 $V$ 的长方形无盖铝盆，怎样设计尺寸，使它的表面积最小？

\paragraph*{五、证明题（10 分）}
证明级数 $\displaystyle \sum_{n=1}^\infty x^n(1-x)^2$ 在 $[0,1]$ 上一致收敛。

\subsection{答案与解析}

\paragraph*{一、叙述定义}
\begin{enumerate}
  \item 若极限
  \[
    \lim_{\Delta x\to 0}\frac{f(x_0+\Delta x,y_0)-f(x_0,y_0)}{\Delta x}
  \]
  存在，则称其为 $f$ 在 $(x_0,y_0)$ 处关于 $x$ 的偏导数，记为 $f_x(x_0,y_0)$；关于 $y$ 的偏导数类似。
  \item 若闭区域 $V$ 的边界为分片光滑闭曲面 $\Sigma$，取外侧，且 $P,Q,R$ 在 $V$ 上有连续一阶偏导数，则
  \[
    \iiint_V\left(\frac{\partial P}{\partial x}+\frac{\partial Q}{\partial y}+\frac{\partial R}{\partial z}\right)\,dV
    =
    \iint_\Sigma P\,dy\,dz+Q\,dz\,dx+R\,dx\,dy .
  \]
  \item 若 $\sum a_n$ 收敛而 $\sum |a_n|$ 发散，则称 $\sum a_n$ 条件收敛。
  \item 若函数列 $\{f_n(x)\}$ 在 $I$ 上逐点收敛于 $f(x)$，且对任意 $\varepsilon>0$，存在 $N$，使当 $n>N$ 时，对一切 $x\in I$ 都有
  \[
    |f_n(x)-f(x)|<\varepsilon,
  \]
  则称 $\{f_n\}$ 在 $I$ 上一致收敛于 $f$。
  \item 若二阶线性齐次方程 $L[y]=0$ 的两个解 $\varphi_1,\varphi_2$ 线性无关，则 $\{\varphi_1,\varphi_2\}$ 称为该方程的一个基础解系。
\end{enumerate}

\paragraph*{二、计算题}
\begin{enumerate}
  \item 令 $F(x,y,z)=xe^x-ye^y-ze^z$。由隐函数求导，
  \[
    z_x=\frac{x+1}{z+1}e^{x-z},\qquad
    z_y=-\frac{y+1}{z+1}e^{y-z}.
  \]
  题面括号中写作 $z\ne 1$，而隐函数求导实际要求 $F_z=-(z+1)e^z\ne0$，即 $z\ne-1$；以下按该条件计算。
  因为 $u=z^2+2z$，所以
  \[
    du=2(x+1)e^{x-z}\,dx-2(y+1)e^{y-z}\,dy,
  \]
  \[
    \frac{\partial^2u}{\partial x^2}
    =
    \frac{2e^{x-z}\bigl[(z+1)(x+2)-(x+1)^2e^{x-z}\bigr]}{z+1}.
  \]
  \item 交换积分次序，积分区域可写为 $1\le y\le 2,\ y\le x\le y^2$，故
  \[
    \int_1^2 dy\int_y^{y^2}\sin\frac{\pi x}{2y}\,dx
    =-\frac{2}{\pi}\int_1^2 y\cos\frac{\pi y}{2}\,dy
    =\frac{4(\pi+2)}{\pi^3}.
  \]
  \item 由高斯公式，取 $P=x,Q=R=0$，则
  \[
    \iint_\Sigma x\,dy\,dz=\iiint_\Omega 1\,dV
    =\iint_{x^2+y^2\le 1}(x+2)\,dx\,dy=2\pi.
  \]
  原答案抽取文本末行疑似将结果误写为 $4\pi$，本稿按高斯公式校为 $2\pi$。
  \item 边界由三段组成：$L_1:y=0,0\le x\le a$；$L_2:y=x,0\le x\le a/\sqrt2$；$L_3:x=a\cos\theta,y=a\sin\theta,0\le\theta\le\pi/4$。于是
  \[
    \int_{L_1}e^{\sqrt{x^2+y^2}}\,ds=e^a-1,\quad
    \int_{L_2}e^{\sqrt{x^2+y^2}}\,ds=e^a-1,
  \]
  \[
    \int_{L_3}e^{\sqrt{x^2+y^2}}\,ds=\frac{\pi a}{4}e^a.
  \]
  因而
  \[
    \int_L e^{\sqrt{x^2+y^2}}\,ds=2(e^a-1)+\frac{\pi a}{4}e^a.
  \]
\end{enumerate}

\paragraph*{三、解答题}
\begin{enumerate}
  \item 设 $P=xy^2,Q=y\varphi(x)$。路径无关给出
  \[
    P_y=Q_x,\qquad 2xy=y\varphi'(x),
  \]
  从而 $\varphi'(x)=2x$，又 $\varphi(0)=0$，故 $\varphi(x)=x^2$。沿折线 $(0,0)\to(1,0)\to(1,1)$ 计算，
  \[
    \int_{(0,0)}^{(1,1)} xy^2\,dx+y\varphi(x)\,dy
    =\int_0^1 y\,dy=\frac12.
  \]
  \item 对 $|x|<1$，设
  \[
    S(x)=\sum_{n=1}^\infty n(n+1)x^n.
  \]
  由 $\sum_{n=1}^\infty nx^{n-1}=1/(1-x)^2$ 可得
  \[
    S(x)=\frac{2x}{(1-x)^3}.
  \]
  取 $x=1/2$ 得
  \[
    \sum_{n=1}^\infty \frac{n(n+1)}{2^n}=S\left(\frac12\right)=8.
  \]
  \item 齐次方程特征根为 $r=-2$（二重根），故
  \[
    y_h=(C_1+C_2x)e^{-2x}.
  \]
  设特解 $y_p=A\sin 2x+B\cos 2x$，代入得 $A=1/8,B=0$，所以
  \[
    y=(C_1+C_2x)e^{-2x}+\frac18\sin 2x.
  \]
\end{enumerate}

\paragraph*{四、应用题}
设长方形无盖铝盆底边为 $x,y$，高为 $h$。约束为 $xyh=V$，表面积
\[
  S=xy+2xh+2yh=xy+\frac{2V}{y}+\frac{2V}{x}.
\]
由对称性和驻点方程可得 $x=y$，令 $x=y=t$，则
\[
  S(t)=t^2+\frac{4V}{t},\qquad S'(t)=2t-\frac{4V}{t^2}=0.
\]
故 $t^3=2V$，$h=V/t^2=t/2$。即底面为边长 $\sqrt[3]{2V}$ 的正方形，高为 $\sqrt[3]{2V}/2$ 时表面积最小。

\paragraph*{五、证明题}
记 $u_n(x)=x^n(1-x)^2$。在 $[0,1]$ 上，
\[
  u_n'(x)=nx^{n-1}(1-x)^2-2x^n(1-x),
\]
其最大值在 $x=n/(n+2)$ 处取得，且
\[
  0\le u_n(x)\le \left(\frac{n}{n+2}\right)^n\left(\frac{2}{n+2}\right)^2\le \frac{4}{n^2}.
\]
由于 $\sum 4/n^2$ 收敛，由 Weierstrass 判别法知
\[
  \sum_{n=1}^\infty x^n(1-x)^2
\]
在 $[0,1]$ 上一致收敛。

\section{2021级工科数学分析（二）期末考试 B 卷}

\subsection{试题}

\paragraph*{一、叙述定义（共 5 小题，每小题 2 分，共 10 分）}
\begin{enumerate}
  \item 二元函数的方向导数。
  \item 斯托克斯公式。
  \item 绝对收敛级数。
  \item 函数项级数的一致收敛。
  \item $n$ 阶线性齐次常微分方程 $L[y]=0$ 的基础解系。
\end{enumerate}

\paragraph*{二、计算题（共 4 小题，每小题 10 分，共 40 分）}
\begin{enumerate}
  \item 设 $z=\ln(u^2+v^2)$，而 $u=e^{x^2+y^2}$，$v=x^2+y$，求 $\dfrac{\partial z}{\partial x}$ 和 $\dfrac{\partial^2z}{\partial x\partial y}$。
  \item 计算
  \[
    \int_{-4}^0 dx\int_0^{2x+8}(8-y)^{1/3}\,dy
    +\int_0^4 dx\int_0^{-2x+8}(8-y)^{1/3}\,dy .
  \]
  \item 计算
  \[
    \iint_\Sigma (z^2+x)\,dy\,dz-z\,dx\,dy,
  \]
  其中 $\Sigma$ 为旋转抛物面 $z=\frac12(x^2+y^2)$ 介于平面 $z=0$ 与 $z=2$ 之间的部分的下侧。
  \item 计算
  \[
    I=\int_C (z-y)\,dx+(x-z)\,dy+(x-y)\,dz,
  \]
  其中 $C$ 是曲线
  \[
    \begin{cases}
      x^2+y^2=1,\\
      x-y+z=2,
    \end{cases}
  \]
  从 $z$ 轴正向往 $z$ 轴负向看 $C$ 的方向是顺时针的。
\end{enumerate}

\paragraph*{三、解答题（共 3 小题，每小题 10 分，共 30 分）}
\begin{enumerate}
  \item 设曲线积分
  \[
    \int_L x\varphi(y)\,dx+x^2y\,dy
  \]
  与路径无关，其中 $\varphi(y)$ 具有连续的导数，且 $\varphi(0)=1$，计算
  \[
    \int_{(0,0)}^{(1,1)}x\varphi(y)\,dx+x^2y\,dy.
  \]
  \item 求级数
  \[
    \sum_{n=2}^\infty (-1)^n\frac{x^n}{n(n-1)}
  \]
  的收敛域，并求和函数。
  \item 求方程 $y''+y'-2y=-2\sin x$ 的通解。
\end{enumerate}

\paragraph*{四、应用题（10 分）}
在椭圆 $x^2+4y^2=4$ 上求一点，使其到直线 $2x+3y-6=0$ 的距离最短。

\paragraph*{五、证明题（10 分）}
证明级数 $\displaystyle \sum_{n=1}^\infty x^2e^{-nx}$ 在 $(0,+\infty)$ 内一致收敛。

\subsection{答案与解析}

\paragraph*{一、叙述定义}
\begin{enumerate}
  \item 若极限
  \[
    \lim_{\rho\to0^+}\frac{f(x_0+\rho\cos\alpha,y_0+\rho\sin\alpha)-f(x_0,y_0)}{\rho}
  \]
  存在，则称其为 $f$ 在 $(x_0,y_0)$ 处沿方向 $\vec l=(\cos\alpha,\sin\alpha)$ 的方向导数。
  \item 若 $\Sigma$ 是分片光滑有向曲面，边界为正向闭曲线 $L$，且 $P,Q,R$ 有连续一阶偏导数，则
  \[
    \oint_L P\,dx+Q\,dy+R\,dz
    =
    \iint_\Sigma
    \begin{vmatrix}
      dy\,dz & dz\,dx & dx\,dy\\
      \partial_x & \partial_y & \partial_z\\
      P & Q & R
    \end{vmatrix}.
  \]
  \item 若级数 $\sum a_n$ 满足 $\sum |a_n|$ 收敛，则称 $\sum a_n$ 绝对收敛。
  \item 若函数项级数 $\sum u_n(x)$ 的部分和 $S_n(x)$ 在区间 $I$ 上一致收敛于 $S(x)$，即对任意 $\varepsilon>0$，存在 $N$，使 $n>N$ 时对一切 $x\in I$ 有
  \[
    |S_n(x)-S(x)|<\varepsilon,
  \]
  则称该级数在 $I$ 上一致收敛。
  \item $n$ 阶线性齐次常微分方程 $L[y]=0$ 的 $n$ 个线性无关解 $y_1,\ldots,y_n$ 构成的解组称为一个基础解系。
\end{enumerate}

\paragraph*{二、计算题}
\begin{enumerate}
  \item 令
  \[
    u=e^{x^2+y^2},\qquad v=x^2+y,\qquad A=u^2+v^2.
  \]
  则
  \[
    z_x=\frac{4x(u^2+v)}{A},
  \]
  且
  \[
    z_{xy}=
    \frac{4x\bigl[(4yu^2+1)A-(u^2+v)(4yu^2+2v)\bigr]}{A^2}.
  \]
  \item 积分区域可改写为
  \[
    0\le y\le8,\qquad \frac{y-8}{2}\le x\le\frac{8-y}{2}.
  \]
  因此原式为
  \[
    \int_0^8(8-y)^{1/3}(8-y)\,dy
    =\int_0^8(8-y)^{4/3}\,dy
    =\frac{384}{7}.
  \]
  \item 对图形 $z=\frac12(x^2+y^2)$ 取下侧，其向量面积元为 $(x,y,-1)\,dx\,dy$，投影区域为 $x^2+y^2\le4$。故
  \[
    \iint_\Sigma (z^2+x)\,dy\,dz-z\,dx\,dy
    =
    \iint_{x^2+y^2\le4}(xz^2+x^2+z)\,dx\,dy.
  \]
  奇函数项积分为 $0$，且 $z=(x^2+y^2)/2$，所以
  \[
    I=\iint_{x^2+y^2\le4}\left(x^2+\frac{x^2+y^2}{2}\right)\,dx\,dy=8\pi.
  \]
  \item 令 $\vec F=(z-y,x-z,x-y)$，则
  \[
    \nabla\times\vec F=(0,0,2).
  \]
  题设方向从 $z$ 轴正向看为顺时针，对应法向量的 $z$ 分量为负。由 Stokes 公式，
  \[
    I=\iint_D 2(-1)\,dx\,dy=-2\pi,
  \]
  其中 $D:x^2+y^2\le1$。
\end{enumerate}

\paragraph*{三、解答题}
\begin{enumerate}
  \item 记 $P=x\varphi(y)$，$Q=x^2y$。路径无关给出
  \[
    P_y=Q_x,\qquad x\varphi'(y)=2xy.
  \]
  故 $\varphi'(y)=2y$，由 $\varphi(0)=1$ 得 $\varphi(y)=y^2+1$。势函数可取
  \[
    \Phi(x,y)=\frac{x^2}{2}(y^2+1),
  \]
  所求积分为 $\Phi(1,1)-\Phi(0,0)=1$。
  \item 令 $t=-x$，则
  \[
    \sum_{n=2}^\infty(-1)^n\frac{x^n}{n(n-1)}
    =\sum_{n=2}^\infty\frac{t^n}{n(n-1)}.
  \]
  当 $|t|<1$ 时，
  \[
    \sum_{n=2}^\infty\frac{t^n}{n(n-1)}
    =(1-t)\ln(1-t)+t.
  \]
  因而
  \[
    S(x)=(1+x)\ln(1+x)-x.
  \]
  端点 $x=1$ 绝对收敛，$x=-1$ 也收敛，故收敛域为 $[-1,1]$。
  \item 齐次方程特征根为 $1,-2$。设特解为 $A\sin x+B\cos x$，代入得
  \[
    A=\frac35,\qquad B=\frac15.
  \]
  通解为
  \[
    y=C_1e^x+C_2e^{-2x}+\frac35\sin x+\frac15\cos x.
  \]
\end{enumerate}

\paragraph*{四、应用题}
椭圆可参数化为 $x=2\cos t,\ y=\sin t$。点到直线距离为
\[
  d=\frac{|2x+3y-6|}{\sqrt{13}}.
\]
在椭圆上有 $2x+3y=4\cos t+3\sin t\le5<6$，故距离最小时需使 $4\cos t+3\sin t$ 最大。取
\[
  \cos t=\frac45,\qquad \sin t=\frac35,
\]
得所求点为
\[
  \left(\frac85,\frac35\right).
\]

\paragraph*{五、证明题}
设余项
\[
  R_N(x)=\sum_{n=N}^\infty x^2e^{-nx}
  =\frac{x^2e^{-Nx}}{1-e^{-x}},\qquad x>0.
\]
函数 $h(x)=x^2/(e^x-1)$ 在 $(0,+\infty)$ 上有界，且当 $x\to0^+$ 与 $x\to+\infty$ 时均趋于 $0$。
给定 $\varepsilon>0$，先取 $\delta>0$ 使 $0<x<\delta$ 时 $h(x)<\varepsilon$。
再对 $x\ge\delta$ 用
\[
  R_N(x)=h(x)e^{-(N-1)x}\le M e^{-(N-1)\delta},
\]
取 $N$ 足够大即可使该项小于 $\varepsilon$。故 $R_N(x)\to0$ 关于 $x\in(0,+\infty)$ 一致成立，原级数一致收敛。

\section{2022级工科数学分析下 A卷}

\subsection{资料来源}
\begin{itemize}
  \item 2022级A卷无答案.pdf（试题，扫描型 PDF，已按页面人工转写）。
  \item 2022级A卷学生答案.pdf（学生答案，扫描型 PDF；能辨认处用于校对，缺口处按教材方法补写解析）。
\end{itemize}

\subsection{试题}

\paragraph*{一、填空题（共 5 题，每题 2 分，共 10 分）}
\begin{enumerate}
  \item 微分方程 $y''+2y'+6y=0$ 的通解为 \underline{\hspace{4cm}}。
  \item 设函数 $u=2x^2+2y^2+3z^2$，则 $\operatorname{div}(\operatorname{grad}u)=$ \underline{\hspace{3cm}}。
  \item 设 $\Gamma$ 为 $y^2=2x$ 上从原点到 $(2,2)$ 的一段，则第一类曲线积分
  \[
    \int_\Gamma y\,ds=
  \]
  \underline{\hspace{3cm}}。
  \item 参数曲线
  \[
    \begin{cases}
      x=t-\cos t,\\
      y=\sin t,\\
      z=2t
    \end{cases}
  \]
  在 $t=0$ 对应的点处的切线方程为 \underline{\hspace{5cm}}。
  \item 设周期为 $2\pi$ 的函数
  \[
    f(x)=
    \begin{cases}
      -1, & -\pi<x\le 0,\\
      1+x, & 0<x\le \pi,
    \end{cases}
  \]
  则 $f(x)$ 的傅里叶级数在 $x=\pi$ 处收敛于 \underline{\hspace{3cm}}。
\end{enumerate}

\paragraph*{二、选择题（共 5 题，每题 2 分，共 10 分）}
\begin{enumerate}
  \item 关于未知函数 $y$ 的微分方程
  \[
    (y-\cos^2 x)\,dx+e^x\,dy=0
  \]
  是（\quad）。

  \begin{tabular}{llll}
    A. 可分离变量方程； & B. 一阶线性非齐次方程； &
    C. 一阶线性齐次方程； & D. 非线性方程。
  \end{tabular}

  \item 若二元函数 $f(x,y)$ 满足
  \[
    \lim_{(x,y)\to(0,0)}
    \frac{f(x,y)-f(0,0)+3x-5y}{\sqrt{x^2+y^2}}=0,
  \]
  则（\quad）。

  \begin{tabular}{ll}
    A. $df(0,0)=0$； & B. $df(0,0)=3dx-5dy$；\\
    C. $df(0,0)=-3dx+5dy$； & D. $df(0,0)$ 不存在。
  \end{tabular}

  \item 设函数 $f(x,y)$ 与 $\varphi(x,y)$ 均可微，且 $\varphi_y(x,y)\ne0$。若 $(x_0,y_0)$ 是 $f(x,y)$ 在约束条件
  $\varphi(x,y)=0$ 下的一个极值点，则下列选项正确的是（\quad）。

  \begin{tabular}{ll}
    A. 若 $f_x(x_0,y_0)=0$，则 $f_y(x_0,y_0)=0$； &
    B. 若 $f_x(x_0,y_0)=0$，则 $f_y(x_0,y_0)\ne0$；\\
    C. 若 $f_x(x_0,y_0)\ne0$，则 $f_y(x_0,y_0)=0$； &
    D. 若 $f_x(x_0,y_0)\ne0$，则 $f_y(x_0,y_0)\ne0$。
  \end{tabular}

  \item 函数 $\ln(1+x)$ 在 $x=0$ 处的泰勒展开式正确的是（\quad）。

  \begin{tabular}{ll}
    A. $\displaystyle \ln(1+x)=\sum_{n=1}^{\infty}\frac{(-1)^n}{n}x^n,\ x\in(-1,1]$； &
    B. $\displaystyle \ln(1+x)=\sum_{n=1}^{\infty}\frac{(-1)^{n-1}}{n}x^n,\ x\in(-1,1]$；\\
    C. $\displaystyle \ln(1+x)=-\sum_{n=1}^{\infty}\frac{1}{n}x^n,\ x\in[-1,1)$； &
    D. $\displaystyle \ln(1+x)=\sum_{n=1}^{\infty}\frac{1}{n}x^n,\ x\in[-1,1)$。
  \end{tabular}

  \item 使得级数
  \[
    \sum_{n=1}^{\infty}\frac{(-1)^n}{n^p}
  \]
  条件收敛的常数 $p$ 的取值范围是（\quad）。

  \begin{tabular}{llll}
    A. $p\le0$； & B. $0<p\le1$； & C. $0<p<1$； & D. $p>1$。
  \end{tabular}
\end{enumerate}

\paragraph*{三、计算题（共 3 题，每题 10 分，共 30 分）}
\begin{enumerate}
  \item 设
  \[
    z=\frac{1}{x}f(x+y,x-y)+g(xy),
  \]
  其中函数 $f$ 有二阶连续的偏导数，且 $g$ 二阶可导，计算 $\dfrac{\partial z}{\partial x}$ 和
  $\dfrac{\partial^2z}{\partial x\partial y}$。

  \item 计算累次积分
  \[
    \int_{-\sqrt2}^{0}dx\int_{-x}^{\sqrt{4-x^2}}\sqrt{x^2+y^2}\,dy
    +\int_{0}^{2}dx\int_{\sqrt{2x-x^2}}^{\sqrt{4-x^2}}\sqrt{x^2+y^2}\,dy.
  \]

  \item 计算锥面 $z=\sqrt{x^2+y^2}$ 被柱面 $x^2+y^2=4x$ 截下的部分锥面面积。
\end{enumerate}

\paragraph*{四、综合题（共 3 题，每题 10 分，共 30 分）}
\begin{enumerate}
  \item 设 $\Sigma$ 是
  \[
    \frac{x^2}{a^2}+\frac{y^2}{b^2}+\frac{z^2}{c^2}=1\quad (z\ge0)
  \]
  的上侧，且 $a,b,c$ 都是正实数。计算第二类曲面积分
  \[
    \iint_\Sigma (x^2-y\sin x)\,dy\,dz+(y^2-z^2)\,dz\,dx+(z^2-x^2)\,dx\,dy.
  \]

  \item 计算曲线积分
  \[
    \int_\Gamma
    \frac{(x+y)\,dx-(x-y)\,dy}{x^2+y^2},
  \]
  其中曲线 $\Gamma$ 是从点 $A(-1,0)$ 到点 $B(1,0)$ 的一条不经过原点的光滑曲线 $y=f(x)$，
  $-1\le x\le1$。

  \item 求幂级数
  \[
    \sum_{n=0}^{\infty}(n+1)x^{2n+1}
  \]
  的收敛域及和函数。
\end{enumerate}

\paragraph*{五、证明题（共 1 题，每题 10 分，共 10 分）}
证明函数项级数
\[
  \sum_{n=1}^{\infty}\sin\frac{n!x^n}{x^2+n^n}
\]
在区间 $[-2,2]$ 上一致收敛。

\paragraph*{六、应用题（共 1 题，每题 10 分，共 10 分）}
制作一个体积为 $2$ 的长方体盒子，利用拉格朗日乘数法求出长、宽和高分别为多少时才能使其表面积最小。

\subsection{答案与解析}

\paragraph*{一、填空题}
\begin{enumerate}
  \item 特征方程为 $r^2+2r+6=0$，根为 $r=-1\pm\sqrt5\,i$，故
  \[
    y=e^{-x}\left(C_1\cos\sqrt5\,x+C_2\sin\sqrt5\,x\right).
  \]
  \item $\operatorname{div}(\operatorname{grad}u)=\Delta u=u_{xx}+u_{yy}+u_{zz}=4+4+6=14$。
  \item 令 $y=t,\ x=\dfrac{t^2}{2}$，$0\le t\le2$，则 $ds=\sqrt{1+t^2}\,dt$，所以
  \[
    \int_\Gamma y\,ds=\int_0^2 t\sqrt{1+t^2}\,dt
    =\frac{(1+t^2)^{3/2}}{3}\bigg|_0^2
    =\frac{5\sqrt5-1}{3}.
  \]
  \item 当 $t=0$ 时点为 $(-1,0,0)$，切向量为 $(1,1,2)$，故切线方程为
  \[
    \frac{x+1}{1}=\frac{y}{1}=\frac{z}{2}.
  \]
  \item 傅里叶级数在跳跃点收敛于左右极限的平均值。因
  \[
    f(\pi-0)=1+\pi,\qquad f(\pi+0)=f(-\pi+0)=-1,
  \]
  所以收敛值为 $\dfrac{\pi}{2}$。
\end{enumerate}

\paragraph*{二、选择题}
\begin{enumerate}
  \item 方程可化为 \(y'+e^{-x}y=e^{-x}\cos^2x\)，是一阶线性非齐次方程。选 B。
  \item 由可微定义，线性主部为 \(-3x+5y\)，故 \(df(0,0)=-3\,dx+5\,dy\)。选 C。
  \item 约束极值必要条件给出 \(f_x\varphi_y-f_y\varphi_x=0\)。因 \(\varphi_y\ne0\)，若 \(f_x\ne0\) 则必有 \(f_y\ne0\)。选 D。
  \item \(\displaystyle \ln(1+x)=\sum_{n=1}^{\infty}\frac{(-1)^{n-1}}{n}x^n\)，收敛区间为 $(-1,1]$。选 B。
  \item 交错 $p$ 级数在 $p>0$ 时收敛，在 $p>1$ 时绝对收敛，所以条件收敛范围为 $0<p\le1$。选 B。
\end{enumerate}

\paragraph*{三、计算题}
\begin{enumerate}
  \item 记 $u=x+y,\ v=x-y$，并把 $f_1,f_2,f_{11},f_{22}$ 均理解为在 $(u,v)$ 处的取值，把 $g',g''$
  理解为在 $xy$ 处的取值。先对 $x$ 求导：
  \[
    \frac{\partial z}{\partial x}
    =-\frac{f}{x^2}+\frac{f_1+f_2}{x}+y\,g'(xy).
  \]
  再对 $y$ 求导，利用 $u_y=1,\ v_y=-1$，得
  \[
    \frac{\partial^2z}{\partial x\partial y}
    =-\frac{f_1-f_2}{x^2}+\frac{f_{11}-f_{22}}{x}+g'(xy)+xy\,g''(xy).
  \]

  \item 第一部分区域为极坐标下
  \[
    \frac{\pi}{2}\le\theta\le\frac{3\pi}{4},\qquad 0\le r\le2,
  \]
  第二部分区域为
  \[
    0\le\theta\le\frac{\pi}{2},\qquad 2\cos\theta\le r\le2.
  \]
  因 $\sqrt{x^2+y^2}=r$，面积元为 $r\,dr\,d\theta$，所以原积分为
  \[
    \int_0^{\pi/2}\int_{2\cos\theta}^{2}r^2\,dr\,d\theta
    +\int_{\pi/2}^{3\pi/4}\int_0^2r^2\,dr\,d\theta
    =2\pi-\frac{16}{9}.
  \]

  \item 锥面写成 $z=\sqrt{x^2+y^2}$ 时，
  \[
    dS=\sqrt{1+z_x^2+z_y^2}\,dx\,dy=\sqrt2\,dx\,dy.
  \]
  投影区域为圆 $x^2+y^2\le4x$，即半径为 $2$ 的圆盘，面积为 $4\pi$。故锥面面积为
  \[
    S=\sqrt2\cdot4\pi=4\sqrt2\,\pi.
  \]
\end{enumerate}

\paragraph*{四、综合题}
\begin{enumerate}
  \item 记
  \[
    P=x^2-y\sin x,\qquad Q=y^2-z^2,\qquad R=z^2-x^2.
  \]
  用底面
  \[
    \Sigma_0:\ z=0,\quad \frac{x^2}{a^2}+\frac{y^2}{b^2}\le1
  \]
  补成上半椭球体 $\Omega$ 的外侧闭曲面，底面取向向下。由高斯公式，
  \[
    I_\Sigma+I_{\Sigma_0}
    =\iiint_\Omega (P_x+Q_y+R_z)\,dV
    =\iiint_\Omega (2x-y\cos x+2y+2z)\,dV.
  \]
  上半椭球关于 $x,y$ 对称，含 $x,y$ 的奇函数项积分为 $0$，故
  \[
    I_\Sigma+I_{\Sigma_0}=\iiint_\Omega 2z\,dV
    =\frac{\pi ab c^2}{2}.
  \]
  底面向下时
  \[
    I_{\Sigma_0}=\iint_{\frac{x^2}{a^2}+\frac{y^2}{b^2}\le1}x^2\,dx\,dy
    =\frac{\pi a^3b}{4}.
  \]
  因而
  \[
    I_\Sigma=\frac{\pi ab c^2}{2}-\frac{\pi a^3b}{4}.
  \]

  \item 积分可拆成
  \[
    \frac{(x+y)\,dx-(x-y)\,dy}{x^2+y^2}
    =\frac{x\,dx+y\,dy}{x^2+y^2}
    +\frac{y\,dx-x\,dy}{x^2+y^2}
    =d\ln r-d\theta.
  \]
  两端点到原点的距离都为 $1$，故 $d\ln r$ 的积分为 $0$。剩下的值由曲线绕过原点的方向决定。若
  $f(0)>0$，曲线从上半平面绕过原点，$\theta$ 可由 $\pi$ 连续变到 $0$，积分为 $\pi$；若
  $f(0)<0$，曲线从下半平面绕过原点，积分为 $-\pi$。因此
  \[
    \int_\Gamma
    \frac{(x+y)\,dx-(x-y)\,dy}{x^2+y^2}
    =
    \begin{cases}
      \pi, & f(0)>0,\\
      -\pi, & f(0)<0.
    \end{cases}
  \]

  \item 令 $t=x^2$，则
  \[
    \sum_{n=0}^{\infty}(n+1)x^{2n+1}
    =x\sum_{n=0}^{\infty}(n+1)t^n.
  \]
  当 $|t|<1$ 即 $|x|<1$ 时，
  \[
    \sum_{n=0}^{\infty}(n+1)t^n=\frac{1}{(1-t)^2},
  \]
  所以和函数为
  \[
    S(x)=\frac{x}{(1-x^2)^2},\qquad |x|<1.
  \]
  当 $x=\pm1$ 时通项不趋于 $0$，故发散。因此收敛域为 $(-1,1)$。
\end{enumerate}

\paragraph*{五、证明题}
对任意 $x\in[-2,2]$，有
\[
  \left|\sin\frac{n!x^n}{x^2+n^n}\right|
  \le
  \left|\frac{n!x^n}{x^2+n^n}\right|
  \le
  \frac{n!\,2^n}{n^n}.
\]
令 $M_n=\dfrac{n!\,2^n}{n^n}$，则
\[
  \frac{M_{n+1}}{M_n}
  =2\left(\frac{n}{n+1}\right)^n\longrightarrow \frac{2}{e}<1.
\]
故 $\sum M_n$ 收敛。由魏尔斯特拉斯判别法，原函数项级数在 $[-2,2]$ 上一致收敛。

\paragraph*{六、应用题}
设长方体长、宽、高分别为 $x,y,z>0$，体积约束为 $xyz=2$，闭盒子的表面积为
\[
  S=2(xy+xz+yz).
\]
设
\[
  L=2(xy+xz+yz)+\lambda(xyz-2).
\]
由拉格朗日乘数法，
\[
  2(y+z)+\lambda yz=0,\qquad
  2(x+z)+\lambda xz=0,\qquad
  2(x+y)+\lambda xy=0.
\]
两两相减得 $x=y=z$。再由 $xyz=2$，得
\[
  x=y=z=\sqrt[3]{2}.
\]
因此长、宽和高都等于 $\sqrt[3]{2}$ 时表面积最小。

\section{2022级工科数学分析下 B卷}

\subsection{试题}

\paragraph*{一、填空题：共 5 题，每题 2 分，共 10 分。}
\begin{enumerate}
  \item 微分方程 $y''+3y'+2y=0$ 的通解为\underline{\hspace{4cm}}。
  \item 函数 $u=2xy-z^2$ 在点 $(2,-1,-1)$ 处方向导数的最大值为\underline{\hspace{4cm}}。
  \item 设 $\Gamma$ 是圆周 $x^2+y^2=1$ 在第一象限的部分，则第一类曲线积分
  \[
    \int_\Gamma (x+y)^2\,ds=\underline{\hspace{4cm}}.
  \]
  \item 级数 $\displaystyle \sum_{n=0}^{\infty}(-1)^n\dfrac{2n+3}{(2n+1)!}$ 的和为\underline{\hspace{4cm}}。
  \item 设周期为 $2\pi$ 的函数
  \[
    f(x)=
    \begin{cases}
      x^2, & -\pi<x\le 0,\\
      0, & 0<x\le \pi,
    \end{cases}
  \]
  则 $f(x)$ 的傅里叶（Fourier）级数在 $x=-\pi$ 处收敛于\underline{\hspace{4cm}}。
\end{enumerate}

\paragraph*{二、选择题：共 5 题，每题 2 分，共 10 分。}
\begin{enumerate}
  \item 下列微分方程中，属于二阶线性常微分方程的是（\quad）
  \begin{enumerate}
    \item $y''+x^3y'+(\ln x)y=\tan x$；
    \item $(x^2+y^2)\,dy+y^2\,dx=0$；
    \item $(y'')^2+x^2y=1$；
    \item $y''+2\ln y=2x$。
  \end{enumerate}
  \item 已知函数 $f(x)$ 具有二阶连续导函数，且 $f(x)>0,\ f'(0)=0$（$f$ 在 $0$ 处的导数为 $0$）。则函数 $z=f(x)\ln f(y)$ 在点 $(0,0)$ 处取得极小值的一个充分条件是（\quad）
  \begin{enumerate}
    \item $f(0)>1,\ f''(0)>0$；
    \item $f(0)>1,\ f''(0)<0$；
    \item $f(0)<1,\ f''(0)>0$；
    \item $f(0)<1,\ f''(0)<0$。
  \end{enumerate}
  \item 函数 $f(x,y)=1+x+y$ 在区域 $\{(x,y)\mid x^2+y^2\le 1\}$ 上的最大值与最小值之积是（\quad）
  \begin{enumerate}
    \item $-1$；
    \item $1$；
    \item $3-\sqrt2$；
    \item $1+\sqrt2$。
  \end{enumerate}
  \item 设 $\Gamma$ 是以 $(1,1),(-1,1),(-1,-1),(1,-1)$ 为顶点的正方形边界，则第一类曲线积分
  \[
    \oint_\Gamma \frac{x+y+1}{|x|+|y|}\,ds
  \]
  等于（\quad）
  \begin{enumerate}
    \item $0$；
    \item $8$；
    \item $4\ln 2$；
    \item $8\ln 2$。
  \end{enumerate}
  \item 设 $a$ 为常数，则级数
  \[
    \sum_{n=1}^{\infty}\left(\frac{\sin(an)}{n^2}-\frac1{\sqrt n}\right)
  \]
  （\quad）
  \begin{enumerate}
    \item 发散；
    \item 绝对收敛；
    \item 条件收敛；
    \item 敛散性与 $a$ 的取值相关。
  \end{enumerate}
\end{enumerate}

\paragraph*{三、计算题：共 3 题，每题 10 分，共 30 分。}
\begin{enumerate}
  \item 设
  \[
    g(x,y)=f\left(xy,\frac12(x^2-y^2)\right),
  \]
  其中 $f(u,v)$ 具有连续的二阶偏导数，且满足
  \[
    \frac{\partial^2 f}{\partial u^2}+\frac{\partial^2 f}{\partial v^2}=1.
  \]
  计算二阶偏导数 $\displaystyle \frac{\partial^2 g}{\partial x^2}+\frac{\partial^2 g}{\partial y^2}$。
  \item 计算累次积分
  \[
    \int_0^1 dx\int_0^x dy\int_0^y \frac{\sin z}{(1-z)^2}\,dz.
  \]
  \item 计算二重积分
  \[
    \iint_D \sqrt{x^2+y^2}\,dx\,dy,
  \]
  其中 $D$ 是由 $x=\sqrt{2y-y^2}$ 与 $y=x$ 所围成的闭区域。
\end{enumerate}

\paragraph*{四、综合题：共 3 题，每题 10 分，共 30 分。}
\begin{enumerate}
  \item 设 $\Sigma$ 是曲面 $z=1-x^2-y^2\ (z\ge 0)$ 的上侧，计算第二型曲面积分
  \[
    \iint_\Sigma 2x^3\,dy\,dz+2y^3\,dz\,dx+3(z^2-1)\,dx\,dy.
  \]
  \item 计算曲线积分
  \[
    \oint_L yz\,dx+3zx\,dy-xy\,dz,
  \]
  其中 $L$ 是曲线
  \[
    \begin{cases}
      x^2+y^2=2y,\\
      y-z=1,
    \end{cases}
  \]
  其方向从 $z$ 轴正向往 $z$ 轴负向看去为逆时针方向。
  \item 将函数
  \[
    f(x)=\frac{x}{x^2-5x+6}
  \]
  展开成 $x-1$ 的幂级数，并指出其收敛域。
\end{enumerate}

\paragraph*{五、证明题：共 1 题，每题 10 分，共 10 分。}
证明函数项级数
\[
  \sum_{n=1}^{\infty}\arctan\frac{2x}{x^2+n^3}
\]
在区间 $(-\infty,+\infty)$ 上一致收敛。

\paragraph*{六、解答题：共 1 题，每题 10 分，共 10 分。}
限制点 $(x,y)$ 在圆周 $(x-1)^2+y^2=1$ 上变化时，求函数 $f(x,y)=xy$ 的最小值与最大值。

\subsection{答案与解析}

\paragraph*{一、填空题}
\begin{enumerate}
  \item 特征方程 $r^2+3r+2=0$，得 $r=-1,-2$，故通解为
  \[
    y=C_1e^{-x}+C_2e^{-2x}.
  \]
  \item $\nabla u=(2y,2x,-2z)$，在 $(2,-1,-1)$ 处为 $(-2,4,2)$，方向导数最大值为
  \[
    |\nabla u|=\sqrt{(-2)^2+4^2+2^2}=2\sqrt6.
  \]
  \item 令 $x=\cos t,\ y=\sin t,\ 0\le t\le \dfrac{\pi}{2}$，则
  \[
    \int_\Gamma (x+y)^2\,ds
    =\int_0^{\pi/2}(\cos t+\sin t)^2\,dt
    =\frac{\pi}{2}+1.
  \]
  \item
  \[
    \sum_{n=0}^{\infty}(-1)^n\frac{2n+3}{(2n+1)!}
    =
    \sum_{n=0}^{\infty}\frac{(-1)^n}{(2n)!}
    +2\sum_{n=0}^{\infty}\frac{(-1)^n}{(2n+1)!}
    =\cos 1+2\sin 1.
  \]
  \item 在 $x=-\pi$ 处，周期延拓函数的右极限为 $\pi^2$，左极限为 $0$，故 Fourier 级数收敛于
  \[
    \frac{\pi^2+0}{2}=\frac{\pi^2}{2}.
  \]
\end{enumerate}

\paragraph*{二、选择题}
\begin{enumerate}
  \item 选 A。A 为二阶线性非齐次常微分方程；B 为一阶方程；C 含 $(y'')^2$，D 含 $\ln y$，均非二阶线性方程。
  \item 选 A。由 $f'(0)=0$，点 $(0,0)$ 为驻点。Hessian 矩阵对角元为 $f''(0)\ln f(0)$ 与 $f''(0)$，交叉项为 $0$；当 $f(0)>1,\ f''(0)>0$ 时正定，故取得极小值。
  \item 选 A。在线性函数 $1+x+y$ 在单位圆盘上，最大值为 $1+\sqrt2$，最小值为 $1-\sqrt2$，二者乘积为 $-1$。
  \item 选 D。沿正方形四条边分段计算，位于上边与右边的积分各为 $4\ln2$，下边与左边的积分和为 $0$，故总积分为 $8\ln2$。
  \item 选 A。$\sum \sin(an)/n^2$ 绝对收敛，而 $\sum 1/\sqrt n$ 发散，故原级数发散。
\end{enumerate}

\paragraph*{三、计算题}
\begin{enumerate}
  \item 记 $u=xy,\ v=\dfrac12(x^2-y^2)$，并用 $f_1,f_2$ 表示 $f$ 对第一、第二个变量的偏导数。则
  \[
    g_x=yf_1+xf_2,\qquad g_y=xf_1-yf_2.
  \]
  继续求导得
  \[
    g_{xx}=y^2f_{11}+2xyf_{12}+x^2f_{22}+f_2,
  \]
  \[
    g_{yy}=x^2f_{11}-2xyf_{12}+y^2f_{22}-f_2.
  \]
  因此
  \[
    g_{xx}+g_{yy}=(x^2+y^2)(f_{11}+f_{22})=x^2+y^2.
  \]
  \item 积分区域为 $0\le z\le y\le x\le 1$。交换积分次序得
  \[
    \int_0^1\frac{\sin z}{(1-z)^2}
    \left(\int_z^1dy\int_y^1dx\right)dz
    =
    \int_0^1\frac{\sin z}{(1-z)^2}\cdot\frac{(1-z)^2}{2}\,dz.
  \]
  故原式为
  \[
    \frac12\int_0^1\sin z\,dz=\frac{1-\cos 1}{2}.
  \]
  \item 曲线 $x=\sqrt{2y-y^2}$ 等价于 $x^2+(y-1)^2=1$，在极坐标下边界为 $r=2\sin\theta$；直线 $y=x$ 为 $\theta=\dfrac{\pi}{4}$。故
  \[
    D:\quad 0\le \theta\le \frac{\pi}{4},\quad 0\le r\le 2\sin\theta.
  \]
  于是
  \[
    \iint_D\sqrt{x^2+y^2}\,dx\,dy
    =
    \int_0^{\pi/4}\int_0^{2\sin\theta}r^2\,dr\,d\theta
    =
    \frac{16-10\sqrt2}{9}.
  \]
\end{enumerate}

\paragraph*{四、综合题}
\begin{enumerate}
  \item 用平面 $\Sigma_1:z=0,\ x^2+y^2\le 1$ 补成闭曲面，取外侧方向。由高斯公式，
  \[
    \iint_{\Sigma+\Sigma_1}2x^3\,dy\,dz+2y^3\,dz\,dx+3(z^2-1)\,dx\,dy
    =
    \iiint_\Omega (6x^2+6y^2+6z)\,dV=2\pi.
  \]
  在补面 $\Sigma_1$ 上取下侧，积分为
  \[
    I_{\Sigma_1}=-\iint_{x^2+y^2\le 1}3(z^2-1)\,dx\,dy=3\pi.
  \]
  因此所求曲面积分为
  \[
    2\pi-3\pi=-\pi.
  \]
  \item 按题设方向可取参数
  \[
    x=\cos t,\qquad y=1+\sin t,\qquad z=\sin t,\qquad 0\le t\le 2\pi.
  \]
  代入得
  \[
    \oint_L yz\,dx+3zx\,dy-xy\,dz
    =
    \int_0^{2\pi}\left[-(1+\sin t)\sin^2t+3\sin t\cos^2t-(1+\sin t)\cos^2t\right]dt.
  \]
  利用一个周期内奇函数项积分为 $0$，得
  \[
    -\int_0^{2\pi}\sin^2t\,dt-\int_0^{2\pi}\cos^2t\,dt=-2\pi.
  \]
  \item 分解为
  \[
    \frac{x}{x^2-5x+6}=\frac{3}{x-3}-\frac{2}{x-2}.
  \]
  令 $t=x-1$，则
  \[
    \frac{3}{x-3}=-\frac32\cdot\frac1{1-t/2}
    =-\frac32\sum_{n=0}^{\infty}\left(\frac{t}{2}\right)^n,\quad |t|<2,
  \]
  \[
    -\frac{2}{x-2}=\frac2{1-t}
    =2\sum_{n=0}^{\infty}t^n,\quad |t|<1.
  \]
  因此
  \[
    f(x)=\sum_{n=0}^{\infty}\left(2-\frac{3}{2^{n+1}}\right)(x-1)^n,
    \qquad |x-1|<1.
  \]
  收敛域为 $(0,2)$。
\end{enumerate}

\paragraph*{五、证明题}
对任意 $x\in\mathbb R$，由 $|\arctan u|\le |u|$ 以及
\[
  x^2+n^3\ge 2|x|n^{3/2}
\]
可得
\[
  \left|\arctan\frac{2x}{x^2+n^3}\right|
  \le
  \frac{2|x|}{x^2+n^3}
  \le
  \frac1{n^{3/2}}.
\]
而级数 $\sum_{n=1}^{\infty}n^{-3/2}$ 收敛，故由 Weierstrass 判别法知
\[
  \sum_{n=1}^{\infty}\arctan\frac{2x}{x^2+n^3}
\]
在 $(-\infty,+\infty)$ 上一致收敛。

\paragraph*{六、解答题}
令约束函数
\[
  \varphi(x,y)=(x-1)^2+y^2-1=0.
\]
对 $F(x,y,\lambda)=xy+\lambda\varphi(x,y)$，有
\[
  y+2\lambda(x-1)=0,\qquad x+2\lambda y=0.
\]
由约束得 $x^2+y^2=2x$。当 $\lambda=0$ 时得点 $(0,0)$，函数值为 $0$，不是最值点。设 $y\ne0$，由第二个方程得
\[
  \lambda=-\frac{x}{2y}.
\]
代入第一个方程，得
\[
  y^2=x(x-1).
\]
再与 $y^2=2x-x^2$ 联立，得
\[
  x(x-1)=2x-x^2,\qquad x=\frac32,\qquad y=\pm\frac{\sqrt3}{2}.
\]
因此
\[
  f\left(\frac32,\frac{\sqrt3}{2}\right)=\frac{3\sqrt3}{4}
  \]
为最大值，
\[
  f\left(\frac32,-\frac{\sqrt3}{2}\right)=-\frac{3\sqrt3}{4}
  \]
为最小值。

\section{2023级工科数学分析（二）期末考试 A 卷}

\subsection{试题}

\paragraph*{一、填空题（共 5 小题，每小题 2 分，共 10 分）}
\begin{enumerate}
  \item 方程 $y'-\dfrac{y}{x}=x^2$ 的通解为 \underline{\hspace{4cm}}。
  \item 设 $z=\ln(x^2+y^2)$，$\vec n=\left(\cos\dfrac{\pi}{3},\sin\dfrac{\pi}{3}\right)$。则
  \[
    \left.\frac{\partial z}{\partial \vec n}\right|_{(1,1)}=\underline{\hspace{4cm}}.
  \]
  \item 设 $\Omega=\{(x,y,z):x^2+y^2+z^2\le 1\}$，则
  \[
    \iiint_\Omega (x^2+y^2+z^2)\,dV=\underline{\hspace{4cm}}.
  \]
  \item 设 $C$ 是从点 $A=(3,2,1)$ 到点 $B=(0,0,0)$ 的直线段，则
  \[
    \int_C x^3\,dx+y^3\,dy+z^3\,dz=\underline{\hspace{4cm}}.
  \]
  \item 设 $f(x)$ 是以 $2\pi$ 为周期的函数，它在 $[-\pi,\pi)$ 上的表达式是
  \[
    f(x)=
    \begin{cases}
      x+1, & -\pi\le x<0,\\
      x^2, & 0\le x<\pi.
    \end{cases}
  \]
  若它的 Fourier 级数的和函数是 $S(x)$，则 $S(-\pi)=\underline{\hspace{4cm}}$。
\end{enumerate}

\paragraph*{二、计算题（共 4 小题，每小题 9 分，共 36 分）}
\begin{enumerate}
  \item 求方程 $y''-y=\sin^2x$ 的通解。
  \item 设 $z=\arctan(x+y)$，$x=2u-v^2$，$y=u^2v$，求 $\dfrac{\partial z}{\partial u}$，$\dfrac{\partial z}{\partial v}$。
  \item 设 $D$ 是以 $(0,0)$，$(\pi,-\pi)$，$(\pi,\pi)$ 为顶点的三角形区域，求
  \[
    \iint_D \cos(x-y)\sin(x+y)\,dx\,dy.
  \]
  \item 设 $\Omega\subset\mathbb R^3$ 是由曲面 $z=5-x^2-y^2$ 和 $z+1=\sqrt{x^2+y^2}$ 所围成的空间立体，且 $\Omega$ 内布满某种流体，其流速为
  \[
    \vec F=\frac{\vec r}{\|\vec r\|^3},\qquad \vec r=(x,y,z).
  \]
  求单位时间内流体通过 $\Omega$ 的表面 $\partial\Omega$ 的流量。
\end{enumerate}

\paragraph*{三、解答题（共 3 小题，每小题 9 分，共 27 分）}
\begin{enumerate}
  \item 求级数 $\displaystyle \sum_{n=1}^\infty \frac{2n-1}{2^n}$ 的和。
  \item 求曲线积分
  \[
    \oint_L y\,dx+z\,dy+x\,dz,
  \]
  其中 $L$ 为圆周
  \[
    \begin{cases}
      x^2+y^2+z^2=a^2,\quad a>0,\\
      x+y+z=0,
    \end{cases}
  \]
  且从 $x$ 轴正向看去，这个圆周取逆时针方向。
  \item 给定
  \[
    du=\frac{(x+2y)\,dx+y\,dy}{(x+y)^2},
  \]
  试验证势函数 $u$ 的存在性，并求出 $u$。
\end{enumerate}

\paragraph*{四、应用题（9 分）}
求曲线
\[
  \begin{cases}
    z=x^2+y^2,\\
    y=\dfrac1x
  \end{cases}
\]
上到 $Oxy$ 平面距离最近的点。

\paragraph*{五、证明题（共 2 小题，每小题 9 分，共 18 分）}
\begin{enumerate}
  \item 证明曲线
  \[
    \begin{cases}
      x^2+y^2+z^2=6,\\
      x+y+z=0
    \end{cases}
  \]
  在点 $(1,-2,1)$ 处的法平面方程为 $x-z=0$。
  \item 证明级数
  \[
    \sum_{n=3}^\infty \ln\left(1+\frac{x}{n\ln^2 n}\right)
  \]
  在 $[-2,2]$ 内一致收敛。
\end{enumerate}

\subsection{答案与解析}

\paragraph*{一、填空题}
\begin{enumerate}
  \item $y=Cx+\dfrac{x^3}{2}$。
  \item $\dfrac{1+\sqrt3}{2}$。
  \item $\dfrac{4\pi}{5}$。
  \item $-\dfrac{49}{2}$。
  \item $\dfrac{\pi^2-\pi+1}{2}$。
\end{enumerate}

\paragraph*{二、计算题}
\begin{enumerate}
  \item 因 $\sin^2x=\dfrac{1-\cos2x}{2}$，通解为
  \[
    y=C_1e^x+C_2e^{-x}-\frac12+\frac1{10}\cos2x.
  \]
  \item 令 $s=x+y=2u-v^2+u^2v$，则
  \[
    \frac{\partial z}{\partial u}
    =\frac{2+2uv}{1+s^2},\qquad
    \frac{\partial z}{\partial v}
    =\frac{-2v+u^2}{1+s^2}.
  \]
  \item 作变换 $u=x-y,\ v=x+y$，则 $dx\,dy=\frac12\,du\,dv$，区域变为
  \[
    u\ge0,\quad v\ge0,\quad u+v\le2\pi.
  \]
  因此
  \[
    I=\frac12\int_0^{2\pi}\int_0^{2\pi-u}\cos u\sin v\,dv\,du
    =-\frac{\pi}{2}.
  \]
  \item 该区域包含原点，而 $\vec F=\vec r/\|\vec r\|^3$ 在原点有孤立奇点。去掉原点附近小球并用高斯公式，外边界通量等于小球面通量
  \[
    \iint_{\partial B_\varepsilon}\frac{\vec r}{\|\vec r\|^3}\cdot\vec n\,dS=4\pi.
  \]
  故通过 $\partial\Omega$ 的流量为 $4\pi$。
\end{enumerate}

\paragraph*{三、解答题}
\begin{enumerate}
  \item
  \[
    \sum_{n=1}^\infty\frac{2n-1}{2^n}
    =2\sum_{n=1}^\infty\frac{n}{2^n}-\sum_{n=1}^\infty\frac1{2^n}
    =4-1=3.
  \]
  \item 令 $\vec F=(y,z,x)$，则 $\nabla\times\vec F=(-1,-1,-1)$。按题设方向取单位法向
  \[
    \vec n=\frac{(1,1,1)}{\sqrt3}.
  \]
  圆盘半径为 $a$，由 Stokes 公式得
  \[
    \oint_L y\,dx+z\,dy+x\,dz
    =(-1,-1,-1)\cdot\frac{(1,1,1)}{\sqrt3}\,\pi a^2
    =-\sqrt3\,\pi a^2.
  \]
  \item 记
  \[
    P=\frac{x+2y}{(x+y)^2},\qquad Q=\frac{y}{(x+y)^2}.
  \]
  有 $P_y=Q_x=-2y/(x+y)^3$，故势函数存在。对 $P$ 关于 $x$ 积分，得
  \[
    u=\ln|x+y|-\frac{y}{x+y}+C.
  \]
\end{enumerate}

\paragraph*{四、应用题}
到 $Oxy$ 平面的距离为 $|z|$。在曲线上 $z=x^2+y^2=x^2+x^{-2}$，故
\[
  z\ge2,
\]
等号在 $x=\pm1$ 时成立。最近点为
\[
  (1,1,2),\qquad (-1,-1,2).
\]

\paragraph*{五、证明题}
\begin{enumerate}
  \item 设
  \[
    F=x^2+y^2+z^2-6,\qquad G=x+y+z.
  \]
  在 $(1,-2,1)$ 处，
  \[
    \nabla F=(2,-4,2),\qquad \nabla G=(1,1,1),
  \]
  切向量可取 $\nabla F\times\nabla G=(-6,0,6)$，即平行于 $(1,0,-1)$。法平面垂直于切向量，故
  \[
    (1,0,-1)\cdot(x-1,y+2,z-1)=0,
  \]
  即 $x-z=0$。
  \item 当 $|x|\le2$ 时，各项均有定义。对充分大的 $n$，
  \[
    \left|\ln\left(1+\frac{x}{n\ln^2n}\right)\right|
    \le C\frac{|x|}{n\ln^2n}
    \le \frac{2C}{n\ln^2n}.
  \]
  因 $\sum_{n=3}^\infty 1/(n\ln^2n)$ 收敛，有限个初始项不影响一致收敛，由 Weierstrass 判别法知原级数在 $[-2,2]$ 上一致收敛。
\end{enumerate}

\section{2023级工科数学分析（二）期末考试 B 卷}

\subsection{试题}

\paragraph*{一、填空题（共 5 小题，每小题 2 分，共 10 分）}
\begin{enumerate}
  \item 方程 $y'+y\tan x=\cos x$ 的通解为 \underline{\hspace{4cm}}。
  \item 设 $z=e^{xy}$，$\vec n=\left(\cos\dfrac{\pi}{4},\sin\dfrac{\pi}{4}\right)$。则
  \[
    \left.\frac{\partial z}{\partial \vec n}\right|_{(1,1)}=\underline{\hspace{4cm}}.
  \]
  \item
  \[
    \int_0^1 dx\int_0^{\sqrt{1-x^2}}x^3\,dy=\underline{\hspace{4cm}}.
  \]
  \item 设曲线的参数方程为
  \[
    x=at,\qquad y=\frac a2t^2,\qquad z=\frac a3t^3,\qquad 0\le t\le 1,
  \]
  且其密度函数为 $\sqrt{2y/a}$，则曲线的质量为 \underline{\hspace{4cm}}。
  \item 设 $f(x)$ 是以 $2\pi$ 为周期的函数，它在 $[-\pi,\pi)$ 上的表达式是 $f(x)=x$。则 $f(x)$ 的 Fourier 级数中 $\sin 3x$ 的系数是 \underline{\hspace{4cm}}。
\end{enumerate}

\paragraph*{二、计算题（共 4 小题，每小题 9 分，共 36 分）}
\begin{enumerate}
  \item 求方程 $y''+y=3xe^x$ 的通解。
  \item 设 $z=\sin\left(xy+\dfrac{x}{y}\right)$，求 $\dfrac{\partial^2z}{\partial x\partial y}$。
  \item 设 $D$ 是由 $x^2+y^2=9$，$x^2+y^2=16$ 以及 $y^2-x^2=1$，$y^2-x^2=9$ 围成的位于第一象限的区域，求
  \[
    \iint_D xy\,dx\,dy.
  \]
  \item 设 $\Omega\subset\mathbb R^3$ 以原点为内点，且边界曲面 $\partial\Omega$ 光滑，又设边界曲面单位外法向为 $\vec n$，向量场
  \[
    \vec F=\frac{\vec r}{\|\vec r\|^3},\qquad \vec r=(x,y,z).
  \]
  求
  \[
    \iint_{\partial\Omega}\vec F\cdot\vec n\,ds.
  \]
\end{enumerate}

\paragraph*{三、解答题（共 3 小题，每小题 9 分，共 27 分）}
\begin{enumerate}
  \item 求级数 $\displaystyle \sum_{n=1}^\infty \frac1{4n(4n+2)}$ 的和。
  \item 原卷配图表示第一卦限内的平面三角形及其法向；本题条件如下。设有向曲面
  \[
    \Sigma=\{(x,y,z):x+y+z=1,\ x\ge 0,\ y\ge 0,\ z\ge 0\},
  \]
  且其方向与 $(1,1,1)$ 同向。求力场 $\vec F=(y^2,z^2,x^2)$ 绕 $\Sigma$ 正向边界 $\partial\Sigma$ 一周所做的功。
  \item 已知 $\phi(\pi)=1$，试确定 $\phi(x)$，使得曲线积分
  \[
    I=\int_A^B \left[\sin x-\phi(x)\right]\frac{y}{x}\,dx+\phi(x)\,dy
  \]
  与路径无关，并求当 $A,B$ 两点分别为 $(1,0)$，$(\pi,\pi)$ 时积分 $I$ 的值。
\end{enumerate}

\paragraph*{四、应用题（9 分）}
在力场 $\vec F=(yz,zx,xy)$ 的作用下，质点由原点沿直线运动到椭球面
\[
  \frac{x^2}{a^2}+\frac{y^2}{b^2}+\frac{z^2}{c^2}=1
\]
上第一卦限的点 $M=(\xi,\eta,\zeta)$。问当 $\xi,\eta,\zeta$ 取何值时力场 $\vec F$ 做的功最大？最大值是多少？

\paragraph*{五、证明题（共 2 小题，每小题 9 分，共 18 分）}
\begin{enumerate}
  \item 证明曲线
  \[
    \begin{cases}
      x^2+y^2+z^2-3x=0,\\
      2x-3y+5z-4=0
    \end{cases}
  \]
  在点 $(1,1,1)$ 处的法平面方程为 $16x+9y-z-24=0$。
  \item 证明级数
  \[
    \sum_{n=1}^\infty \sin\frac{\ln n^x}{n^3+x^2}
  \]
  在 $(-\infty,+\infty)$ 内一致收敛。
\end{enumerate}

\subsection{答案与解析}

\paragraph*{一、填空题}
\begin{enumerate}
  \item $y=(x+C)\cos x$。
  \item $\sqrt2\,e$。
  \item $\dfrac{2}{15}$。
  \item
  \[
    \frac{a(3\sqrt3-1)}{8}
    +\frac{3a}{16}\ln\frac{3+2\sqrt3}{3}.
  \]
  \item $\dfrac23$。
\end{enumerate}

\paragraph*{二、计算题}
\begin{enumerate}
  \item 通解为
  \[
    y=C_1\cos x+C_2\sin x+\frac32(x-1)e^x.
  \]
  \item 令
  \[
    w=xy+\frac{x}{y}.
  \]
  因 $z_x=(y+y^{-1})\cos w$，故
  \[
    z_{xy}
    =\left(1-\frac1{y^2}\right)\cos w
    -\left(y+\frac1y\right)\left(x-\frac{x}{y^2}\right)\sin w.
  \]
  \item 作变换
  \[
    u=x^2+y^2,\qquad v=y^2-x^2.
  \]
  第一象限内 $du\,dv=8xy\,dx\,dy$，区域变为 $9\le u\le16,\ 1\le v\le9$，故
  \[
    \iint_D xy\,dx\,dy
    =\frac18(16-9)(9-1)=7.
  \]
  \item 因 $\Omega$ 以原点为内点，而 $\vec F=\vec r/\|\vec r\|^3$ 在原点有孤立奇点，去掉原点附近小球并用高斯公式可得
  \[
    \iint_{\partial\Omega}\vec F\cdot\vec n\,dS=4\pi.
  \]
\end{enumerate}

\paragraph*{三、解答题}
\begin{enumerate}
  \item
  \[
    \sum_{n=1}^\infty\frac1{4n(4n+2)}
    =\frac18\sum_{n=1}^\infty\frac1{n(2n+1)}
    =\frac{1-\ln2}{4}.
  \]
  \item 令 $\vec F=(y^2,z^2,x^2)$，则
  \[
    \nabla\times\vec F=(-2z,-2x,-2y).
  \]
  在平面 $x+y+z=1$ 上取与 $(1,1,1)$ 同向的单位法向量，则
  \[
    (\nabla\times\vec F)\cdot\vec n=-\frac2{\sqrt3}.
  \]
  三角形面积为 $\sqrt3/2$，由 Stokes 公式得所做的功为 $-1$。本题原卷有图，当前转写按题干给出的正向约定处理。
  \item 记
  \[
    P=\left[\sin x-\phi(x)\right]\frac{y}{x},\qquad Q=\phi(x).
  \]
  路径无关要求
  \[
    P_y=Q_x,\qquad \phi'(x)+\frac{\phi(x)}{x}=\frac{\sin x}{x}.
  \]
  因而 $(x\phi(x))'=\sin x$。由 $\phi(\pi)=1$ 得
  \[
    \phi(x)=\frac{\pi-1-\cos x}{x}.
  \]
  势函数可取 $U(x,y)=\phi(x)y$，故
  \[
    I=U(\pi,\pi)-U(1,0)=\pi.
  \]
\end{enumerate}

\paragraph*{四、应用题}
沿直线从原点到 $M=(\xi,\eta,\zeta)$，取参数 $\vec r(t)=tM$，$0\le t\le1$。则
\[
  W=\int_0^1 3t^2\xi\eta\zeta\,dt=\xi\eta\zeta.
\]
在约束
\[
  \frac{\xi^2}{a^2}+\frac{\eta^2}{b^2}+\frac{\zeta^2}{c^2}=1
\]
下，令 $X=\xi/a,\ Y=\eta/b,\ Z=\zeta/c$，由对称性或拉格朗日乘数法知 $XYZ$ 在
\[
  X=Y=Z=\frac1{\sqrt3}
\]
处最大。故
\[
  M=\left(\frac a{\sqrt3},\frac b{\sqrt3},\frac c{\sqrt3}\right),\qquad
  W_{\max}=\frac{abc}{3\sqrt3}.
\]

\paragraph*{五、证明题}
\begin{enumerate}
  \item 设
  \[
    F=x^2+y^2+z^2-3x,\qquad G=2x-3y+5z-4.
  \]
  在 $(1,1,1)$ 处，
  \[
    \nabla F=(-1,2,2),\qquad \nabla G=(2,-3,5).
  \]
  切向量可取
  \[
    \nabla F\times\nabla G=(16,9,-1).
  \]
  法平面垂直于切向量，故
  \[
    16(x-1)+9(y-1)-(z-1)=0,
  \]
  即 $16x+9y-z-24=0$。
  \item 因 $\ln n^x=x\ln n$，对任意 $x\in\mathbb R$ 有
  \[
    \left|\sin\frac{\ln n^x}{n^3+x^2}\right|
    \le \frac{|x|\ln n}{n^3+x^2}
    \le \frac{\ln n}{2n^{3/2}}.
  \]
  而 $\sum_{n=1}^\infty \ln n/n^{3/2}$ 收敛，故由 Weierstrass 判别法，原级数在 $(-\infty,+\infty)$ 上一致收敛。
\end{enumerate}


